\documentclass[11pt]{article}

\usepackage[T1]{fontenc}
\usepackage[utf8]{inputenc}
\usepackage{lmodern}
\usepackage{microtype}
\usepackage{amsmath,amssymb}
\usepackage[margin=1.1in]{geometry}
\usepackage{booktabs}
\usepackage{tabularx}
\usepackage{array}
\usepackage[colorlinks=true,linkcolor=blue,urlcolor=blue,citecolor=blue]{hyperref}

\title{Bayesian Inference for Nonlinear Models with a Genuine Zero Lower
Bound\\[0.6ex]\large A survey of exact and approximate sampling methods,
with two case studies}
\author{}
\date{August 2026}

\setcounter{tocdepth}{2}

\begin{document}
\maketitle

\begin{abstract}
\noindent
A zero (or effective) lower bound on nominal interest rates changes the geometry of Bayesian inference. Calling every such change ``discontinuous'' obscures the distinctions that determine sampler choice. The censoring rule $i_t=\max\{\ell,i_t^\star\}$ is continuous with a kink: its slope changes at the bound. A deterministic solution map that selects between equilibrium branches can create a genuine jump in the density. A stochastic regime yields a mixed posterior containing both discrete and continuous variables. When equilibrium is not unique, no posterior is defined until a selection law specifies how an equilibrium is chosen. These four geometries call for different samplers; in the last case, the target must first be defined. Misclassifying one geometry as another is the most common design error in this literature. The survey covers the linear benchmark (OccBin and piecewise Kalman filtering), particle-filter likelihoods and particle MCMC, and Hamiltonian Monte Carlo (HMC) on kinked and piecewise-smooth targets. It also covers event-driven HMC (reflection/refraction), discontinuous and mixed variants of HMC, and pseudo-marginal and differentiable-particle-filter methods. For these methods, the treatment derives the key constructions rather than merely citing them. Two case studies apply the classification. In a two-country shadow-rate term-structure model, the hard-bound variant produces a kink target with a tractable particle likelihood. In a small New Keynesian model with an occasionally binding policy rate, the genuine discontinuity lies in the solution correspondence: verified paths can fail to exist or can be multiple. The max operator itself remains continuous.
\end{abstract}

\tableofcontents
\newpage

\section{Introduction}\label{sec:intro}

In a linearized New Keynesian model, a zero lower bound (ZLB), or effective lower bound (ELB), can be handled by solving two linear systems: one with the policy constraint slack and one with it binding. OccBin and its successors iterate over a candidate sequence of these regimes. They apply a Kalman filter conditional on that sequence and then verify that the resulting path satisfies the inequalities defining the proposed regimes (Guerrieri and Iacoviello 2015; Giovannini, Pfeiffer, and Ratto 2021). The approach is valuable and fast, but its scope is limited by conditional Gaussian closure. This closure enters twice: conditional on the regime sequence, the approximation uses Gaussian state transitions and maintains Gaussian filtering distributions in the likelihood recursion. Integrating over states under these assumptions therefore gives the likelihood of a piecewise-linear Gaussian approximation, not the likelihood of the nonlinear model.

The survey therefore addresses a narrower and harder problem: sampling a posterior that retains an actual ZLB regime change in a nonlinear state-space model. The question is how to do so while keeping the sampling transition kernel exact for the model being estimated and making every approximation explicit.

Sampler choice first requires care with the word ``discontinuous'', which is routinely used for at least three different mathematical situations. The map
\[
 i_t=\max\{\ell,i_t^\star\}
\]
is continuous and has a kink at \(i_t^\star=\ell\). Below the bound, changes in the shadow rate leave the policy rate unchanged; above the bound, the policy rate follows the shadow rate. The change in slope creates a kink without a jump in the rate. A model with an explicit regime variable \(r_t\in\{n,b\}\) can instead produce a target with discrete regime labels alongside continuous variables. A solver that selects between distinct equilibrium branches can produce a genuine jump in the density. These are three different posterior geometries. The modern Hamiltonian Monte Carlo (HMC) literature addresses them with three different tools: reflection/refraction HMC, discontinuous HMC, and mixed HMC. Treating all three geometries as one kind of nonsmoothness is the first design error a practitioner can make, and much of the survey is devoted to keeping them distinct.

The survey has four goals. First, it derives the linear benchmark, including the piecewise-linear solution and its conditional Kalman likelihood, so that readers can see precisely what the fast standard method computes and what it leaves out. Second, it derives the HMC kernels appropriate for piecewise-smooth and mixed-support targets. Third, it derives the nonlinear state-space likelihood and its particle approximation, making explicit what an exact method must ultimately integrate. Fourth, it puts the comparison into practice through two case studies: a shadow-rate term-structure model (Section~\ref{sec:shadow}) and a small New Keynesian model with an occasionally binding policy rate (Section~\ref{sec:nk}). Throughout the survey, results stated by a paper are cited to that paper. The authors' own deductions use the notation introduced below and are labelled as such.

\section{A common model and notation}\label{sec:model}

To compare methods developed in different literatures, the survey first fixes a single notation broad enough to contain them all. In this notation, \(\vartheta\) denotes structural parameters, \(x_t\in\mathbb R^{d_x}\) the latent state, \(r_t\) the ZLB regime, \(\epsilon_t\) a structural shock, and \(y_t\) the observation. A general model links the regime rule to the state transition and observation law:
\begin{equation}\label{eq:1}
\begin{aligned}
 r_t &= \mathcal R(x_{t-1},\epsilon_t,\vartheta),\\
 x_t &= F_{r_t}(x_{t-1},\epsilon_t,\vartheta),\\
 y_t &\sim g_{r_t}(y_t\mid x_t,\vartheta).
\end{aligned}
\end{equation}
The first line determines the regime from the previous state, the current shock, and the structural parameters. The second line uses those same inputs in the transition function for the selected regime to obtain the current state. The third line gives the observation distribution conditional on the current state and the parameters, again indexed by the same regime. The regime can therefore affect both state dynamics and the observation law.

For a monetary lower bound, the shadow or notional nominal rate is defined by \(i_t^\star=\iota(x_{t-1},\epsilon_t,\vartheta)\), and \(\ell\) denotes the bound. The shadow rate thus depends on the previous state, the current shock, and the structural parameters. The regime rule is
\begin{equation}\label{eq:2}
 r_t=n \iff i_t^\star>\ell,\qquad
 r_t=b \iff i_t^\star\leq\ell,
\end{equation}
and the observed policy rate may be \(i_t=i_t^\star\) in the nonbinding regime \(n\) and \(i_t=\ell\) in the binding regime \(b\). The nonbinding regime applies exactly when the shadow rate is strictly above the bound; the binding regime applies when it is at or below the bound, including equality.

An equivalent complementarity form introduces a nonnegative multiplier \(\lambda_t\geq0\):
\begin{equation}\label{eq:3}
 i_t-\ell\geq0,\qquad \lambda_t\geq0,\qquad
 (i_t-\ell)\lambda_t=0.
\end{equation}
Equation \eqref{eq:3} requires the gap between the policy rate and the bound, and the multiplier, to be nonnegative, with their product equal to zero. A rate strictly above the bound therefore requires a zero multiplier, while a positive multiplier requires the rate to equal the bound. The gap and multiplier may both be zero at the boundary. These conditions describe which inequalities hold in the slack and binding cases, but they do not by themselves specify how the rest of the equilibrium is solved. That distinction matters for inference: the likelihood is defined by the complete transition and observation law, not by the bound in isolation.

Conditional on the observations, the joint posterior over parameters, the state trajectory, and the regime path is
\begin{equation}\label{eq:4}
 \pi(\vartheta,x_{0:T},r_{1:T}\mid y_{1:T})
 \propto p(\vartheta)p(x_0\mid\vartheta)
 \prod_{t=1}^T p(r_t,x_t\mid x_{t-1},\vartheta)
 g_{r_t}(y_t\mid x_t,\vartheta).
\end{equation}
The factorization combines the parameter prior and the initial-state prior with the joint state-and-regime transition and the observation density at each date.

Interpreting \eqref{eq:4} as a density requires specifying a dominating measure: the reference measure against which the density is integrated to obtain probabilities. A deterministic state map need not induce a distribution with a density relative to Lebesgue measure. For example, if \(x_t\) is a deterministic function of \((x_{t-1},\epsilon_t)\) and the shock has fewer dimensions than the state, \(\dim\epsilon<\dim x\), then, given \(x_{t-1}\), the state \(x_t\) lies on a lower-dimensional manifold. This manifold has zero Lebesgue measure, yet the conditional state law assigns it all its probability. A Lebesgue density would assign zero probability to that set, so no transition density with respect to Lebesgue measure exists in this case. Even when the shock and state dimensions match, transforming the shock density into a state density requires a change-of-variables Jacobian.

The Markov transition kernel is therefore the appropriate primitive, rather than a transition density:
\begin{equation}\label{eq:4a}
 K_\vartheta(x_{t-1},A,r)
 =\int \mathbf 1\{\mathcal R(x_{t-1},\epsilon,\vartheta)=r\}\,
       \mathbf 1\{F_r(x_{t-1},\epsilon,\vartheta)\in A\}\,
       \phi_\vartheta(d\epsilon),
\end{equation}
where the first indicator selects shocks that generate the specified regime and the second selects shocks that place the next state in the specified set. Integrating their product under the shock law gives the probability of both events, conditional on the previous state and the parameters. This defines transition probabilities without presuming that the state law has a Lebesgue density.

The factor \(p(r_t,x_t\mid x_{t-1},\vartheta)\) in \eqref{eq:4} denotes a density of \(K_\vartheta\) with respect to an explicitly chosen dominating measure: Lebesgue measure in \(x_t\) times counting measure in \(r_t\) when the shock map is a smooth bijection onto its range. These measures account for the continuous state and the discrete regime, respectively. When the shock map is not a smooth bijection onto its range, the non-centered shock chart of \eqref{eq:56} is the natural default coordinate system. There the target is a density in \((\vartheta,x_0,\epsilon_{1:T})\): the coordinates are the parameters, the initial state, and the shocks. No transition density is needed at all.

The initial condition requires similar care in keeping the probability law consistent. The factor \(p(x_0\mid\vartheta)\) in \eqref{eq:4} and the first-step factor \(p(X_1\mid\vartheta)\) used by the particle recursion \eqref{eq:12a} must be reconciled through one of two conventions. Under the first, \(x_0\) is carried explicitly in every particle, and \eqref{eq:12a} begins from the prior \(p(x_0\mid\vartheta)\) followed by a first transition draw. Under the second, the recursion uses the induced initial law \(p(x_1\mid\vartheta)=\int K_\vartheta(x_0,\cdot)\,
p(x_0\mid\vartheta)\,dx_0\), with \(x_0\) integrated out. This induced law already averages the first transition over the initial-state prior. Both conventions are correct when used consistently; mixing them double-counts or drops the initial condition.

At fixed parameters, summing over the regimes and integrating over the state trajectory gives the marginal likelihood:
\begin{equation}\label{eq:5}
 p(y_{1:T}\mid\vartheta)=
 \int p(x_0\mid\vartheta)
 \prod_{t=1}^T\sum_{r_t}p(r_t,x_t\mid x_{t-1},\vartheta)
 g_{r_t}(y_t\mid x_t,\vartheta)\,dx_{0:T}.
\end{equation}
The regime summation in \eqref{eq:5}, or the corresponding nonlinear integral, is exactly the calculation that the piecewise Kalman filter avoids by conditioning on a verified regime path. Conditioning on that path is different from averaging over the possible paths. Keeping the summation in view distinguishes the exact nonlinear methods developed later from the linear benchmark in Section~\ref{sec:linear}.

\subsection{Four different posterior geometries}\label{sec:geometries}

The measure in \eqref{eq:4} determines the sampling problem; the economic word ``regime'' does not. Four cases must be kept separate because they call for different sampling choices. One of them requires further modeling rather than a sampler.

First, a censored rule such as \(i_t=\max\{\ell,i_t^\star\}\) is continuous. Its kink is a change in slope, not a jump in the rule's value. If all other model equations and densities are continuous, the posterior is continuous too, although its gradient may jump at the kink. Ordinary Metropolized HMC remains a valid MCMC construction when the nondifferentiable set has Lebesgue measure zero, but an integrator that ignores the kink can be inefficient (Afshar and Domke 2015, p.~1). Validity and efficiency are therefore distinct concerns here. Because the density is continuous, there is no potential-energy jump to refract across.

Second, suppose a deterministic solution map produces different one-sided density limits at a boundary. The discontinuity is now in the density itself, rather than only in its gradient. Here \(q\) collects all continuous unknowns, \(R(q)\) denotes the regime path they imply, and \(\pi_r(q)\) is the density obtained from branch \(r\). Every \(\pi_r\) must be expressed in one named coordinate system shared by the branches. If different branches are naturally parameterized in different coordinates, the change-of-variables Jacobian into this common chart belongs inside \(\pi_r\). With respect to Lebesgue measure in \(q\) and counting measure in \(r\), the joint density is
\begin{equation}\label{eq:5a}
 \Pi(q,r)=\pi_r(q)\,\mathbf 1\{r=R(q)\}.
\end{equation}
The indicator in \eqref{eq:5a} restricts support to the regime path implied by the continuous unknowns. At fixed \(q\), every proposal \(r'\ne R(q)\) therefore has zero target density. A regime-only Metropolis, mixed-HMC, or Gibbs move cannot cross this deterministic boundary, because changing only the regime would leave the target's support. The alternatives are to move \(q\) across the boundary with an event-aware continuous kernel, to make a joint \((q,r)\) proposal with its density and Jacobian included, or to integrate the states with a particle method.

Event-driven dynamics must account for the one-sided potential jump:
\begin{equation}\label{eq:5b}
 \Delta U(q_b)=-\log \pi_{r^+}(q_b^+)+\log \pi_{r^-}(q_b^-),
\end{equation}
where \(q_b^-\) and \(q_b^+\) denote limits as the boundary is approached from the two regions. The formula measures the change in negative log density across the boundary, with both branch densities expressed in the common chart. For a pure continuous \(\max\), the one-sided density limits agree and \(\Delta U=0\), recovering the first case.

Third, suppose the statistical model assigns positive probabilities to both regimes, say \(p(r_t\mid r_{t-1},x_{t-1},\vartheta)>0\) for both values. The posterior then genuinely has mixed continuous and discrete support. In contrast to the deterministic-boundary case, a discrete move can have positive target probability, so mixed HMC and particle Gibbs are valid candidates.

Fourth, if a nonlinear equilibrium solver can return several solutions, an algorithmic tie-break such as ``first root found'' depends on solver order. It does not supply probabilities for the competing solutions. The target therefore remains undefined until the model supplies a solution-selection distribution. This case requires further modeling: no sampler can draw from a distribution that has not been defined. The New Keynesian case study in Section~\ref{sec:nk} returns to this solution-selection issue at length.

\section{What OccBin and piecewise Kalman filtering actually do}\label{sec:linear}

Before developing exact nonlinear methods, the survey sets out the linear benchmark in detail. Understanding what the standard piecewise-linear machinery computes and where its two Gaussian closures enter makes it possible to state exactly what a nonlinear method adds.

\subsection{The two-regime linear system}\label{sec:linear-tworegime}

Guerrieri and Iacoviello (2015) represent an occasionally binding constraint with a reference regime \(n\) and an alternative regime \(b\). The regime labels distinguish the slack and binding cases. Linearizing both regimes around a common steady state gives a linear rational-expectations system for each. A candidate regime path \(r_{1:T}\) specifies which regime applies at each date and hence which coefficients enter the reduced state-space equations. The transition equation propagates the previous state and adds a regime-specific intercept and a transformed Gaussian shock:
\begin{equation}\label{eq:6}
 x_t=T_{r_t}x_{t-1}+C_{r_t}+R_{r_t}\epsilon_t,
 \qquad \epsilon_t\sim N(0,Q),
\end{equation}
The shock has mean zero and a covariance common to the regimes; the regime-specific loading matrix determines its effect on the state. The observation equation maps the current state to the data and adds an intercept and a Gaussian measurement disturbance:
\begin{equation}\label{eq:7}
 y_t=Z_{r_t}x_t+D_{r_t}+u_t,
 \qquad u_t\sim N(0,H_{r_t}).
\end{equation}
The measurement disturbance also has mean zero, but its covariance is regime-specific. The matrices in these equations come from solving each regime's equations as a linear rational-expectations system. A candidate path is accepted only when the implied binding and slack inequalities agree with its regime labels. Splicing arbitrary policy rules together would skip the expectations consistency that this solution step enforces.

For a finite horizon, OccBin starts with a guessed regime path and imposes the reference regime forever from a terminal date onward. It solves backward from that date, then simulates forward. The simulated trajectory is checked against the binding and slack inequalities, and the regime path is updated accordingly. If repeating these steps converges, the result is a fixed point: the path used in the solution agrees with the path implied by the simulation. Multiple fixed points can exist, so convergence of a particular iteration does not establish global uniqueness (Guerrieri and Iacoviello 2015, pp.~2--7).

\subsection{The conditional Kalman likelihood}\label{sec:linear-kalman}

Conditional on a verified regime path, the model \eqref{eq:6}--\eqref{eq:7} is time-varying but Gaussian. Its coefficients change with the prescribed regime, while its conditional linear-Gaussian structure makes the likelihood available in closed form through the Kalman recursions. The filtered mean \(a_{t-1}\) and covariance \(P_{t-1}\) summarize the state distribution after incorporating the preceding date's observations. Prediction propagates these through the current regime's transition equation. The mean uses the transition matrix and intercept; the covariance combines propagated state uncertainty with the covariance contributed by the new shock:
\begin{equation}\label{eq:8}
 a_t^- =T_{r_t}a_{t-1}+C_{r_t},
\qquad
 P_t^- =T_{r_t}P_{t-1}T_{r_t}^{\mathsf T}+R_{r_t}QR_{r_t}^{\mathsf T}.
\end{equation}
The innovation is the observed data minus the predicted observation, including its intercept. Its covariance combines predicted state uncertainty mapped into observation space with the measurement-disturbance covariance:
\begin{equation}\label{eq:9}
 v_t=y_t-Z_{r_t}a_t^- -D_{r_t},
 \qquad
 F_t=Z_{r_t}P_t^-Z_{r_t}^{\mathsf T}+H_{r_t}.
\end{equation}
The Kalman gain uses the covariance between the state and the observation, multiplied by the inverse innovation covariance, to weight the innovation. The update adds this weighted innovation to the predicted state mean and subtracts the uncertainty resolved by the observation from the predicted covariance:
\begin{equation}\label{eq:10}
 K_t=P_t^-Z_{r_t}^{\mathsf T}F_t^{-1},\qquad
 a_t=a_t^-+K_tv_t,\qquad
 P_t=P_t^- -K_tF_tK_t^{\mathsf T}.
\end{equation}
For the observed components \(\mathcal O_t\), summing the Gaussian innovation log densities over time gives the log likelihood conditional on the regime path and parameters:
\begin{equation}\label{eq:11}
 \log p(y_{1:T}\mid r_{1:T},\vartheta)
 =-\frac12\sum_{t=1}^T
 \left[|\mathcal O_t|\log(2\pi)+\log|F_t|+v_t^{\mathsf T}F_t^{-1}v_t\right].
\end{equation}
The first term supplies the Gaussian normalization for the number of observed components. The determinant term accounts for the scale of innovation uncertainty, and the quadratic term measures the innovation relative to that uncertainty.

Giovannini, Pfeiffer, and Ratto (2021, pp.~5--11) make the regime path endogenous to the filter: the filtering procedure also revises the path on which its recursions depend. Their piecewise Kalman filter (PKF) first performs the prediction, innovation, and update recursions \eqref{eq:8}--\eqref{eq:10}. It then estimates shocks and smoothed states, calls an OccBin simulation to infer a new regime path, replaces the transition matrices accordingly, and repeats the procedure. If this fixed-point iteration fails to converge, the parameter proposal is treated as a bad proposal.

The publicly available Dynare implementation follows the same structure. A standard Kalman update is followed by an OccBin solve and a new regime history; the loop continues until the proposed and implied paths agree. This implementation detail corroborates the described procedure but does not substitute for the paper's derivation.

\subsection{Exactness conditional on a linearization}\label{sec:linear-limits}

The PKF is exact for its conditional time-varying linear-Gaussian model. This exactness is conditional in a strong sense because the model itself is a piecewise-linear approximation. That approximation omits precautionary effects and higher-order nonlinearities that alter the probability of future binding (Giovannini, Pfeiffer, and Ratto 2021, pp.~3--4, 26~ff.). OccBin also needs a reference regime, a finite path iteration, and a rule for resolving multiple fixed points. These are model and solver choices, not mere numerical details.

For nonlinear estimation, the appropriate role of the PKF is therefore a \emph{limiting benchmark}. For this comparison, the nonlinear model is replaced by its verified piecewise-linear local representation. A correct nonlinear implementation applied to that representation should reproduce the conditional log likelihood \eqref{eq:11}, the regime path, and the OccBin inequality checks. Passing this comparison rules out a large class of bookkeeping errors shared by the two implementations. Validating the nonlinear sampler on the model actually being estimated remains a separate task, taken up in the sections that follow.

\section{A nonlinear likelihood without regime smoothing}
\label{sec:particle}

The task is to evaluate the marginal likelihood \eqref{eq:5} of the genuinely nonlinear model \eqref{eq:1} without smoothing away the regime. The discussion proceeds in three stages: a generic particle filter, the conditionally optimal particle filter built for piecewise-linear ZLB dynamics, and a family of branch-aware sigma-point constructions. The last family sits between a single Gaussian filter and a fully particle-based filter.

\subsection{Generic particle filtering}
\label{sec:pf-generic}

Particle filtering approximates the sequence of filtering distributions by
updating a weighted collection of sampled states, or particles. As a
by-product, it yields an unbiased estimator of the likelihood. This
unbiasedness underlies all the particle MCMC methods in
Section~\ref{sec:pmcmc}. For the nonlinear transition in
\eqref{eq:1}, the filter includes the regime in the particle
state when the regime is stochastic. When the regime is deterministic, the
filter computes it from each proposed state and shock.

The proposal densities $q_1(x_1\mid y_1,\vartheta)$ and
$q_t(x_t\mid x_{t-1},y_t,\vartheta)$ specify how the initial and subsequent
states are sampled. At $t=1$, the filter draws
$X_1^j\sim q_1$ and calculates an importance weight: the
prior state density times the observation density, divided by the proposal
density. Dividing each importance weight by the sum of these weights gives
the normalized filtering weights:
\begin{equation}\label{eq:12a}
 v_1^j=\frac{p(X_1^j\mid\vartheta)\,g(y_1\mid X_1^j,\vartheta)}
 {q_1(X_1^j\mid y_1,\vartheta)},\qquad
 W_1^j=\frac{v_1^j}{\sum_k v_1^k}.
\end{equation}
For $t\geq2$, the filter first resamples an ancestor
index $A_t^j\sim\operatorname{Categorical}(W_{t-1}^{1:M})$, selecting a previous particle
according to its normalized weight. It then draws
$X_t^j\sim q_t(\cdot\mid X_{t-1}^{A_t^j},y_t,\vartheta)$, conditional on the selected ancestor
and the current observation. The incremental importance weight now uses the
transition density from that ancestor in place of the initial state density.
The proposal correction and normalization are
\begin{equation}\label{eq:12b}
 v_t^j=\frac{p(X_t^j\mid X_{t-1}^{A_t^j},\vartheta)\,
                  g(y_t\mid X_t^j,\vartheta)}
 {q_t(X_t^j\mid X_{t-1}^{A_t^j},y_t,\vartheta)},\qquad
 W_t^j=\frac{v_t^j}{\sum_kv_t^k}.
\end{equation}
The resample-propagate likelihood estimator multiplies the mean unnormalized
weight from each period:
\begin{equation}\label{eq:13}
 \widehat p(y_{1:T}\mid\vartheta)=
 \prod_{t=1}^T\left(\frac1M\sum_{j=1}^M v_t^j\right).
\end{equation}

Unbiasedness of \eqref{eq:13} follows from a single
conditional-expectation calculation, iterated backward. Fixing the particle
system at time $t-1$ means conditioning on both
the particles $X_{t-1}^{1:M}$ and their normalized
weights $W_{t-1}^{1:M}$. With this system held fixed, each
incremental weight $v_t^j$ depends on two fresh
random draws: the ancestor $A_t^j$ and the
proposed state $X_t^j$. Conditional on a particular
ancestor, averaging over the proposed state cancels the proposal density
against the denominator of the importance weight:
\[
 \mathbb E\!\left[v_t^j \,\middle|\, A_t^j=k\right]
 =\int \frac{p(x\mid X_{t-1}^{k},\vartheta)\,g(y_t\mid x,\vartheta)}
            {q_t(x\mid X_{t-1}^{k},y_t,\vartheta)}\,
       q_t(x\mid X_{t-1}^{k},y_t,\vartheta)\,dx
 =p(y_t\mid X_{t-1}^{k},\vartheta),
\]
which leaves the one-step predictive density from particle
$k$. Averaging next over the categorical
ancestor draw, still conditional on the fixed previous particle system,
gives $\mathbb E[v_t^j]=\sum_k W_{t-1}^k\,p(y_t\mid X_{t-1}^{k},\vartheta)$. This weighted sum is exactly the
particle approximation to the predictive integral
$p(y_t\mid y_{1:t-1},\vartheta)$, whose integration uses the
time-$(t-1)$ filter.

In the likelihood product, multiplication by the preceding period's mean
unnormalized weight cancels the denominator in the normalized ancestor
weights. This leaves an average of unnormalized weights multiplied by the
predictive densities. The next conditional integration carries this
predictive factor inside the proposal integral. Repeating the integration
and cancellation backward to $t=1$ makes the
expected product telescope to the Feynman--Kac normalizing constant
$p(y_{1:T}\mid\vartheta)$, the likelihood of the full observation
sequence. Thus \eqref{eq:13} is nonnegative and unbiased
under the usual support and integrability conditions (Andrieu,
Doucet, and Holenstein 2010, Sec.~2.2; Aruoba et al. 2021, Sec.~6). Taking
logarithms does not preserve this unbiasedness. Because the logarithm is
concave, Jensen's inequality gives
$\mathbb E[\log\widehat p]\leq\log p$, so the log-likelihood estimate is biased
downward. Particle MCMC constructions must handle this bias at the level of
the acceptance ratio rather than by averaging log-likelihood estimates.

\subsection{The piecewise-linear conditionally optimal particle filter}
\label{sec:copf}

A generic bootstrap-style filter can be arbitrarily inefficient near the
bound because its blind proposals rarely land in the regime favored by the
data. The conditionally optimal particle filter (COPF) of Aruoba et al.
(2021) remedies this problem for piecewise-linear
ZLB dynamics by using the regime structure in the proposal itself. Their
canonical form uses continuous piecewise-linear decision rules,
\begin{equation}\label{eq:14}
 s_t=
 \begin{cases}
 \Phi_0(n)+\Phi_1(n)s_{t-1}+\Phi_\eta(n)\eta_t,&
 \eta_{1,t}<\zeta(s_{t-1}),\\[2pt]
 \Phi_0(b)+\Phi_1(b)s_{t-1}+\Phi_\eta(b)\eta_t,&
 \eta_{1,t}\geq\zeta(s_{t-1}),
 \end{cases}
\end{equation}
with a linear measurement equation. The first branch applies when the first
component of the innovation lies below the threshold determined by the
previous state; the second applies at or above that threshold. Conditional
on a previous particle, the observation density is then a mixture of two
truncated Gaussian integrals. Each integral covers the innovations
consistent with one branch of the decision rule.

Let $D_t^j(n)$ and
$D_t^j(b)$ denote the two regime contributions to
this conditional observation density. The probability assigned to the first
regime is its contribution divided by the total contribution:
\begin{equation}\label{eq:15}
 \lambda_t^j=\frac{D_t^j(n)}{D_t^j(n)+D_t^j(b)}.
\end{equation}
The conditionally optimal proposal selects the first regime with probability
$\lambda_t^j$ and the second with the remaining
probability. It then draws the structural innovation from the truncated
Gaussian corresponding to the selected regime and assigns the incremental
weight equal to the sum of both regime contributions:
\begin{equation}\label{eq:16}
 \widetilde w_t^j=D_t^j(n)+D_t^j(b).
\end{equation}

The proposal and its weight follow directly from Bayes' rule. Conditioning
the transition on the current observation gives the optimal proposal
\begin{equation}\label{eq:17}
 g_t^\star(s_t\mid y_t,s_{t-1})
 \propto p(y_t\mid s_t)\,p(s_t\mid s_{t-1}),
\end{equation}
so dividing the transition density times the observation density by this
normalized proposal cancels their common dependence on the proposed state.
The importance ratio is therefore the normalizing constant of the product
on the right-hand side of \eqref{eq:17}, namely the
predictive density $p(y_t\mid s_{t-1})$. This density is exactly
the sum of the regime contributions in
\eqref{eq:16}. With vanishing measurement error, one,
two, or no regime solutions may explain a given observation. If every
particle falls in the no-solution case, the likelihood estimate is zero
(Aruoba et al. 2021,
pp.~26--31).

This construction is the closest existing bridge to nonlinear ZLB inference
because it preserves an endogenous regime and avoids replacing the bound
by a softplus. Its limitations are equally important for choosing a sampler.
The transition remains a piecewise-linear approximation. The filter returns
a likelihood value but no parameter gradients, so an HMC kernel requires a
separate gradient construction. When vanishing measurement error allows two
regime solutions to explain the same observation, both contributions must
remain in the mixture \eqref{eq:16}. Dropping either
contribution changes the predictive density and hence the likelihood.

\subsection{A piecewise, truncated, and mixture UKF construction}
\label{sec:ukf}

The unscented Kalman filter (UKF) literature contains a useful intermediate family between a single Gaussian filter and a fully particle-based filter. Gaussian-sum filtering represents the distribution as a weighted sum of Gaussian components and updates each component with a Kalman-type step. It goes back to Alspach and Sorenson (1972). Within the sigma-point family, the truncated UKF conditions the Gaussian approximation on bounded-support measurement information, using the measurement bounds to restrict its support (Garcia-Fernandez, Morelande, and Grajal 2012a). Its mixture variant applies this truncated update separately to each component while retaining a Gaussian mixture (Garcia-Fernandez, Morelande, and Grajal 2012b). A UKF-aided Gaussian-sum filter combines the Gaussian-sum representation with these sigma-point and truncation ingredients (Gokce and Kuzuoglu 2015). These papers, together with Alspach and Sorenson (1972), are cited for orientation. The survey's authors could not obtain their full texts, so every equation in this section is derived in the survey's own notation rather than imported from them.

Two constrained-UKF papers that the authors studied in full complete this account. Kandepu, Imsland, and Foss (2008, Secs.~II--III) project transformed sigma points onto the feasible region. Teixeira et al. (2010, Secs.~4--5, Table~1) formulate interval-constrained unscented filtering, compare eight projection-, constraint-, and truncation-based UKF variants, and note that their motivating examples have multimodal densities. Together, these methods suggest a principled construction: retain the binding and nonbinding pieces of the distribution explicitly, and use a particle correction as the final step to restore exactness. None of the surveyed methods alone evaluates the exact nonlinear-ZLB likelihood, so preserving the branches does not remove the need for that correction.

\subsubsection{Why an ordinary UKF loses a branch}
\label{sec:ukf-branchloss}

To see where a single-Gaussian filter loses information at the bound, let $z_t$ collect the uncertain variables that determine the next regime. For example, $z_t=(x_{t-1}^{\mathsf T},\epsilon_t^{\mathsf T})^{\mathsf T}$ stacks the previous state and the current shock. Conditional on observations through the previous period, suppose this vector has the Gaussian predictive approximation
\begin{equation}\label{eq:18}
 z_t\mid y_{1:t-1}\approx N(m_t,P_t).
\end{equation}
Let $D_b$ and $D_n$ denote the binding and nonbinding regions. The switching statistic is the gap between the unconstrained policy rate and the lower bound, so its sign determines the regime:
\begin{equation}\label{eq:19}
 D_b=\{z:b_t(z)\leq0\},
 \qquad
 D_n=\{z:b_t(z)>0\},
 \qquad
 b_t(z)=i_t^\star(z)-\ell .
\end{equation}
The exact predictive state density adds the contributions from both regions. Within each region, it integrates the corresponding state-transition density against the predictive density of the uncertain inputs:
\begin{equation}\label{eq:20}
 p(x_t\mid y_{1:t-1})
 =\sum_{r\in\{b,n\}}
 \int_{D_r}p(x_t\mid z_t,r,\vartheta)\,
       p(z_t\mid y_{1:t-1})\,dz_t .
\end{equation}
Even when $p(z_t\mid y_{1:t-1})$ is Gaussian, the density in \eqref{eq:20} is generally not a single Gaussian. Conditioning on $D_r$ truncates the input by excluding values outside that regime. The branch maps $F_b$ and $F_n$ can then produce different means and covariances. The output is therefore a two-component mixture at best; neither branch component is guaranteed to be Gaussian.

A conventional UKF instead propagates a single sigma-point cloud through the branch function and replaces the transformed distribution by one Gaussian specified by its first two moments. This moment matching is a legitimate approximation in smooth problems. At the ZLB, however, it discards the separate branch masses, skewness, and multimodality precisely where the bound is informative. Evaluating the $\max$ function at every sigma point does not repair this loss: regardless of how the points are mapped, the final reduction to one mean and covariance still removes that information.

\subsubsection{Exact branch probabilities and truncated moments for an
affine boundary}
\label{sec:ukf-truncmoments}

The affine case provides the fully analytic building block for everything later in this section. The derivation first reduces the branch event to a scalar Gaussian truncation, then uses Gaussian conditioning to recover the full vector moments. Suppose the switching statistic is affine and the input is Gaussian, with positive variance in the direction that determines the branch:
\begin{equation}\label{eq:21}
 b_t(z)=a^{\mathsf T}z-c,
 \qquad z\sim N(m,P),
 \qquad \sigma_b^2=a^{\mathsf T}Pa>0 .
\end{equation}
The scalar $u=a^{\mathsf T}z$ then has distribution $N(\mu_b,\sigma_b^2)$, with $\mu_b=a^{\mathsf T}m$. Thus the branch event depends on $z$ only through this one Gaussian scalar. Expressing the distance from its mean to the threshold in standard-deviation units gives the following standardized cutoff and branch probabilities:
\begin{equation}\label{eq:22}
 \gamma=\frac{c-\mu_b}{\sigma_b},
 \qquad
 \alpha_b=\Pr(b_t(z)\leq0)=\Phi(\gamma),
 \qquad
 \alpha_n=\Pr(b_t(z)>0)=1-\Phi(\gamma),
\end{equation}
where $\Phi$ and $\phi$ are the standard normal cumulative distribution function and density, respectively. The binding probability is the normal mass below the standardized cutoff; the nonbinding probability is the complementary mass above it.

To derive the conditional moments, standardize the scalar as $v=(u-\mu_b)/\sigma_b$, so the binding event becomes $v\leq\gamma$. The elementary identity $\frac{d}{dv}\bigl[-\phi(v)\bigr]=v\,\phi(v)$ shows that the negative of $\phi$ is an antiderivative of the density multiplied by $v$. Direct integration over the binding region, divided by the probability of that region, therefore gives the conditional mean:
\begin{equation}\label{eq:23}
 \mathbb E[v\mid v\leq\gamma]
 =\frac{\int_{-\infty}^{\gamma}v\,\phi(v)\,dv}{\Phi(\gamma)}
 =\frac{\bigl[-\phi(v)\bigr]_{-\infty}^{\gamma}}{\Phi(\gamma)}
 =-\frac{\phi(\gamma)}{\Phi(\gamma)}.
\end{equation}
The denominator renormalizes the density after truncation, and the boundary term at negative infinity vanishes.

For the second moment, integration by parts starts from $v^2\phi(v)=v\cdot v\phi(v)$, using $-\phi$ as the antiderivative for the second factor:
\[
 \int_{-\infty}^{\gamma}v^2\phi(v)\,dv
 =\bigl[-v\,\phi(v)\bigr]_{-\infty}^{\gamma}
  +\int_{-\infty}^{\gamma}\phi(v)\,dv
 =\Phi(\gamma)-\gamma\,\phi(\gamma),
\]
so dividing this integral by the binding probability gives $\mathbb E[v^2\mid v\leq\gamma]=1-\gamma\lambda_b$, with $\lambda_b=\phi(\gamma)/\Phi(\gamma)$. Subtracting the squared conditional mean $\lambda_b^2$ from this second moment, using the mean in \eqref{eq:23}, gives the conditional variance:
\begin{equation}\label{eq:24}
 \operatorname{Var}(v\mid v\leq\gamma)
 =1-\lambda_b(\lambda_b+\gamma),
 \qquad
 \lambda_b=\frac{\phi(\gamma)}{\Phi(\gamma)}.
\end{equation}
For the nonbinding event $v>\gamma$, the same two calculations run over $(\gamma,\infty)$. Normalizing by the nonbinding probability and subtracting the squared mean for the variance gives
\begin{equation}\label{eq:25}
 \mathbb E[v\mid v>\gamma]=\lambda_n,
 \qquad
 \operatorname{Var}(v\mid v>\gamma)
 =1-\lambda_n(\lambda_n-\gamma),
 \qquad
 \lambda_n=\frac{\phi(\gamma)}{1-\Phi(\gamma)}.
\end{equation}
The change of integration region reverses the direction of the conditional mean shift.

These scalar moments can now be lifted back to the full vector $z$ through Gaussian regression. Since $\operatorname{Cov}(z,u)=Pa$, the linear regression of $z$ on $u$ has coefficient $Pa/\sigma_b^2$. The Gaussian conditioning identity therefore decomposes $z$ into the part that moves with the switching scalar and an independent residual:
\begin{equation}\label{eq:26}
 z=m+\frac{Pa}{\sigma_b^2}(u-\mu_b)+\xi,
 \qquad
 \xi\perp u,
 \qquad
 \operatorname{Var}(\xi)=P-\frac{Paa^{\mathsf T}P}{\sigma_b^2}.
\end{equation}
Conditioning on $\{v\leq\gamma\}$ or $\{v>\gamma\}$ truncates only the $u$-term in \eqref{eq:26}. The residual $\xi$ is independent of either event and retains its mean and covariance. Only the scalar term's moments therefore need to change. Substituting the scalar results \eqref{eq:23}--\eqref{eq:25} into \eqref{eq:26} gives the exact first two moments of each truncated branch. For the binding branch,
\begin{equation}\label{eq:27}
 m_b=m-\frac{Pa}{\sigma_b}\lambda_b,
 \qquad
 P_b=P-\frac{Paa^{\mathsf T}P}{\sigma_b^2}\,
             \lambda_b(\lambda_b+\gamma),
\end{equation}
and for the nonbinding branch,
\begin{equation}\label{eq:28}
 m_n=m+\frac{Pa}{\sigma_b}\lambda_n,
 \qquad
 P_n=P-\frac{Paa^{\mathsf T}P}{\sigma_b^2}\,
             \lambda_n(\lambda_n-\gamma).
\end{equation}
The conditional means shift in opposite directions along the regression direction, while each covariance adjusts for the reduction in scalar variance caused by truncation.

The branch probabilities in \eqref{eq:22} and the conditional means and covariances in \eqref{eq:27} and \eqref{eq:28} are exact for the original Gaussian law and its conditioning on each branch event. They follow from integrating that law, not from a heuristic projection. They are therefore the correct starting point for a branch-aware sigma-point filter whenever the switching statistic is affine in the augmented state and shock.

Implementing these formulas requires stable evaluation of the Mills ratios $\lambda_b=\phi(\gamma)/\Phi(\gamma)$ and $\lambda_n=\phi(\gamma)/(1-\Phi(\gamma))$. Direct evaluation from $\Phi$ can underflow catastrophically in the tails, where the numerator and denominator both vanish at very different rates. Implementations should therefore use log-cdf, log-survival, and log-density arithmetic, or an independently tested inverse-Mills-ratio routine. Any variance clamp added near $\gamma\to\pm\infty$ is a safeguard against roundoff; it cannot substitute for stable probability evaluation.

For a nonlinear switching statistic $b_t$, no corresponding general closed form exists. One approximation linearizes $b_t$ at the predictive mean $m$ and reuses \eqref{eq:21}--\eqref{eq:28} with $a=\nabla b_t(m)$ and $c=a^{\mathsf T}m-b_t(m)$. This replaces the curved boundary by a local affine approximation, whose error grows with the curvature of the boundary inside the predictive ellipsoid. A sigma-point implementation of this approximation also inherits a blind spot: a standard $2d+1$-point set concentrated around the predictive mean can entirely miss a low-probability binding region. The branch probability $\alpha_b$ estimated from those points can then be exactly zero even when its true value is small but positive. There are two defensible responses: split components along the estimated boundary normal, which refines the approximation, or apply the particle correction of Section~\ref{sec:ukf-proposal}, which restores exactness. Whichever response is used, the approximate and exact ingredients should remain clearly distinguished.

\subsubsection{The pure censored policy-rate likelihood}
\label{sec:ukf-censored}

A single censored Gaussian observation already has mixture structure without any UKF approximation. Working out this case completely shows how censoring creates that structure. The model consists of a Gaussian shadow rate, a policy rate censored at the lower bound, and an observed rate with additive Gaussian measurement error:
\begin{equation}\label{eq:29}
 i_t^\star\sim N(\mu,\sigma^2),
 \qquad i_t=\max\{\ell,i_t^\star\},
 \qquad y_t=i_t+u_t,
 \qquad u_t\sim N(0,V).
\end{equation}
The bound binds when the shadow rate lies at or below it. Standardizing the shadow rate gives the binding probability and its complementary nonbinding probability:
\begin{equation}\label{eq:30}
 \alpha_b=\Phi\!\left(\frac{\ell-\mu}{\sigma}\right),
 \qquad
 \alpha_n=1-\alpha_b.
\end{equation}
The censoring map $\max\{\ell,\cdot\}$ sends the entire lower tail of $i_t^\star$ to the single point $\ell$. Thus, with no measurement error, the distribution of $i_t$ has an atom, or point mass, of size $\alpha_b$ at $\ell$, together with a continuous Gaussian upper tail:
\begin{equation}\label{eq:31}
 p(i_t)=\alpha_b\,\delta_\ell(i_t)
 +\frac{1}{\sigma}\phi\!\left(\frac{i_t-\mu}{\sigma}\right)
  \mathbf 1\{i_t>\ell\}.
\end{equation}
The upper-tail term retains the original Gaussian density above the bound, so its total mass is the nonbinding probability.

Gaussian observation error often hides this atom in the observed data. The point mass remains in the censored rate's distribution; its contribution to the observed density is obtained by convolution with the noise. The derivation starts with the binding event, on which $i_t=\ell$ holds exactly. The observation is then the bound plus Gaussian noise, so weighting its density by the binding probability gives the contribution $\alpha_b\,N(y_t;\ell,V)$.

The nonbinding contribution requires integrating the observation density $N(y_t;i,V)$ against the upper Gaussian tail in \eqref{eq:31}. The standard product-of-Gaussians identity evaluates this integral. For two Gaussian densities in $i$, collecting the quadratic and linear terms in the exponent and completing the square gives another Gaussian in $i$ multiplied by a factor independent of $i$:
\begin{equation}\label{eq:32}
 N(y_t;i,V)\,N(i;\mu,\sigma^2)
 =N(y_t;\mu,V+\sigma^2)\,N(i;\widetilde\mu_t,\widetilde\sigma^2),
\end{equation}
where matching coefficients gives the variance and mean in the familiar precision-weighted form of a Gaussian update:
\begin{equation}\label{eq:33}
 \widetilde\sigma^2=\frac{\sigma^2V}{\sigma^2+V},
 \qquad
 \widetilde\mu_t=\frac{V\mu+\sigma^2y_t}{\sigma^2+V}.
\end{equation}
The factor $N(y_t;\mu,V+\sigma^2)$ is independent of the integration variable and can be taken outside the integral. What remains is the probability mass that the Gaussian $N(\widetilde\mu_t,\widetilde\sigma^2)$ assigns to $(\ell,\infty)$. Integrating \eqref{eq:32} over $i>\ell$ and adding the binding contribution therefore gives the complete predictive density:
\begin{equation}\label{eq:34}
 p(y_t)=
 \alpha_b\,N(y_t;\ell,V)
 +N(y_t;\mu,V+\sigma^2)
  \left[1-\Phi\!\left(
   \frac{\ell-\widetilde\mu_t}{\widetilde\sigma}\right)\right].
\end{equation}
The bracketed term is the upper-tail probability from the nonbinding integral. By Bayes' rule, the posterior probability that a given observation came from the binding branch is the binding contribution divided by the complete predictive density:
\begin{equation}\label{eq:35}
 \Pr(b\mid y_t)=
 \frac{\alpha_b\,N(y_t;\ell,V)}{p(y_t)}.
\end{equation}

Equation \eqref{eq:34} is the one-period analogue of the two-regime truncated-Gaussian mixture underlying the COPF in Section~\ref{sec:copf}. It also shows why replacing the max censoring map with a softplus changes more than the numerical evaluation: the softplus removes the atom in \eqref{eq:31}, which changes the predictive density in \eqref{eq:34} and hence the likelihood itself. In a DSGE model, the other observables and the state transition replace the scalar Gaussian factors in \eqref{eq:34}. The decomposition into binding and nonbinding branches remains the same.

\subsubsection{Branchwise UKF update}
\label{sec:ukf-branchwise}

A full filtering step that carries both branches forward can now be assembled from the affine building block in Section~\ref{sec:ukf-truncmoments}. The construction assumes that, before the next observation, the filtering approximation is a Gaussian mixture with weights summing to one:
\begin{equation}\label{eq:36}
 p(z_t\mid y_{1:t-1})
 \approx\sum_{j=1}^{J_{t-1}}\omega_{t-1,j}\,
 N(z_t;m_{t-1,j},P_{t-1,j}),
 \qquad \sum_j\omega_{t-1,j}=1.
\end{equation}
For each mixture component $j$, the calculations in \eqref{eq:22}, \eqref{eq:27}, and \eqref{eq:28} give a component $(j,r)$ for each branch. Its prior mass is the original component weight multiplied by the probability of that branch within the component:
\begin{equation}\label{eq:37}
 \rho_{j,r}=\omega_{t-1,j}\,\alpha_{j,r},
 \qquad r\in\{b,n\}.
\end{equation}
For each branch component, the branch-conditioned moments $(m_{j,r},P_{j,r})$ provide the mean and covariance used to form sigma points $\chi_{j,r}^{(k)}$ with mean and covariance weights $w_m^{(k)},w_c^{(k)}$,
and propagate each branch-conditioned sigma point through its branch's transition map, then apply the same branch's observation map to the resulting state:
\begin{equation}\label{eq:38}
 X_{j,r}^{(k)}=F_r(\chi_{j,r}^{(k)},\vartheta),
 \qquad
 Y_{j,r}^{(k)}=h_r(X_{j,r}^{(k)},\vartheta).
\end{equation}
If process and observation noises are included in the augmented sigma point, their covariances are already accounted for, so the separate noise-covariance additions below must be omitted. Otherwise, the prediction adds $Q_r$ to the state covariance and $V_r$ to the observation covariance.

For each branch component, the mean weights give the predicted state mean as a weighted sum of the transitioned sigma points. The covariance weights combine the outer products of deviations from that mean to give the predicted state covariance:
\begin{equation}\label{eq:39}
 \bar x_{j,r}=\sum_k w_m^{(k)}X_{j,r}^{(k)},
 \qquad
 P^x_{j,r}=\sum_k w_c^{(k)}
 (X_{j,r}^{(k)}-\bar x_{j,r})(X_{j,r}^{(k)}-\bar x_{j,r})^{\mathsf T}+Q_r,
\end{equation}
Applying the same calculations to the observation sigma points gives the predicted observation mean and covariance, with the noise contribution handled as above:
\begin{equation}\label{eq:40}
 \bar y_{j,r}=\sum_k w_m^{(k)}Y_{j,r}^{(k)},
\end{equation}
\begin{equation}\label{eq:41}
 S_{j,r}=\sum_k w_c^{(k)}
 (Y_{j,r}^{(k)}-\bar y_{j,r})(Y_{j,r}^{(k)}-\bar y_{j,r})^{\mathsf T}+V_r,
\end{equation}
The cross-covariance pairs each state deviation with the observation deviation generated from the same sigma point:
\begin{equation}\label{eq:42}
 C_{j,r}=\sum_k w_c^{(k)}
 (X_{j,r}^{(k)}-\bar x_{j,r})(Y_{j,r}^{(k)}-\bar y_{j,r})^{\mathsf T}.
\end{equation}
The Gaussian UKF predictive likelihood evaluates the observed data under a Gaussian distribution with the predicted observation mean and covariance. The branch gain combines the cross-covariance with the inverse observation covariance:
\begin{equation}\label{eq:43}
 \widehat L_{j,r}=N(y_t;\bar y_{j,r},S_{j,r}),
 \qquad K_{j,r}=C_{j,r}S_{j,r}^{-1},
\end{equation}
The component update adds the gain-weighted observation prediction error to the predicted state mean. Its covariance update subtracts the observation covariance transformed by the gain:
\begin{equation}\label{eq:44}
 m^+_{j,r}=\bar x_{j,r}+K_{j,r}(y_t-\bar y_{j,r}),
 \qquad
 P^+_{j,r}=P^x_{j,r}-K_{j,r}S_{j,r}K_{j,r}^{\mathsf T}.
\end{equation}
Bayes' rule then updates the branch weights. Multiplying each prior branch mass by its predictive likelihood gives an unnormalized weight. Summing these weights over every prior mixture component and both the binding and nonbinding branches gives the evidence normalizer. Dividing each unnormalized weight by this sum gives its posterior weight:
\begin{equation}\label{eq:45}
 \widetilde\omega_{j,r}=\rho_{j,r}\widehat L_{j,r},
 \qquad
 Z_t=\sum_{j=1}^{J_{t-1}}\sum_{r\in\{b,n\}}\widetilde\omega_{j,r},
 \qquad
 \omega^+_{j,r}=\frac{\widetilde\omega_{j,r}}{Z_t}.
\end{equation}
After observing $y_t$, the branch-aware UKF approximation is therefore a Gaussian mixture with the updated component means, covariances, and posterior weights:
\begin{equation}\label{eq:46}
 p(x_t\mid y_{1:t})
 \approx\sum_{j,r}\omega^+_{j,r}\,
 N(x_t;m^+_{j,r},P^+_{j,r}).
\end{equation}

If the mixture in \eqref{eq:46} is to be replaced by a single Gaussian, a moment-preserving collapse matches its mean and covariance. The collapsed mean averages the component means using their posterior weights:
\begin{equation}\label{eq:47}
 \bar m_t=\sum_{j,r}\omega^+_{j,r}m^+_{j,r},
\end{equation}
The collapsed covariance includes both the component covariances and the dispersion of the component means around the collapsed mean:
\begin{equation}\label{eq:48}
 \bar P_t=\sum_{j,r}\omega^+_{j,r}
 \left[P^+_{j,r}+
 (m^+_{j,r}-\bar m_t)(m^+_{j,r}-\bar m_t)^{\mathsf T}\right].
\end{equation}
The second term in \eqref{eq:48} is the between-branch variance, contributed by the deviations of the component means from the collapsed mean. This contribution matters for prediction: dropping it makes the next period's threshold probability too small or too large whenever the binding and nonbinding components are well separated. Retaining all components instead lets the mixture grow to as many as $2^t$ histories, the classical Gaussian-sum growth problem. Practical filters limit this growth through pruning, merging, or an IMM-style moment collapse. Each operation changes the filtering density, so its effect belongs in the run's error budget alongside the moment-matching step itself.

\subsubsection{Relation to IMM-UKF and to deterministic ZLB support}
\label{sec:ukf-imm}

It is natural to ask how the branchwise construction relates to classical multiple-model filters. The answer depends critically on whether the regime is genuinely stochastic. For a stochastic Markov regime, the construction becomes a multiple-model UKF. If $p(r_t=r\mid r_{t-1}=s)=\Pi_{sr}$, the mode-prediction probability sums the previous mode probabilities weighted by their transition probabilities into the destination mode. Dividing each contribution by that sum gives the conditional probability of the previous mode given the destination mode:
\begin{equation}\label{eq:49}
 c_{t,r}=\sum_s\Pi_{sr}\mu_{t-1,s},
 \qquad
 \mu_{t-1,s\mid r}=\frac{\Pi_{sr}\mu_{t-1,s}}{c_{t,r}},
\end{equation}
These conditional weights form the IMM mixing mean and covariance supplied to the $r$-specific UKF. The mixing mean averages the previous modes' state means:
\begin{equation}\label{eq:50}
 m_{t-1\mid r}=\sum_s\mu_{t-1,s\mid r}m_{t-1,s},
\end{equation}
The mixing covariance includes both the weighted state covariances within those modes and the weighted dispersion of their means around the mixing mean:
\begin{equation}\label{eq:51}
 P_{t-1\mid r}=\sum_s\mu_{t-1,s\mid r}
 \left[P_{t-1,s}+
 (m_{t-1,s}-m_{t-1\mid r})(m_{t-1,s}-m_{t-1\mid r})^{\mathsf T}\right].
\end{equation}
After the branchwise update, the mode probability is proportional to $c_{t,r}\widehat L_{t,r}$, the product of the predicted mode probability and that mode's predictive likelihood. Equation \eqref{eq:49} expresses Bayes' rule on the mode chain. Equations \eqref{eq:50}--\eqref{eq:51} apply the moment-preserving collapse \eqref{eq:47}--\eqref{eq:48} within each destination mode. This is the interacting-multiple-model (IMM) logic introduced for Markovian switching systems by Blom and Bar-Shalom (1988), here with UKF components replacing linear Kalman components. A constrained multiple-model UKF examined by the authors applies the same Markov-jump structure with truncated noise supports and model-conditioned feasible regions (Zhang et al. 2020, Secs.~II--III).

The deterministic ZLB differs precisely because of the support restriction in \eqref{eq:5a}. There is no free $\Pi_{sr}$ to tune: the branch is determined by $r_t=R(z_t)$, and at a fixed $z_t$ the other label has zero support. Equations \eqref{eq:49}--\eqref{eq:51} are therefore a valid approximation only if a stochastic regime law is deliberately introduced. Introducing that law changes the regime specification. For the original complementarity model, $\alpha_{j,r}$ in \eqref{eq:37} must instead be computed from the continuous predictive law. The regime label must remain a quantity derived from the continuous variables rather than a free state with its own transition probabilities.

\subsubsection{UKF proposals inside a corrected particle filter}
\label{sec:ukf-proposal}

Every construction in Section~\ref{sec:ukf} up to this point approximates the likelihood. A more defensible role for the branchwise UKF is to propose particles, with importance weighting correcting the proposal rather than letting the UKF define the likelihood by itself. Even a crude proposal can be used because the weights correct the discrepancy between the proposal and the target. Let $q_{t,j,r}(z_t\mid y_t)$ be a branch-conditioned UKF proposal supported on $D_r$, and let $\alpha_{j,r}$ be the branch-selection probability. The algorithm first draws branch $r$ with probability $\alpha_{j,r}$, then draws $z_t\sim q_{t,j,r}$. For a stochastic regime, the target contribution is the product of the regime-transition probability, the conditional transition density, and the observation density. Dividing this product by the branch-selection probability and the conditional proposal density gives the incremental importance weight:
\begin{equation}\label{eq:52}
 w_t=
 \frac{p(r_t\mid r_{t-1},x_{t-1},\vartheta)\,
       p(z_t\mid x_{t-1},r_t,\vartheta)\,
       g_{r_t}(y_t\mid x_t,\vartheta)}
      {\alpha_{j,r_t}\,
       q_{t,j,r_t}(z_t\mid x_{t-1},y_t,\vartheta)}.
\end{equation}
For a deterministic regime, $r_t=R(z_t)$ is determined by the augmented state and is not an independently sampled variable under the model. If the shock-augmented transition is used, the state follows from the branch-specific transition map. The weight numerator then contains the exact shock density and the observation density, without a separate regime-transition probability. The branch-selection probability remains in the denominator because selecting a branch is still part of the proposal procedure. The corresponding weight and state update are
\begin{equation}\label{eq:53}
 w_t=
 \frac{p(\epsilon_t\mid\vartheta)\,
       g_{R(z_t)}(y_t\mid x_t,\vartheta)}
      {\alpha_{j,R(z_t)}\,
       q_{t,j,R(z_t)}(z_t\mid x_{t-1},y_t,\vartheta)},
 \qquad x_t=F_{R(z_t)}(x_{t-1},\epsilon_t,\vartheta).
\end{equation}
In both cases, the exact numerator and the proposal density determine $w_t$. The UKF's Gaussian moment approximation enters only through the proposal, so the importance ratio corrects its approximation error.

The density ratio in \eqref{eq:53} requires an explicit coordinate convention. Its numerator density is written in the shock coordinate $\epsilon_t$, whereas its proposal density is written in the augmented coordinate $z_t$. These densities are consistent only under one of two conventions. Under convention (a), the proposal draws the shock directly, so $q$ is a density in $\epsilon_t$. With $x_{t-1}$ fixed, the map $\epsilon_t\mapsto z_t$ is an affine embedding, and both densities live in the same shock-coordinate chart. Under convention (b), the proposal density in $z_t$ must be converted to the $\epsilon_t$ chart using the absolute Jacobian of the transformation. Convention (a) is recommended: propose shocks, then derive states. Under either convention, and under the usual support and integrability conditions, inserting \eqref{eq:52} or \eqref{eq:53} into the particle estimator \eqref{eq:13} preserves the nonnegative unbiased likelihood property of Section~\ref{sec:pf-generic}. A UKF-only filter has no such importance-weight correction and retains an approximate Gaussian-sum likelihood.

The iterative truncated unscented particle filter of Wang et al. (2020, Secs.~3.1--3.2, Algorithm~1) provides a precedent for this use of a UKF to construct particle proposals. It first truncates an iterated-UKF update to the constraint region. The first two moments of the truncated density define a Gaussian proposal, and each particle is reweighted by the ratio of target factors to the proposal density. Two differences from \eqref{eq:52}--\eqref{eq:53} matter for the ZLB. First, Wang et al.'s constraint region is exogenous, whereas the model itself selects the ZLB branch. This endogenous selection explains why the branch probability $\alpha_{j,r}$ and the branch-specific transition and observation maps $F_r,h_r$ enter the weights above. Second, their weight numerator uses a Gaussian approximation of the predictive density (their eq.~36) rather than the exact transition density. Their filter is therefore a benchmark for proposal quality, not an exactly weighted likelihood estimator. For a filter's likelihood estimate to be trusted as unbiased, its weight numerator must retain the exact shock density, as in \eqref{eq:53}.

These constructions establish a hierarchy of roles. A split or truncated mixture UKF is a fast, useful diagnostic and a good source of particle proposals. A branch-corrected particle filter provides the authoritative nonlinear likelihood. The COPF construction of Section~\ref{sec:copf} is the exact piecewise-linear benchmark. Correcting the likelihood in this way does not make the UKF proposal an exact source of HMC gradients. An HMC scheme must target either the corrected extended particle state or an explicitly acknowledged approximation (Sections~\ref{sec:hmc} and~\ref{sec:pmcmc}). The shadow-rate case study of Section~\ref{sec:shadow} shows how these constructions apply when switching occurs in the observation map rather than in the transition.

\section{Ordinary HMC and how the lower bound changes its geometry}\label{sec:hmc}

Hamiltonian Monte Carlo is the natural starting point for posterior sampling in high-dimensional latent-variable models. The organizing question for the rest of this survey is precisely where an occasionally binding constraint breaks its assumptions. For a smooth posterior $\pi(q)\propto e^{-U(q)}$, the potential is the negative log posterior up to a constant. HMC introduces momentum $p\sim e^{-K(p)}$ and forms the Hamiltonian $H=U+K$ by adding potential and kinetic energy. The joint position and momentum, or phase point, then evolve according to Hamilton's equations (Neal 2011, Sec.~5.2)
\begin{equation}\label{eq:54}
 \dot q=\nabla_pK(p),\qquad \dot p=-\nabla_qU(q).
\end{equation}
For Gaussian momentum, the kinetic energy is $K(p)=\frac12p^{\mathsf T}M^{-1}p$. The leapfrog integrator discretizes \eqref{eq:54} through a half-step momentum update, a full-step position update, and a second half-step momentum update:
\begin{equation}\label{eq:55}
 p\leftarrow p-\tfrac\epsilon2\nabla U(q),\quad
 q\leftarrow q+\epsilon M^{-1}p,\quad
 p\leftarrow p-\tfrac\epsilon2\nabla U(q).
\end{equation}
The second momentum update uses the gradient at the updated position. The leapfrog map is reversible and volume preserving. These properties allow a final Metropolis correction, with acceptance probability $1\wedge e^{-\Delta H}$, to account exactly for the integration error.

An attractive way to use HMC in a nonlinear state-space model is to sample the parameters and latent shocks jointly, avoiding the filter entirely. Childers et al.\ (2022, pp.~10--16) pursue this approach. For the state equation $x_t=F(x_{t-1},\epsilon_t,\vartheta)$, the state path is generated recursively from the initial state and shocks. The joint posterior is proportional to the product of the parameter prior, the conditional initial-state density, and each period's shock and observation densities. Taking logarithms gives, up to a constant,
\begin{equation}\label{eq:56}
 \log p(\vartheta,x_0,\epsilon_{1:T}\mid y)
 =\log p(\vartheta)+\log p(x_0\mid\vartheta)
 +\sum_{t=1}^T\left[\log p(\epsilon_t\mid\vartheta)
 +\log g(y_t\mid x_t,\vartheta)\right].
\end{equation}
This construction can remove a costly nonlinear filter, but HMC must then differentiate the complete simulation map $F$ composed over $T$ periods. The lower bound complicates that derivative calculation. A hard $\max$ is only piecewise differentiable: automatic differentiation supplies the derivative of the active branch, but reports neither a boundary event nor information about the kink. A branch-dependent equilibrium solver may also jump or be multi-valued across the boundary. Within a particle filter, discrete resampling changes ancestor indices discontinuously. Replacing the bound with a softplus restores smoothness but removes the censoring atom, the probability mass at the bound. It therefore changes the likelihood itself (Section~\ref{sec:ukf-censored} computes this change exactly) and defines a separate approximate model rather than a smooth route to the constrained model.

\subsection{Validity of Metropolized HMC on a kinked target}\label{sec:kink-validity}

A kink means that the potential is continuous while its gradient is discontinuous. The issue is what this loss of smoothness costs ordinary leapfrog HMC. In the two-country shadow-rate case study of Section~\ref{sec:shadow}, the target is continuous with a kinked gradient rather than discontinuous. The validity of ordinary leapfrog HMC on such a target therefore needs a careful statement. The authors state and sketch the following proposition because they could not locate a reference at this level of generality. Its assumptions are chosen to cover a continuous piecewise-quadratic-plus-smooth potential, the form produced by a piecewise-affine observation map inside a Gaussian likelihood.

\paragraph{Assumptions.}
The potential $U$ is continuous on $\mathbb R^d$. There is a closed Lebesgue-null set $N$ arising from a finite (or locally finite) piecewise-regular partition of $\mathbb R^d$. For the target of Section~\ref{sec:shadow}, this set is a finite union of hyperplanes. The potential $U$ is continuously differentiable on $\mathbb R^d\setminus N$, with $\nabla U$ locally bounded. On each piece, the branch gradient is deterministic and single-valued, and a fixed boundary convention assigns its value on $N$. Finally, $e^{-U}$ is integrable. These assumptions are \emph{not} guaranteed by almost-everywhere differentiability alone. A potential evaluated through an implicit multi-root equilibrium solver, through an algorithm with a state-dependent iteration count, or under any branch rule that is not a fixed partition of $\mathbb R^d$ lies outside this proposition until a separate theorem covers it.

\paragraph{Claim.}
Fix a step size $\epsilon$ and a step count $L$. Under the stated assumptions, the set of initial phase points $(q,p)$ whose leapfrog trajectory \eqref{eq:55} evaluates $\nabla U$ at a point of $N$ has Lebesgue measure zero. Outside this exceptional set, the $L$-step leapfrog map is well defined, volume preserving, and reversible in the usual momentum-flip sense. Flipping the momentum at the endpoint and applying the same steps retraces the trajectory with reversed momentum. Consequently, the Metropolized chain that accepts with probability $1\wedge e^{-\Delta H}$ leaves $\pi\propto e^{-U}$ invariant.

\emph{Sketch.} Each half-kick $p\mapsto p-\tfrac\epsilon2\nabla U(q)$ at fixed $q\notin N$ translates momentum while holding position fixed. Each drift $q\mapsto q+\epsilon M^{-1}p$ translates position while holding momentum fixed. Both updates are measure-preserving shears. At a given substep, the starting points that land on $N$ form the preimage of a null set under a composition of measure-preserving maps defined off a null set. This preimage is therefore null. Taking the finite union over the $2L+1$ substeps still gives a null set. On the remaining full-measure set, reversibility and volume preservation follow from the standard shear argument (Neal 2011, Sec.~5.2), which does not require smoothness across $N$. The piecewise-regular partition assumption makes the shear composition well defined off a null set. The argument is \emph{not} valid for an arbitrary nonsmooth composition merely because its individual shears preserve measure. $\square$

A sampler can target the correct posterior and still sample inefficiently. When a drift substep crosses a kink, the usual $O(\epsilon^3)$ local energy error is replaced, at the crossing, by
\begin{equation}\label{eq:56a}
 \Delta H \;=\; O\!\left(\epsilon\,
 \bigl\|\nabla U^{+}-\nabla U^{-}\bigr\|\right)
\end{equation}
where $\nabla U^{\pm}$ are the one-sided gradients of the potential on opposite sides of the kink. The crossing error depends on both the step size and the gradient jump. More frequent crossings and larger gradient jumps therefore degrade acceptance. This loss of efficiency is what Afshar and Domke (2015, p.~1) observed; it does not make the sampler invalid.

Two consequences follow for integrator choice. First, the reflection and refraction machinery of Section~\ref{sec:event} degenerates at a kink. The potential jump in \eqref{eq:58} is $\Delta U=0$, so \eqref{eq:59} leaves the momentum unchanged. Detecting the boundary can improve integration accuracy, but it is not needed to make the sampler correct.

Second, a continuous \emph{piecewise-quadratic} potential admits a closed-form Hamiltonian flow on each piece. This piecewise construction gives exact HMC for truncated multivariate Gaussians (Pakman and Paninski 2014, cited here for orientation; the authors of this survey were unable to obtain the full text). It is directly relevant to sampling the cell-restricted Gaussians of Section~\ref{sec:shadow-filter}. A boundary-exact integrator is therefore an optimization for the kinked case and a requirement only for a genuine density jump.

\section{The UKF filtering collapse problem and marginal-likelihood gradient discontinuities}\label{sec:ukf-gradient}

The preceding section established that ordinary leapfrog HMC remains valid on continuous-kink targets, although its efficiency deteriorates. A more insidious geometry can arise when a Gaussian filter (Kalman, extended Kalman, or unscented) computes the marginal likelihood $p(y_{1:T}\mid\vartheta)$ in a model with occasionally binding constraints. Sampling then depends on the composed map from parameters through the filter to the log marginal likelihood,
$\vartheta\mapsto\log p(y_{1:T}\mid\vartheta)$.
This map can develop steep or discontinuous gradients even when the economic model uses a smooth softplus constraint rather than a hard max. The reason is that the filter's variance collapse near the bound amplifies gradient jumps in parameter space. This geometry occurs in the two-country shadow-rate model of Section~\ref{sec:shadow} when an unscented Kalman filter (UKF) replaces the particle filter. It is the primary obstacle to making gradient-based MCMC tractable for realistic DSGE models with zero lower bounds.

\subsection{Why UKF variance collapses at the bound}\label{sec:ukf-collapse-mechanism}

Suppose the shadow nominal rate $i_t^\star$ is a deterministic function of a latent state $x_t$. The constrained rate obeys
\begin{equation}\label{eq:ukf1}
 i_t = \ell + \operatorname{softplus}(i_t^\star - \ell, \alpha)
     = \ell + \frac{1}{\alpha}\log(1 + e^{\alpha(i_t^\star-\ell)}),
\end{equation}
and $i_t$ is observed with small Gaussian noise. When the data keep $i_t$ near the lower bound $\ell$ for many consecutive periods, a UKF, or any Gaussian filter, conditions the filtering distribution over $x_t$ on these near-bound observations. The predictive variance enters both the Kalman gain and the innovation covariance. When the observations are informative and consistent, the filtering updates shrink the variance.

The observations do not constrain every direction equally. For a state with a deeply negative shadow rate, all shadow-rate values below a threshold produce approximately the same observed rate $i_t\approx\ell$. The likelihood therefore provides no information about \emph{how far} below the bound the shadow rate lies. Together, variance shrinkage in the informative direction and the absence of information about shadow depth produce a one-sided collapse of the filtering distribution. Variance in the direction that affects the observed rate shrinks to near zero. In the orthogonal ``shadow-depth'' direction, variance can grow or remain large, depending on the transition dynamics.

The softplus in \eqref{eq:ukf1} does not prevent this collapse. It makes the transition from the informative region ($i_t^\star\approx\ell$) to the flat plateau ($i_t^\star\ll\ell$) continuous rather than abrupt. Increasing the softplus temperature $\alpha$ steepens this transition, so the geometry approaches that of a hard max.

\subsection{How variance collapse creates parameter-space gradient jumps}\label{sec:collapse-gradient}

A UKF propagates a deterministic set of sigma points through the nonlinear transition and observation maps. When the filtering variance is large, these points are spread out. Small parameter perturbations typically move all the points by small amounts, and the marginal log-likelihood gradient $\nabla_\vartheta\log p(y\mid\vartheta)$ varies smoothly. When the variance collapses, the points cluster tightly near the filtering mean. A small parameter change can then move the entire cluster across the steep part of the softplus at once. All sigma points change their regime contribution simultaneously, causing sharp changes in the predicted observations, innovation covariance, and log-likelihood contribution. These changes produce a large jump in $\nabla_\vartheta\log p(y\mid\vartheta)$.

The resulting target has a \emph{continuous-kink geometry}. The log likelihood $\log p(y\mid\vartheta)$ is continuous because both the softplus and the Gaussian observation density are continuous. Its gradient, however, can have jumps of size $O(\alpha)$ when a collapsed sigma-point cloud crosses the steep softplus region. The gradient jump comes from variance collapse in the filter combined with the deterministic sigma-point rule; it is not a modeling artifact.

\subsection{Why this differs from a value jump}\label{sec:kink-vs-jump-ukf}

UKF filtering collapse creates a continuous-kink target, rather than a value-jump target. The distinction concerns what changes at the boundary. At a parameter-space boundary where the sigma points cross the steep softplus region, the one-sided log-likelihood values match:
\[
 \lim_{\vartheta\to\vartheta_b^-}\log p(y\mid\vartheta)
 = \lim_{\vartheta\to\vartheta_b^+}\log p(y\mid\vartheta),
\]
whereas the one-sided gradients differ:
\[
 \lim_{\vartheta\to\vartheta_b^-}\nabla\log p(y\mid\vartheta)
 \ne \lim_{\vartheta\to\vartheta_b^+}\nabla\log p(y\mid\vartheta).
\]

The argument of Section~\ref{sec:kink-validity} establishes why ordinary Metropolized HMC remains \emph{valid} for this target. The leapfrog map is well defined almost everywhere, reversible, and volume preserving. These properties allow the Metropolis correction to target the correct posterior. Leapfrog is nevertheless \emph{inefficient}: each trajectory that crosses a steep-gradient region incurs an $O(\epsilon\,\Delta
g)$ energy error. Here $\Delta g$ denotes the gradient-jump magnitude and $\epsilon$ the step size. When $\Delta g$ is large, as with a high softplus temperature $\alpha$ or severe variance collapse, acceptance rates deteriorate. The chain can then become stuck near boundaries or within one filtering regime.

The practical obstacle is thus computational: the sampler is mathematically correct but explores the posterior poorly. Reflection/refraction HMC (Section~\ref{sec:event}) does not resolve this difficulty. Because the target is continuous, there is no potential jump $\Delta U$ to offset with kinetic energy. Improving efficiency instead requires either a different filter that maintains variance, such as particle methods, or an event-aware integrator that detects steep-gradient crossings and adjusts the trajectory accordingly.

\subsection{Why integer regime solvers create value jumps}\label{sec:integer-regime-jumps}

The preceding analysis assumed a smoothed constraint. The geometry changes fundamentally when a model instead uses an integer solver to select discrete regime paths, for example with a quantitative-easing on/off indicator or an OccBin-style spell-duration selector. The solver checks equilibrium consistency inequalities to map parameters to a discrete regime path $r_{1:T}\in\{0,1\}^T$. At a parameter-space boundary where it switches from path $j$ to path $k$, the regime-dependent state-space matrices $(F_{r_t},
Q_{r_t}, H_{r_t})$ change discontinuously even though the parameter $\vartheta$ changes smoothly.

The filtered likelihood therefore combines regime selection, construction of the state-space matrices, and filtering:
\[
 \vartheta \xrightarrow{\text{solver}} r_{1:T}
          \xrightarrow{\text{state-space}} (F_t, Q_t, H_t)
          \xrightarrow{\text{UKF}} \log p(y\mid\vartheta).
\]
When the regime path jumps, the UKF evaluates a different piecewise-linear system. The selected regime determines the innovation sequence, covariance updates, and cumulative log-likelihood. Consequently, approaching the boundary from opposite sides generally gives different log-likelihood values:
\[
 \lim_{\vartheta\to\vartheta_b^-}\log p(y\mid\vartheta, r^-)
 \ne \lim_{\vartheta\to\vartheta_b^+}\log p(y\mid\vartheta, r^+).
\]

This is a \emph{value-jump geometry}, for which reflection/refraction HMC and discontinuous HMC were designed. The potential has the finite jump
$\Delta U = -[\log p(y\mid\vartheta_b,r^+) - \log p(y\mid\vartheta_b,r^-)]$.
Standard leapfrog is invalid for this geometry: without an event map, a trajectory crossing the boundary systematically under- or over-estimates the Hamiltonian. This biases acceptance and produces the wrong stationary distribution.

\subsection{Measurement protocol for UKF-based posteriors}\label{sec:ukf-measurement}

Choosing an HMC extension for a UKF-filtered ZLB model requires measuring its actual boundary geometry first. The following protocol distinguishes continuous kinks from value jumps and quantifies gradient variation.

\paragraph{Step 1: Identify candidate boundaries.}
The protocol begins with a short HMC warm-up and records parameter values at which any of the following occurs:
\begin{itemize}
\item the softplus input $i_t^\star - \ell$ is near zero, within a few multiples of $1/\alpha$;
\item the UKF filtering variance for the shadow-rate state is small, orders of magnitude below the prior variance;
\item an integer regime solver, if present, changes its selected path.
\end{itemize}
These criteria locate the softplus transition, detect filter concentration, and identify discrete changes in the selected regime path, respectively.

\paragraph{Step 2: Value-continuity test.}
The value-continuity test compares likelihood values across a candidate boundary before examining gradients. For parameter pairs $(\vartheta^-, \vartheta^+)$ on opposite sides of the boundary, the absolute difference in log likelihood is:
\[
 \Delta\ell = |\log p(y\mid\vartheta^+) - \log p(y\mid\vartheta^-)|.
\]
The protocol classifies the boundary as a continuous kink if $\Delta\ell / |\log p(y\mid\vartheta^-)| < 10^{-6}$, and as a value jump if $\Delta\ell = O(1)$ or larger. The first criterion measures the difference relative to the magnitude of the log likelihood; the second identifies a difference of order unity or more.

\paragraph{Step 3: Gradient-jump magnitude.}
For a continuous-kink boundary, the protocol evaluates the one-sided gradients and measures their Euclidean difference:
\[
 \Delta g = \|\nabla_\vartheta\log p(y\mid\vartheta^+)
             - \nabla_\vartheta\log p(y\mid\vartheta^-)\|_2.
\]
If $\Delta g$ is large relative to the typical gradient norm, crossing the boundary changes the gradient substantially on the scale encountered by the sampler. This indicates poor leapfrog geometry.

\paragraph{Step 4: Boundary crossing rate.}
The effect on sampling also depends on how often trajectories encounter the boundary. The protocol runs HMC trajectories of length $L$ steps and records:
\begin{itemize}
\item the fraction of trajectories that cross a steep-gradient or regime-switch boundary;
\item the acceptance rate conditional on crossing, compared with the acceptance rate conditional on not crossing;
\item the average Hamiltonian error on crossing trajectories, compared with that on non-crossing trajectories.
\end{itemize}
A high crossing rate combined with low acceptance among crossing trajectories indicates a bottleneck: the sampler frequently proposes trajectories that cross the boundary, but few of those proposals are accepted.

\paragraph{Step 5: Gradient Lipschitz estimation.}
To estimate how rapidly the gradient varies with the parameters, the protocol samples pairs satisfying $\|\vartheta_1 - \vartheta_2\| = \delta$ and computes:
\[
 L_{\text{est}} = \max_{\text{pairs}}
   \frac{\|\nabla\log p(y\mid\vartheta_1) - \nabla\log p(y\mid\vartheta_2)\|}
        {\|\vartheta_1 - \vartheta_2\|}.
\]
Each ratio measures gradient change per unit parameter displacement, and the maximum selects the largest observed ratio among the sampled pairs. For a softplus model, $L_{\text{est}}$ grows approximately linearly with the temperature $\alpha$.

\subsection{Why particle filters smooth the gradient}\label{sec:pf-smoothing}

A particle filter approximates the filtering distribution by a weighted cloud of samples. Unlike a UKF's deterministic sigma points, particles retain diversity even when the likelihood strongly favors one region. This matters near the bound: when some particles have shadow rates slightly above it and others slightly below, the marginal likelihood is approximated by a mixture of predictive densities:
\[
 p(y_t\mid y_{1:t-1},\vartheta)
 \approx \sum_{i=1}^M W_{t-1}^i\, p(y_t\mid X_{t-1}^i,\vartheta),
\]
where each particle $X_{t-1}^i$ contributes its own predictive density, weighted by its particle weight. A small parameter change moves each particle slightly and changes the mixture weights smoothly, rather than flipping an entire cluster. The log-likelihood gradient $\nabla_\vartheta\log p(y\mid\vartheta)$ is then an average over particle contributions. This averaging smooths the sharp transitions that a collapsed UKF exhibits.

This smoothing mechanism is the fundamental reason differentiable particle filters, including LEDH-PFPF-OT and similar constructions, are pursued for ZLB models. They maintain diversity by design, which bounds gradient variation and makes HMC geometrically tractable. The computational cost is the price of that gradient smoothness.

\subsection{Candidate solutions for the UKF gradient problem}\label{sec:ukf-solutions}

The choice among the following strategies depends on the geometry diagnosed by the measurement protocol in Section~\ref{sec:ukf-measurement}. Continuous kinks and value jumps call for different remedies.

\paragraph{For continuous-kink targets (softplus ZLB, no integer regime).}
\begin{enumerate}
\item \textbf{Smaller step size.} Reducing $\epsilon$ lowers the energy error incurred at each boundary crossing, as described by \eqref{eq:56a}. This is the simplest fix, but it increases cost per effective sample.
\item \textbf{Event-aware integration (Tran and Kleppe 2025).} The integrator monitors the softplus input or a related margin to detect when a leapfrog step would cross a steep-gradient region. It truncates the step at the boundary, recomputes the gradient on the new side, and then continues. Bisection and re-evaluation make this computationally intensive, but the method reduces integration error.
\item \textbf{Neural surrogate proposal.} A neural network $s_\phi(\vartheta)$ is trained to approximate the log-likelihood gradient $\nabla_\vartheta\log p(y\mid\vartheta)$. Automatic differentiation through the UKF supplies the gradient values at the training points. The surrogate $s_\phi$ then drives leapfrog proposals, while Metropolis acceptance evaluates the \emph{true} UKF likelihood. With this correction, a poor surrogate leads to low acceptance while the sampler still targets the UKF posterior (Section~\ref{sec:pmcmc}).
\item \textbf{UKF variance inflation.} The predicted covariance $P_{t|t-1}$ is artificially inflated when the variance is below a threshold and the state is near a bound. This prevents full collapse and smooths the gradient. It also changes the filtering distribution and hence the posterior, so it constitutes a modified model rather than a computational trick.
\end{enumerate}

\paragraph{For value-jump targets (integer regime solver).}
\begin{enumerate}
\item \textbf{Reflection/refraction HMC (Section~\ref{sec:event}).} After detecting a regime-path change, the method computes the one-sided likelihoods $\log p(y\mid\vartheta_b, r^\pm)$ at the boundary and applies the energy-conserving event map \eqref{eq:59}--\eqref{eq:60}. The method is exact, but locating the implicit boundary is expensive: bisection requires full solver and filter re-evaluation at intermediate points.
\item \textbf{Discontinuous HMC (Section~\ref{sec:dhmc}).} This method uses Laplace momentum for parameters that affect regime switches. It avoids explicit event detection, but requires random coordinate order, randomized step size, and evaluation of the potential on both sides of proposed moves.
\item \textbf{Regime-conditional HMC.} Conditional on a regime path $r_{1:T}$, standard HMC targets $p(\vartheta\mid y,r)$. This conditional target is smooth, and the HMC step is fast. A discrete sampler then resamples $r$ from $p(r\mid\vartheta,y)$. The algorithm alternates these parameter and regime updates. Its mixing depends on the regime-parameter coupling: if $p(r\mid\vartheta,y)$ is peaked, regime updates have low acceptance.
\item \textbf{Particle MCMC.} This approach replaces the UKF with a particle filter, integrates out the regime path, and uses particle Metropolis-Hastings (Section~\ref{sec:pmcmc}). It avoids parameter-space discontinuities entirely, at the cost of running a particle filter at every proposal.
\end{enumerate}

\subsection{Why DHMC does not solve the continuous-kink problem}\label{sec:dhmc-limitation}

The UKF variance-collapse problem can make discontinuous HMC (Section~\ref{sec:dhmc}) seem appealing: ``the target is nonsmooth, and DHMC is for nonsmooth targets.'' That reasoning is incorrect because it conflates jumps in the potential with jumps in its gradient. DHMC with Laplace momentum and the reflection/refraction event map are designed for \emph{value jumps}: finite discontinuities in the potential $U(\vartheta)$ itself. For the UKF continuous-kink target, $U$ remains continuous everywhere; only its gradient jumps.

The DHMC reflection condition checks whether the available kinetic energy $|p_i|/m_i$ can pay the potential increase $\Delta U$. For a continuous target, $\Delta U = 0$ at every boundary, so there is no finite potential increase to pay. The reflection test is therefore trivially satisfied, and DHMC reduces to ordinary coordinate-wise leapfrog with Laplace momentum. Laplace momentum may improve mixing in some high-dimensional settings, where constant-velocity motion is cheaper to integrate than Gaussian trajectories. It does not, however, address the inefficiency caused by gradient jumps. That inefficiency arises from the $O(\epsilon\,\Delta g)$ leapfrog error term, which remains present regardless of the momentum distribution.

The distinction is sharp: DHMC and reflection/refraction HMC solve the value-jump problem (Case 3 of Section~\ref{sec:geometries}), but they do not solve the continuous-kink problem (Case 1). The latter requires smoother filters, event-aware gradient switching, or surrogate proposals.

\subsection{Open research questions}\label{sec:ukf-open}

The UKF gradient-discontinuity problem for ZLB models remains only partially solved. The following questions are active research directions.

\paragraph{Can event-aware integration be made practical?}
Tran and Kleppe's (2025) truncated-step integrator is mathematically valid for continuous-kink targets. Its practical cost comes from detecting an implicit OBC boundary: each bisection step requires full solver and filter re-evaluation. For a finite-horizon DSGE model, this can mean solving a 40-period piecewise-linear system dozens of times per HMC trajectory. Is this cost competitive with using smaller step sizes, or do the repeated evaluations dominate the cost per effective sample?

\paragraph{Is the marginal likelihood Lipschitz continuous?}
Even with softplus smoothing, the composed map $\vartheta\mapsto\log p(y\mid\vartheta)$ may fail to be globally Lipschitz continuous. Covariance updates and log determinants in the filter can amplify small changes in the solver output. Can a local Lipschitz constant be derived as a function of the softplus temperature, filter type, and time horizon? Or must OBC models be accepted as inherently non-Lipschitz, with integrators designed accordingly?

\paragraph{Can neural surrogates learn implicit boundaries?}
A neural score approximation $s_\phi(\vartheta)\approx\nabla_\vartheta\log
p(y\mid\vartheta)$, trained on branch-interior and near-boundary points, can drive HMC proposals with hard-target correction. The difficulty is that implicit OBC boundaries are high-codimension surfaces defined by solver consistency checks, rather than simple hyperplanes. Can a network generalize across these surfaces, or does it overfit to the training regime and fail when the chain explores a new boundary?

\paragraph{Does UKF variance inflation preserve scientific validity?}
Inflating $P_{t|t-1}$ to prevent collapse smooths the gradient, but also changes the filtering distribution and therefore the posterior. Can a principled calibration, such as tempering, robust Bayes, or conservative updating, make the inflated posterior remain a valid answer to the original inference question? Or is inflation purely a computational heuristic that must be validated post hoc by comparison with a particle-filter benchmark?

\paragraph{When is regime-conditional HMC competitive?}
Alternating between smooth HMC targeting $p(\vartheta\mid y,r)$ and a discrete sampler targeting $p(r\mid\vartheta,y)$ avoids joint regime-parameter boundaries. The discrete sampler can nevertheless mix poorly when the regime path is long and tightly constrained by the data. Under what coupling structure does this decomposition outperform joint methods?

\subsection{Interim recommendation}\label{sec:ukf-recommendation}

Until the research questions of Section~\ref{sec:ukf-open} are resolved, the safest path for production inference with UKF-filtered ZLB models is:

\begin{enumerate}
\item \textbf{Measure the geometry first.} The protocol of
      Section~\ref{sec:ukf-measurement} should be used to classify the target as
      continuous-kink or value-jump and to quantify gradient variation.
\item \textbf{For continuous-kink targets:} Standard HMC is recommended,
      with an aggressively tuned step size. Crossing diagnostics should report
      acceptance conditional on steep-gradient crossing, and mixing should be
      expected to be slower than for smooth targets. If the computational
      budget allows, a neural surrogate with hard-target correction should
      also be tested.
\item \textbf{For value-jump targets:} Standard leapfrog is invalid.
      Reflection/refraction HMC (Section~\ref{sec:event}) is exact but
      expensive. DHMC (Section~\ref{sec:dhmc}) is cheaper, but requires
      randomization of both coordinate order and step size. Regime-conditional
      HMC is a fallback if event methods prove intractable. Particle MCMC
      (Section~\ref{sec:pmcmc}) avoids the discontinuity entirely.
\item \textbf{Validate reversibility.} For any event-aware or DHMC method,
      validation should record a trajectory, flip its momentum, replay it,
      and check that the replay returns to the initial parameter position
      within tolerance:
      $\|\vartheta_{\text{final}} - \vartheta_{\text{initial}}\| < \epsilon$.
      Reversibility failure indicates a sampler bug or boundary-handling error.
\item \textbf{Report boundary telemetry.} Every UKF-based inference result
      should state its geometry classification (smooth/kink/jump), boundary
      crossing rate, acceptance rate conditional on crossing, and softplus
      temperature, if applicable. Without these diagnostics, acceptance rates
      and ESS alone cannot distinguish poor mixing caused by target geometry
      from poor mixing caused by sampler tuning.
\end{enumerate}

The particle-filter route (Section~\ref{sec:pmcmc}) remains the robust fallback
for production ZLB inference. By maintaining diversity and smoothing the
parameter-space likelihood, it sidesteps the UKF gradient problem. Current
differentiable implementations have biased gradients, and the particle-filter
route has high computational cost. These limitations motivate continued work
on UKF-based methods, but do not yet justify abandoning particle methods for
production ZLB inference.

\section{Reflection and refraction HMC for piecewise-smooth targets}\label{sec:event}

A value jump means that the target density itself is discontinuous across a
boundary, rather than only its gradient. At such a crossing, leapfrog's energy
error no longer vanishes as the step size shrinks, so the sampler needs an
explicit event mechanism. The physical analogy is optics: a particle hitting
an interface between media either refracts through it, with a kinetic-energy
loss equal to the potential jump, or reflects when it lacks the energy to
cross.

\subsection{The event calculation}\label{sec:event-calc}

Let the potential $U$ be smooth within regions separated by an affine
boundary with unit normal $\nu$. During a free-flight step, momentum
splits into a component perpendicular to the boundary and a tangential
remainder:
\begin{equation}\label{eq:57}
 p_\perp=(p^{\mathsf T}\nu)\nu,
 \qquad p_\parallel=p-p_\perp.
\end{equation}
At the boundary, the potential jump $\Delta U=U(q^+)-U(q^-)$ compares the potential on
the destination side with that on the approach side. Conservation of total
energy requires the kinetic-energy change to offset this potential change:
\begin{equation}\label{eq:58}
 K(p^+)-K(p^-)=-\Delta U.
\end{equation}
For unit Gaussian mass, if $\|p_\perp\|^2>2\Delta U$, the trajectory has enough normal
kinetic energy to cross. Refraction preserves the direction of the normal
momentum and adjusts its magnitude to conserve energy:
\begin{equation}\label{eq:59}
 p_\perp^+=
 \sqrt{\|p_\perp^-\|^2-2\Delta U}\,
 \frac{p_\perp^-}{\|p_\perp^-\|};
\end{equation}
otherwise the trajectory reflects, reversing its normal momentum:
\begin{equation}\label{eq:60}
 p_\perp^+=-p_\perp^-.
\end{equation}
The tangential momentum is unchanged in both cases. Afshar and Domke (2015,
pp.~2--4) first detect the earliest affine-boundary intersection along
the drift. They then apply \eqref{eq:59} or \eqref{eq:60}, continue the remaining
flight, and repeat the procedure for later boundaries.

Volume preservation must be attached to the correct map, a distinction that
is easy to miss. The reflection map \eqref{eq:60} is an orthogonal reflection
with unit absolute Jacobian. The refraction map \eqref{eq:59}, considered
\emph{in isolation}, has a local normal-momentum rescaling factor of
$s_-/s_+$, the ratio of incoming to outgoing signed normal momentum. This
factor describes the change in the spacing of nearby momenta, and the map
is not volume preserving when $\Delta U\neq0$.

Afshar and Domke prove volume preservation and reversibility for the
\emph{assembled} flight--event--flight trajectory step. The change in normal
velocity at the event also changes the position leg traversed during the
remaining flight time. The resulting stretch or compression of the set of
positions reached is precisely reciprocal to the local momentum rescaling,
so the Jacobian factors cancel. Combining this assembled step with leapfrog
away from boundaries gives detailed balance after the Metropolis correction
(Afshar and Domke 2015, pp.~5--7).

Equations \eqref{eq:57}--\eqref{eq:60} are stated for unit Gaussian mass and
cannot be used verbatim in a preconditioned sampler. Any practical HMC
implementation uses a positive-definite mass matrix $M$, with kinetic
energy $K(p)=\tfrac12 p^{\mathsf T}M^{-1}p$. The mass matrix determines the geometry in which the
momentum must be decomposed: the correct decomposition is orthogonal in the
$M^{-1}$ inner product, rather than the Euclidean inner product.
For a boundary with normal $n$, which need not have unit length, define
\begin{equation}\label{eq:60a}
 a=n^{\mathsf T}M^{-1}n,
 \qquad
 s_-=\frac{n^{\mathsf T}M^{-1}p_-}{\sqrt a}.
\end{equation}
The normalization accounts for the mass matrix and the length of the normal;
the signed coordinate measures the incoming normal momentum in this geometry.
For a potential jump $\Delta U$, the trajectory crosses when
$s_-^2>2\Delta U$, with the outgoing normal coordinate and full momentum given by
\begin{equation}\label{eq:60b}
 s_+=\operatorname{sign}(s_-)\sqrt{s_-^2-2\Delta U},
 \qquad
 p_+=p_-+n\,\frac{s_+-s_-}{\sqrt a},
\end{equation}
and reflects otherwise with
\begin{equation}\label{eq:60c}
 p_+=p_--2n\,\frac{n^{\mathsf T}M^{-1}p_-}{n^{\mathsf T}M^{-1}n}.
\end{equation}
Both maps change only the momentum component along $M^{-1}n$, conserve
$K(p)+U(q)$ across the event, and reduce to \eqref{eq:59}--\eqref{eq:60}
when $M=I$ and $\|n\|=1$.

The component description uses the normal transformed by the inverse mass
matrix, whereas the displayed crossing and reflection updates express the
momentum increment as a scalar multiple of the boundary normal itself.

Their Jacobians differ. The reflection \eqref{eq:60c} is an
$M^{-1}$-orthogonal reflection with unit absolute Jacobian, while the
refraction \eqref{eq:60b} alone has absolute Jacobian
$|s_-/s_+|\neq1$ for $\Delta U\neq0$. For refraction, volume preservation again
holds only for the composed flight--event--flight step: the local momentum
rescaling is cancelled by the reciprocal change in the position leg. A
correctness test must therefore target the composed trajectory map (or
account for the event Jacobian explicitly in the acceptance ratio), never
the isolated refraction.

An implementation must also fix the boundary orientation convention,
including the sign of $n$ and hence of $\Delta U$. It must specify
how to handle exact-equality events, $s_-^2=2\Delta U$, and grazing events,
$s_-=0$. Tests should exercise both dense and diagonal nonidentity
metrics, not only $M=I$.

\subsection{What transfers to a ZLB}\label{sec:event-zlb}

Reflection/refraction HMC is a good match when the posterior in the sampled
variables is piecewise smooth and the switching surfaces are explicit affine
boundaries, or at least surfaces whose first intersection with a trajectory
can be found reliably. A nonlinear endogenous regime solver makes these
requirements harder to meet. In that case, the boundary is an implicit surface
\begin{equation}\label{eq:61}
 b(q)=0,
\end{equation}
and a trajectory needs (i) a first-root solver for $b(q(t))$, (ii) the correct
one-sided values of the full log posterior, and (iii) an argument that the
selected equilibrium branch is single-valued at the event. The first
requirement locates the earliest boundary event; the others establish the
posterior values on either side of that event.

The Metropolis correction absorbs integration error relative to a well-defined
target. When a root search fails or the regime solver lands on a different
fixed point, the two one-sided posterior values in \eqref{eq:5b} are no longer
defined. Such a failure requires repair of the problem statement itself,
rather than of the sampler.

\section{Discontinuous HMC with Laplace momentum}\label{sec:dhmc}

Reflection--refraction HMC incurs a separate cost at every boundary crossing.
When a discrete or ordinal variable is embedded into the continuous state,
a trajectory may cross hundreds of density jumps per proposal. In this
setting, changing the momentum distribution makes each crossing dramatically
cheaper.

\subsection{Embedding an ordinal regime}\label{sec:dhmc-embed}

Nishimura, Dunson, and Lu (2020, pp.~366--368) embed an integer parameter
$N$ in a real variable $\widetilde N$ by assigning each integer value
an interval:
\begin{equation}\label{eq:62}
 N=n \iff \widetilde N\in(a_n,a_{n+1}],\qquad
 \widetilde\pi(\widetilde n)=
 \sum_n\frac{\pi_N(n)}{a_{n+1}-a_n}
 \mathbf 1\{a_n<\widetilde n\leq a_{n+1}\}.
\end{equation}
Each interval excludes its lower endpoint and includes its upper endpoint.
Within the interval, the embedded density is constant, with height equal to
the original probability divided by the interval length. The interval-length
denominator is essential: it is a Jacobian-like factor that preserves
probability mass. Integrating the constant density across an interval
multiplies it by that same length, leaving total embedded mass
$\pi_N(n)$. Thus integrating $\widetilde\pi$ over $(a_n,a_{n+1}]$ recovers the
original probability.

The embedding also imposes an ordering on the discrete values. An arbitrary
ordering can place high-probability values in separated intervals, so
energy-conserving dynamics may mix poorly between the resulting modes.

\subsection{Why Gaussian momentum is expensive, and the Laplace alternative}\label{sec:dhmc-laplace}

When a Gaussian-momentum trajectory crosses many density jumps, every
crossing must be located and corrected as in Section~\ref{sec:event-calc}.
For a discrete coordinate with posterior uncertainty of hundreds of units,
this means hundreds of target evaluations per trajectory. Ignoring the jumps
and relying on the Metropolis step does not remove the need for this event
handling: the leapfrog error at a jump does not vanish as
$\epsilon\to0$, because the unaccounted potential change remains finite
regardless of the step size (Nishimura, Dunson, and Lu 2020,
pp.~367--370).

The alternative is to choose independent Laplace momentum,
\begin{equation}\label{eq:63}
 K(p)=\sum_i |p_i|/m_i,
\qquad
 \dot q_i=m_i^{-1}\operatorname{sign}(p_i).
\end{equation}
With independent Laplace momentum, kinetic energy sums the absolute coordinate momenta divided by their respective masses. Each coordinate's velocity depends on the momentum's sign, not its magnitude. Losing kinetic energy therefore does not slow the coordinate, so its proposed position after a step is known in advance regardless of how much kinetic energy it loses along the way. For coordinate $i$, the sampler proposes the endpoint and computes the potential increase relative to the current position:
\begin{equation}\label{eq:64}
 q_i^\star=q_i+\epsilon m_i^{-1}\operatorname{sign}(p_i),
\qquad
 \Delta U=U(q^\star)-U(q).
\end{equation}
If $|p_i|/m_i>\Delta U$, the coordinate has enough kinetic energy for the move. The sampler accepts the endpoint and subtracts the potential increase from the coordinate's kinetic energy by updating its momentum,
\begin{equation}\label{eq:65}
 p_i\leftarrow p_i-\operatorname{sign}(p_i)m_i\Delta U;
\end{equation}
otherwise it keeps $q_i$ unchanged and reflects the momentum, $p_i\leftarrow-p_i$. For an accepted move, a rise in potential reduces the momentum magnitude, while a fall increases it.

The sampler applies the coordinate maps sequentially in a randomly permuted order. This random permutation is needed for reversibility in distribution. Nishimura et al.\ prove that the coordinate map is volume preserving and reversible almost everywhere, and that the assembled integrator is valid with a final Metropolis correction (Nishimura, Dunson, and Lu 2020, pp.~371--374). They also randomize the step size to avoid a fixed-grid reducibility problem, in which updates remain confined to a grid of positions.

This algorithm is attractive when the regime is an explicit coordinate: a move in \eqref{eq:64} can cross several ordinal intervals using a single endpoint likelihood evaluation. Gaussian momentum lacks this economy. The algorithm is a poor fit, however, for a solver that hides a complicated regime path inside one black-box function. The endpoint energy difference $\Delta U$ is meaningful only when both endpoints evaluate the same well-defined posterior, namely the density of the model actually being estimated. All ambiguity about which branch the solver selects must therefore be resolved.

\section{Mixed HMC for an explicit regime path}\label{sec:mixed}

Embedding a discrete variable into the real line, as in Section~\ref{sec:dhmc-embed}, imposes an ordering that may not exist naturally. Zhou's mixed HMC (2020, pp.~2--4) instead keeps the variable discrete. Let $z$ denote the discrete state and $q$ the continuous coordinates, with joint target $\pi(z,q)\propto e^{-U(z,q)}$. The sampler attaches to each discrete coordinate $j$ a clock $s_j$ living on a flat torus and a discrete momentum $p_j^D$. The clocks determine when to attempt discrete moves, while the continuous coordinates evolve by ordinary Hamiltonian dynamics. When a clock reaches its boundary, the sampler proposes $\widetilde z$ from an irreducible single-site proposal $Q_j(\widetilde z\mid z)$ and treats the move as an event.

At this event, the continuous state is held fixed. The energy cost of the proposed discrete move accounts for both the change in target density and the forward and reverse proposal probabilities. Define
\begin{equation}\label{eq:66}
 \Delta E=\log\frac{\pi(z,q)Q_j(\widetilde z\mid z)}
                         {\pi(\widetilde z,q)Q_j(z\mid\widetilde z)}.
\end{equation}
For discrete kinetic energy $k_j^D(p_j^D)$, the sampler accepts the move (the trajectory ``crosses'') and changes the momentum to satisfy
\begin{equation}\label{eq:67}
 k_j^D(p_j^{D,+})=k_j^D(p_j^{D,-})-\Delta E
\end{equation}
provided that the right-hand side is nonnegative. This condition requires the kinetic energy remaining after paying the proposal-corrected energy cost to be nonnegative. Otherwise, the momentum reflects and $z$ is kept unchanged. This is the refraction--reflection rule of Section~\ref{sec:event-calc} applied to a discrete jump: the proposal-corrected energy difference in \eqref{eq:66} plays the role of $\Delta U$.

The sampler advances the continuous coordinates with a reversible, volume-preserving integrator and applies a final Metropolis correction to make the joint distribution invariant. The invariance theorem depends on the proposal ratio in \eqref{eq:66}; it does not require a conditional Gibbs draw.

For a \emph{stochastic} ZLB path, the discrete state can be $z=r_{1:T}$, the full regime path, or the sampler can use a blocked path proposal. The energy difference in equation \eqref{eq:66} then includes the complete regime-path probability and the nonlinear state-space likelihood. This construction is conceptually clean, but evaluating it can be expensive: changing even one regime indicator may require re-solving a long state path to evaluate the new likelihood. For the deterministic indicator of \eqref{eq:5a}, the regime path is instead a function of the continuous state. With that state held fixed, a proposal that changes only the regime has zero support, so this mixed-HMC construction does not apply at all. In either case, defining the energy difference in \eqref{eq:66} first requires an explicit probability law over regime paths, including a stance on how multiple equilibrium paths are treated.

\section{Smooth categorical relaxations and nonsmooth energies}\label{sec:relaxations}

If exact discrete dynamics are too expensive, an alternative is to smooth away the discreteness and run ordinary HMC on the relaxed target. This choice changes the target: the relaxed posterior is not the original posterior. The Concrete / Gumbel-softmax method (Torgander, Magnusson, and Wallin 2024, pp.~2--3) maps categorical probabilities to a continuous simplex variable. As the temperature $\tau\to0$, the variable becomes nearly one-hot, placing almost all its weight on a single category. At the same time, the transformed posterior develops steep geometry: gradients with respect to simplex coordinates can diverge as those coordinates approach zero, and the temperature direction becomes ill-conditioned. The method is therefore a useful relaxation baseline, but it is neither an exact categorical sampler nor an exact ZLB sampler. Comparisons against it should report both the temperature used and the posterior discrepancy induced by the relaxation.

A different smoothing tradition addresses energies that are nonsmooth but continuous. Chaari et al.'s nonsmooth HMC (2016, pp.~1--5) uses subgradients or a proximity operator
\begin{equation}\label{eq:68}
 \operatorname{prox}_{\phi}(x)=
 \arg\min_y\left\{\phi(y)+\tfrac12\|y-x\|^2\right\}
\end{equation}
for convex nondifferentiable energies such as an $\ell_1$ penalty. The minimization balances the energy against a penalty for squared distance from the input point.

The distinction in Section~\ref{sec:kink-validity} applies fully here. The proximity operator \eqref{eq:68} is built for convex \emph{continuous} energies: their nondifferentiability is a kink, where the energy remains continuous even though a derivative does not exist. A density discontinuity, by contrast, is a finite jump in $U$ that no subgradient or proximal smoothing removes. Equilibrium regime selection is a modeling choice, not an energy. Prox-HMC therefore enters the ZLB problem exactly when a convex nonsmooth approximation of the constrained model is deliberately chosen.

\section{Particle MCMC, pseudo-marginal HMC, and differentiable state-space HMC}\label{sec:pmcmc}

Every gradient method above faces the same obstacle: at the zero lower bound, the exact likelihood is an integral over latent paths whose integrand is only piecewise smooth. Particle methods take a different approach. They retain the model rather than smoothing it and replace the intractable likelihood with an unbiased Monte Carlo estimate. They then build a Markov chain that is \emph{exactly} invariant for the posterior despite the estimation noise. The following discussion develops this construction, asks when it can be combined with Hamiltonian dynamics, and examines what modern ``differentiable'' particle filters do and do not deliver.

\subsection{Particle MCMC's exact extended target}\label{sec:pm-target}

The whole construction rests on a single, disarmingly simple identity. Let $u$ collect \emph{all} random variables consumed by a run of the particle filter, including every proposal draw and every resampling variable, with joint density $m_\vartheta(u)$, and let
$\widehat L(\vartheta,u)\geq 0$ be the resulting nonnegative likelihood estimate. For each fixed parameter value, the standard particle filter is unbiased: averaging its estimate over the filter randomness gives the exact likelihood,
\begin{equation}\label{eq:69}
 \int \widehat L(\vartheta,u)\,m_\vartheta(u)\,du = L(\vartheta).
\end{equation}
The \emph{extended} target treats the parameter and the filter randomness as a joint state, the pair $(\vartheta,u)$:
\begin{equation}\label{eq:70}
 \bar\pi(\vartheta,u) \propto p(\vartheta)\,\widehat L(\vartheta,u)\,m_\vartheta(u).
\end{equation}
When \eqref{eq:70} is integrated over $u$, the parameter is held fixed, so the prior factors out of the integral. The unbiasedness identity \eqref{eq:69} then replaces the integral of the likelihood estimate against the filter's random-number density by the exact likelihood. Thus the $\vartheta$-marginal has density proportional to $p(\vartheta)L(\vartheta)$. After normalization, this is exactly the desired posterior, with no bias from the finite particle count. This is the pseudo-marginal identity underlying particle MCMC (Andrieu, Doucet, and
Holenstein 2010). It gives a practical sampling method: a Metropolis--Hastings chain targets \eqref{eq:70}, pairs each proposed parameter $\vartheta'$ with a fresh filter run $u'$, and uses the \emph{estimated} likelihoods in its acceptance ratio. This is a valid MCMC algorithm for the exact posterior. Likelihood-estimation noise belongs to the enlarged state space on which the chain is exact; it leaves no approximation error in the parameter marginal.

\subsection{What pseudo-marginal HMC additionally requires}\label{sec:pmhmc}

The target in \eqref{eq:70} suggests a further step: if $\log\widehat L(\vartheta,u)$ can be evaluated, it is tempting to differentiate it and run HMC on the extended target. Alenl\"ov, Doucet, and Lindsten (2021, Secs.~2--3) construct precisely such an extended Hamiltonian and specify what this step requires:
\begin{enumerate}
\item a nonnegative unbiased estimator $\widehat L(\vartheta,U)$;
\item standard Gaussian auxiliary variables $U\sim N(0,I)$;
\item continuous differentiability, in both the parameter and auxiliary coordinates, of the map
      $(\vartheta,U)\mapsto\widehat L(\vartheta,U)$; and
\item pointwise evaluability of the log-estimator gradient $\nabla\log\widehat L(\vartheta,U)$.
\end{enumerate}
Under these conditions, the algorithm is exact for the extended target. As the number of importance samples grows, its trajectories approach those of ideal HMC on the marginal. Exactness concerns the target distribution; the limiting statement concerns the trajectories.

The obstacle is condition (iii). As Alenl\"ov, Doucet, and Lindsten explicitly note, ordinary particle-filter likelihoods typically \emph{fail} this requirement. In multinomial resampling, the selected ancestor index jumps when a uniform random number crosses a cumulative-weight threshold. That discrete change makes $\widehat L$ discontinuous in $(\vartheta,U)$, even when the underlying model is smooth.

This creates a sharp distinction for ZLB inference. An unbiased bootstrap particle filter supports particle Metropolis--Hastings through the identity in Section~\ref{sec:pm-target}. Pseudo-marginal HMC also requires continuous differentiability, condition (iii), which ordinary resampling breaks. A differentiable resampling scheme or a transport-based relaxation can restore condition (iii), but doing so changes what the computed gradient represents. The gradient of a smoothed or transport-resampled computation is either an approximation or the exact gradient of a \emph{different} extended target. The latter interpretation is justified only once that target's full law and Metropolis correction have been written down. Pseudo-marginal HMC therefore requires a separately constructed differentiable extended estimator. Constructing it is a substantive design task beyond implementing a working particle Metropolis--Hastings sampler.

The correlated pseudo-marginal method (Deligiannidis, Doucet, and Pitt
2018) addresses a different issue: it improves particle Metropolis--Hastings mixing by correlating successive $U$ draws. This makes it the appropriate variance-reduction comparator. The differentiability repair needed for pseudo-marginal HMC must come from the resampling scheme itself.

\subsection{Joint latent-state HMC}\label{sec:joint-hmc}

An alternative to marginalizing the states by particle filtering is to sample the parameters, initial state, and shocks, $(\vartheta,x_0,\epsilon_{1:T})$, directly using the joint density \eqref{eq:56}. This avoids a particle likelihood altogether and works well for smooth nonlinear models (Childers et al.\ 2022, pp.~10--16). At the ZLB, however, the structure of the joint density depends on how the regime is treated.

With a stochastic explicit regime path, \eqref{eq:56} is a mixed target with continuous and discrete coordinates. The state transition and observation density are piecewise in $r_t$, and the discrete path requires a transition kernel that assigns positive probability to regime changes. With a deterministic regime, the density instead concentrates on the support described in \eqref{eq:5a}, and the discontinuity lies in the continuous coordinates. Joint HMC therefore supplies the continuous-coordinate component of a complete sampler for this problem class. The regime treatment determines what must be composed with that component.

\subsection{Particle Gibbs and ancestor sampling}\label{sec:pgas}

Particle Gibbs samples an entire latent trajectory as a Gibbs block, rather than estimating the marginal likelihood. Given a reference path $x'_{1:T}$, conditional sequential Monte Carlo runs a particle filter with one particle pinned to that path. Plain particle Gibbs has a well-known mixing problem: when the particle genealogy collapses, the paths share early ancestors. A newly sampled trajectory then tends to retain that early history, so mixing over the early part of the trajectory stalls. Particle Gibbs with ancestor sampling (PGAS) repairs this by resampling the fixed particle's ancestor at each time step, allowing that particle to inherit a different past (Lindsten, Jordan,
and Sch\"on 2014, Secs.~3--5).

For a Markov state-space model, one PGAS sweep proceeds as follows. At $t=1$, the sweep draws particles $1{:}M-1$ from the initial proposal, sets $X_1^M=x'_1$ to fix the remaining particle at the reference value, and weights all $M$ particles. For each $t=2{:}T$:
\begin{enumerate}
\item For $j<M$, the sweep draws $A_t^j\sim\operatorname{Categorical}(W_{t-1})$ and propagates $X_t^j$ from the proposal in \eqref{eq:12b}.
\item It sets $X_t^M=x'_t$, keeping this particle at the reference state, but draws its ancestor with probabilities
\begin{equation}\label{eq:70a}
 \Pr(A_t^M=j)\propto W_{t-1}^j\,
 p(x'_t\mid X_{t-1}^j,\vartheta).
\end{equation}
\item It recomputes the importance weights. At $T$, it draws a terminal particle according to the weights $W_T$ and traces that particle's ancestors to obtain the new path.
\end{enumerate}
Equation \eqref{eq:70a} is the Markov simplification of the general ancestor weight. The general weight multiplies the old particle weight by the target density after attaching the fixed future to a candidate past, divided by the target density before that attachment. In the displayed Markov probabilities, the old importance weight is multiplied by the transition density to the fixed reference state, which measures how well the candidate ancestor connects to that state. Lindsten, Jordan, and Sch\"on prove that this kernel leaves the joint smoothing distribution invariant for \emph{every} particle count. Their proof treats the sweep as a partially collapsed Gibbs update on an extended particle system (Lindsten, Jordan, and Sch\"on 2014, Algorithm~2 and Theorem~1).

At the bound, the choice of particle state depends on whether the regime is deterministic or stochastic. For a deterministic ZLB, the particle state is the continuous state or shock path, and the regime $r_t=R(q)$ is derived from it. PGAS then mixes over paths, with the regime sequence following from those paths. For a stochastic regime, the particle state is $(r_t,x_t)$, so PGAS mixes over regimes and states jointly. PGAS does \emph{not} permit an independent regime update when the regime is a function of the continuous path: there is no separate coordinate to resample, and a label inconsistent with that path has zero support.

\subsection{What differentiable particle filters do and do not supply}\label{sec:dpf}

Two recent lines of work connect particle filters with automatic differentiation. Each supplies a precise result, but they change different parts of the computation.

\'Scibior and Wood keep the ordinary particle filter's forward values entirely unchanged. They insert stop-gradient weight ratios that alter the differentiation of the computation, so automatic differentiation returns the Fisher-identity score estimator associated with the particle genealogy (\'Scibior and Wood 2021,
Algorithm~1 and Sec.~3). This score estimates the gradient of the marginal log likelihood. Poyiadjis, Doucet, and Singh derive the underlying path-space and forward-smoothing score estimators and show that they differ in computational cost and variance growth (Poyiadjis, Doucet, and Singh 2011, Secs.~2--3). The particle-genealogy score estimator is consistent as the particle count grows. At any finite particle count, however, it carries Monte Carlo noise. Feeding that noisy score to a leapfrog integrator as though it were the deterministic force of the marginal log likelihood produces a stochastic-gradient chain. The usual HMC exactness proof no longer covers that chain's invariant distribution.

Corenflos et al.\ instead replace discrete resampling with a differentiable entropy-regularized optimal-transport map. Their construction changes the finite-particle forward algorithm itself and introduces a regularization parameter. Their consistency analysis requires the particle count to grow \emph{and} the regularization to shrink (Corenflos et al.\ 2021, Sec.~4). Thus the modification affects the values computed by the filter as well as their differentiation.

Both constructions are important comparison methods. Exact use of either inside HMC requires the full pseudo-marginal discipline of Section~\ref{sec:pmhmc}: an explicitly written extended target, random numbers held fixed during each reversible trajectory, and a Metropolis energy computed for that same extended density. Inserting either gradient as a noisy force into ordinary HMC skips all three requirements and therefore loses the exactness proof.

\section{What the literature supports}\label{sec:synthesis}

The literature surveyed so far supports a layered architecture for designing Bayesian inference at an occasionally binding bound. Its four layers are connected, but each has a different job. Much practical difficulty comes from asking one layer to perform another's work.

\textbf{Layer A: a verified linear benchmark.} This layer calls for implementing OccBin-style piecewise path iteration and the conditional Kalman likelihood \eqref{eq:8}--\eqref{eq:11}. Both implementations should be validated against a published example and against synthetic data with known regime paths. Although this layer may appear routine, it establishes timing conventions, inequality signs, likelihood normalization, the treatment of missing observations, and multiple-fixed-point diagnostics. Errors here propagate undetected into every later layer.

\textbf{Layer B: a nonlinear likelihood authority.} Layer B starts with an implementation of the bootstrap
particle filter \eqref{eq:12a}--\eqref{eq:13} and, where the decision rules are
piecewise linear and continuous, the Aruoba et al.\ conditionally optimal
special case. The bootstrap filter handles any nonlinear transition that can
be simulated. For piecewise-linear continuous decision rules, the
conditionally optimal filter provides the low-variance filtering alternative.
Particle Metropolis--Hastings then supplies, via
\eqref{eq:70}, an extended target with a finite number of particles.
Its parameter marginal is the exact posterior. This exactness makes particle
Metropolis--Hastings the ``ground truth'' sampler against which faster methods
are judged.

\textbf{Layer C: deterministic complementarity.} The regime is the derived
quantity $R(q)$, so it is not an independent unknown. Joint-state
random-walk or slice updates and particle MCMC are the safe starting methods.
Event-aware HMC should be added only after the complete nonlinear solution map
provides a first-boundary oracle, which locates the earliest crossing, and
both one-sided posterior values at that crossing. A fixed-branch HMC update
explores efficiently within one region. Entering the other regime requires
moving the continuous state across the boundary: by
\eqref{eq:5a}, a proposal that changes only the regime label has zero target
density.

\textbf{Layer D: stochastic or explicitly selected regimes.} When both regime
values carry positive probability, the initial kernel should alternate HMC at
a fixed regime path with PGAS or another valid discrete Metropolis kernel.
Zhou's mixed HMC becomes an acceleration candidate only after this simpler
composition has been checked against exact enumeration on a toy model. Its
clock changes when updates occur while leaving the target distribution
unchanged. Agreement on a small exactly solvable case is therefore the
appropriate test before considering mixed HMC as an acceleration method.

Throughout, the Concrete-relaxation and softplus routes of
Section~\ref{sec:relaxations} should be understood as approximation arms. In
particular, the shadow-rate case study of Section~\ref{sec:shadow} uses a smooth
softplus bound and is therefore a useful comparator for the hard-bound methods,
not evidence about genuinely discontinuous semantics.

The two case studies later in this survey instantiate these layers in
complementary ways. The two-country shadow-rate model of Section~\ref{sec:shadow} is a
hard-bound observation-map kink target: Layer B specializes to the bootstrap
particle-filter likelihood of Section~\ref{sec:shadow-filter}, with the ideal cell
decomposition \eqref{eq:84}--\eqref{eq:87} as its proposal and diagnostic
layer; Layer C specializes to joint kinked-target HMC under the validity
conditions of Section~\ref{sec:kink-validity}; and Layers A and D do not apply to that
model as specified. The New Keynesian model of Section~\ref{sec:nk} is the converse:
Layers A, C, and D carry the weight, the genuine discontinuity lives in the
solution correspondence and its selection rule (Section~\ref{sec:nk-mult}), and Layer D
becomes available exactly when selection is made stochastic
(Section~\ref{sec:nk-sunspot}).

\section{Exact targets and implementable transition kernels}\label{sec:targets}

A sampling method cannot be fully specified without its target distribution.
This section therefore states the exact posterior density and a transition
kernel that provably preserves it for each of the three regime treatments:
deterministic threshold, stochastic regime, and multiple equilibria. This
explicit specification is useful later because many subtle errors in
bound-aware samplers concern \emph{which} distribution is being sampled,
rather than the dynamics used to sample it.

\subsection{Deterministic threshold target}\label{sec:target-det}

Let $q=(\vartheta,x_0,\epsilon_{1:T})$ collect the parameters, initial state, and innovation
sequence. The regime and state are defined recursively by
$r_t=\mathcal R(x_{t-1},\epsilon_t,\vartheta)$ and
$x_t=F_{r_t}(x_{t-1},\epsilon_t,\vartheta)$, respectively. Assume that this recursion has a unique
solution. The regime is then a deterministic function of
$q$, so it adds no independent unknown. The posterior is therefore
a density on the continuous variables alone:
\begin{equation}\label{eq:71}
 \Pi_D(q)\propto p(\vartheta)\,p(x_0\mid\vartheta)
 \prod_{t=1}^T p(\epsilon_t\mid\vartheta)\,
 g_{r_t(q)}\bigl(y_t\mid x_t(q),\vartheta\bigr).
\end{equation}
The density combines the parameter and initial-state priors with the
innovation densities and observation likelihoods evaluated along the uniquely
determined state and regime paths.

An event-aware HMC step targeting \eqref{eq:71} proceeds within a region
with fixed $r_{1:T}$ by differentiating the same finite program that
evaluates \eqref{eq:71}. The gradient must correspond to the density
actually being computed, rather than to a surrogate. At the first event
$b_k(q(t))=0$ along the trajectory, the step evaluates the two one-sided
values of the \emph{complete} posterior. These values determine the potential
jump through \eqref{eq:5b}. The step then applies the general-mass event
map \eqref{eq:60b} or \eqref{eq:60c}; the simpler maps
\eqref{eq:59}/\eqref{eq:60} apply only in the unit-mass special case.
After the event update, the flight continues for the unused part of the
position step. At the end of the trajectory, the momentum is reversed and
the proposal is accepted with probability
\begin{equation}\label{eq:72}
 \alpha_D=1\wedge\exp\{H(q,p)-H(q',p')\}.
\end{equation}

Implementing this full step requires the event locator to return the
\emph{earliest} crossing and the branch solver to be deterministic and
single-valued on each side of the boundary. The \emph{composed}
flight--event--flight trajectory map must be reversible. It must also preserve
volume, or the acceptance ratio must explicitly account for the event
Jacobian. The isolated refraction map rescales the normal momentum
(Section~\ref{sec:event-calc}), so volume preservation holds only for the
composed map. A correctness test must therefore exercise the full
flight--event--flight composition. If any of these conditions is missing,
particle Metropolis--Hastings using \eqref{eq:13}, or a non-gradient
joint-state kernel, remains the exact baseline.

\subsection{Stochastic mixed target}\label{sec:target-stoch}

If regimes carry genuine transition mass, the regime path is an unknown in its
own right. The target therefore includes the regime path alongside the
parameters and continuous state path:
\begin{equation}\label{eq:73}
 \Pi_S(\vartheta,x_{0:T},r_{1:T})\propto
 p(\vartheta)\,p(x_0\mid\vartheta)
 \prod_{t=1}^T p(r_t\mid r_{t-1},x_{t-1},\vartheta)\,
 p(x_t\mid x_{t-1},r_t,\vartheta)\,
 g_{r_t}(y_t\mid x_t,\vartheta).
\end{equation}
At each date, the factors account for the regime transition, the state
transition conditional on that regime, and the observation likelihood.

A transparent first kernel composes two invariant updates in sequence:
\begin{enumerate}
\item With $r_{1:T}$ held fixed, ordinary Metropolized HMC updates
      $(\vartheta,x_{0:T})$ using the gradient of
      $\log\Pi_S$. With the regime path frozen, the target is smooth
      in the continuous coordinates, so standard HMC theory applies.
\item With the continuous variables held fixed, the second component updates
      one regime, a regime block, or the whole latent path with Metropolis,
      forward-filter backward-sampling, or PGAS. For a proposal
      $Q(r'\mid r,q)$, a regime-only move is accepted with probability
\begin{equation}\label{eq:74}
 \alpha_R=1\wedge
 \frac{\Pi_S(q,r')\,Q(r\mid r',q)}
      {\Pi_S(q,r)\,Q(r'\mid r,q)}.
\end{equation}
      The ratio compares the joint posterior at the proposed and current
      regimes and includes the reverse and forward proposal probabilities.
\end{enumerate}
Each component preserves $\Pi_S$, so applying them in sequence also
preserves it. This is the standard Metropolis-within-Gibbs argument; it
requires no smoothness across regime changes. Mixed HMC replaces the second,
regime-update component with the clock dynamics
\eqref{eq:66}--\eqref{eq:67} and interleaves those dynamics with the first,
continuous-variable HMC component.

A regime proposal that also transforms the continuous coordinates requires
additional care. Such a transformation might keep an economic constraint
satisfied after the regime flips. The move is then no longer discrete-only:
its reverse proposal density and the absolute Jacobian of the continuous
transformation must be incorporated into \eqref{eq:74}.

\subsection{Multiple equilibria}\label{sec:target-multi}

Finally, suppose the model admits multiple solutions
$s_t\in\mathcal S(x_{t-1},\epsilon_t,\vartheta)$ at some states, as models with occasionally binding
constraints can. The model must then specify a selection mass
$p(s_t\mid x_{t-1},\epsilon_t,\vartheta)$, which determines how it selects among the available
solutions. The selected solution $s_t$ must either be included as
a latent variable in \eqref{eq:73} or integrated out. Without that mass,
neither \eqref{eq:71} nor \eqref{eq:73} defines a target, because the
structural parameters do not determine the likelihood of the data. Supplying
the selection mass is therefore part of the economic model specification.
The sampler's job begins only once the target distribution from which it must
draw exists.

The two case studies that follow instantiate these abstract targets.
Section~\ref{sec:shadow} develops a two-country shadow-rate term-structure
model whose bound enters through the observation map. Section~\ref{sec:nk}
develops a small New Keynesian model whose bound enters through the solution
correspondence itself.

\section{Case study I: a two-country shadow-rate term-structure model}
\label{sec:shadow}

Near a lower bound on nominal rates, observed yield curves flatten and
compress against the bound while the economically meaningful latent state
keeps moving. This latent state is the ``shadow'' curve: the curve that would
prevail in the absence of the bound. A yield pinned at the bound for two years
is consistent with a shadow rate slightly below it or far below it.
Monetary-policy analysis at the bound requires distinguishing these two
situations. This need motivates Bayesian inference on the latent shadow
factors and model parameters from bounded yield observations. It is a
state-space problem in which the bound enters the observation map as a
censoring nonlinearity.

This case study builds a concrete two-country shadow-rate term-structure model
from scratch, places it in the literature, and then applies the methods of the
preceding sections. It identifies the model's posterior geometry and examines
which filters and samplers are exact for which versions of the model. It also
examines where the smooth approximation that makes gradient-based sampling
easy ceases to be innocuous for inference.

\subsection{The model}
\label{sec:shadow-model}

The eight-dimensional latent state is
\(x_t=(x_t^{d\mathsf T},x_t^{f\mathsf T},x_t^{b\mathsf T})^{\mathsf T}
\in\mathbb R^8\). It contains domestic (\(d\)) and foreign (\(f\)) dynamic Nelson--Siegel (DNS) factor triples, each comprising level, slope, and curvature, together with two directed FX-basis factors. The transition is a stationary Gaussian VAR(1) toward a long-run mean,
\begin{equation}\label{eq:75}
 x_t=(I-\Phi)\bar\theta+\Phi x_{t-1}+\eta_t,
 \qquad \eta_t\sim N(0,Q),
\end{equation}
where the lagged-state term carries persistence and the Gaussian innovation adds new state uncertainty. In the working example below, the estimated parameters \(\vartheta\) consist of components of the long-run mean \(\bar\theta\) together with grouped log measurement-noise scales. The transition matrix \(\Phi\), the innovation covariance \(Q\), and the initial covariance are held fixed. This deliberately small nine-parameter chart keeps attention on the bound rather than on the VAR dynamics.

For country \(c\in\{d,f\}\), the shadow instantaneous forward curve at horizon \(s\) is affine in that country's factors, with a fixed decay parameter \(\lambda_c\):
\begin{equation}\label{eq:76}
 f_c(s;x)=a_c(s)^{\mathsf T}x^c,
 \qquad
 a_c(s)=\bigl(1,\;e^{-\lambda_c s},\;\lambda_c s\,e^{-\lambda_c s}
 \bigr)^{\mathsf T}.
\end{equation}
The loadings multiply the level, slope, and curvature factors, respectively.

The model imposes the bound on the \emph{forward} curve, not directly on yields. This is the Krippner-style architectural choice whose significance is explained in Section~\ref{sec:shadow-lit}. The immediate reason for distinguishing the two choices is that a yield is a maturity average of forwards. Bounding each forward and then averaging generally gives a different yield from averaging first and bounding the result. The survey regards bounding forwards before averaging as economically more natural. The bound enters through one of two scalar maps applied to the forward curve: the hard censoring map and its softplus smoothing,
\begin{equation}\label{eq:77}
 m_\ell(u)=\max\{\ell,u\},
 \qquad
 s_{\ell,\alpha}(u)=\ell+\alpha\log\!\left(1+e^{(u-\ell)/\alpha}\right),
 \qquad \alpha>0.
\end{equation}
The hard map raises values below the bound to the bound and leaves values above it unchanged. Softplus replaces the kink in this map with a smooth transition. The working example uses \(s_{\ell,\alpha}\) exclusively, with domestic constants
\((\lambda_d,\ell_d,\alpha_d)=(0.65,\,0,\,1.5\times10^{-3})\) and foreign constants
\((\lambda_f,\ell_f,\alpha_f)=(0.45,\,-0.005,\,1.0\times10^{-3})\). The foreign bound is a negative effective lower bound. The smoothing parameter \(\alpha\) is a fixed model constant, not an estimated parameter.

Observed zero-coupon yields are maturity averages of the bounded forward curve. The integral has no closed form once the softplus (or the max) is inside it, so it is evaluated by an order-\(K\) Gauss--Legendre rule on \([0,1]\), after rescaling the integration horizon by maturity. The working example fixes \(K=40\):
\begin{equation}\label{eq:78}
 y_c(\tau_i;x)=\frac1{\tau_i}\int_0^{\tau_i}
 s_{\ell_c,\alpha_c}\bigl(f_c(u;x)\bigr)\,du
 \;\approx\;\sum_{k=1}^{K}w_k\,
 s_{\ell_c,\alpha_c}\bigl(f_c(\tau_i v_k;x)\bigr),
\end{equation}
with six maturities per country. Each quadrature node selects a fixed fraction of the maturity at which to evaluate the bounded forward. Rescaling the integral cancels the maturity normalization, leaving the weighted sum shown above. Quadrature is suitable because the integrand in the softplus example is smooth and cheap to evaluate. A fixed high-order rule also keeps the observation map deterministic in \(x\), an essential property when differentiating through it.

The sole FX observation is the raw log forward with the observed log spot subtracted. The covered-interest-parity identity ties this observation to the two yield curves and the basis factors:
\begin{equation}\label{eq:79}
 \log F_t(\tau)-\log S_t=\tau\bigl(
 y_d(\tau;x_t)-y_f(\tau;x_t)+b(\tau;x_t^b)\bigr),
\end{equation}
where \(b(\tau;\cdot)\) is affine in the basis factors. The log forward--spot difference thus equals maturity times the sum of the domestic--foreign yield differential and the basis adjustment. All measurement noise is diagonal Gaussian with grouped scales.

In the implementation examined by the survey, the active inference route evaluates a marginal-likelihood approximation for \eqref{eq:75}--\eqref{eq:79} with a square-root unscented Kalman filter (UKF). It runs HMC and NeuTra-HMC over the nine-parameter chart using that filter's value and score. Two immediate consequences frame the case study. First, the resulting posterior is a \emph{smooth-model} posterior: \(s_{\ell,\alpha}\) is \(C^\infty\), so it has derivatives of every order. No discontinuity of any kind is present in the model actually being estimated. Second, the UKF likelihood is itself a deterministic, well-defined function of the parameters, so HMC on it is exact \emph{for that target}. Exact sampling of the UKF target leaves open how far its likelihood lies from the exact-model likelihood. Section~\ref{sec:shadow-filter} and Section~\ref{sec:shadow-routes} address that question.

\subsection{Where this sits in the shadow-rate literature}
\label{sec:shadow-lit}

This design is a reduced-form member of the shadow-rate term-structure family, and the survey's earlier machinery maps precisely onto that family. Black (1995) introduced the observation that currency makes the nominal short rate an option, \(r_t=\max\{0,r_t^\star\}\). The survey cites this result for orientation; its authors were unable to obtain the full text. Krippner (2012, Secs.~1--2) made the multi-factor Gaussian case tractable by bounding the \emph{forward} curve and obtaining yields through elementary numerical integration of the bounded curve. In that construction, the ZLB forward is the shadow forward plus a closed-form currency-option effect. The working example follows exactly this architecture: bound the forward curve, then integrate it to obtain yields. It uses softplus in place of Krippner's Gaussian-based map. Priebsch (2013) developed a cumulant-based arbitrage-free approximation for the same class.

Wu and Xia (2016, Sec.~2.3, eqs.~(6)--(8)) write the one-period forward observation as
\(\underline r+\sigma^Q\,g\!\left((a_n+b_n^{\mathsf T}X_t-\underline r)/
\sigma^Q\right)\), using the closed-form function \(g(z)=z\Phi(z)+\phi(z)\). They estimate the model by extended Kalman filtering and note that monotonicity of \(g\) keeps the likelihood surface well behaved. Christensen and Rudebusch (2015, Appendix~B) estimate arbitrage-free shadow-rate Nelson--Siegel models with an extended Kalman filter built on the discretized Ornstein--Uhlenbeck transition. Its conditional mean,
\((I-e^{-K\Delta})\theta+e^{-K\Delta}X_{t-1}\), has precisely the form in \eqref{eq:75}.

Kim and Priebsch (2013, Sec.~4, eqs.~(8)--(9)) replace the extended filter with the \emph{unscented} Kalman filter for quasi-maximum-likelihood estimation. They document that first-order linearization can be numerically unstable. Every earlier zero-bound study they name (Ichiue--Ueno, Kim--Singleton, and the 2013 Christensen--Rudebusch draft) had used the extended filter. Lemke and Vladu (2017) use the extended filter to estimate a euro-area shadow-rate model with a time-varying, estimated lower bound.

Finally, Opschoor and van der Wel (2024, Secs.~2.1--2.3, eqs.~(2)--(7)) build the reduced-form smooth shadow-rate DNS model closest to the working example. It uses AR(1) DNS factors and applies a lower bound through a smooth approximation of the max. Their menu comprises the Gaussian-based map
\(f_G(z)=z\Phi(z)+\phi(z)\), the softplus \(f_S\), and an inverse-exponential map, each with a sharpness scaling. They use two-step least-squares estimation and identify nonlinear state-space estimation via the extended Kalman filter as the alternative.

Three points of contrast matter for what follows. First, Opschoor and van der Wel bound the \emph{yield} directly,
\(y^Z(\tau)=r_{LB}+f(y(\tau)-r_{LB})\), whereas Krippner-type models, including the working example, bound the \emph{forward} curve and integrate. For hard censoring, pointwise convexity gives
\begin{equation}\label{eq:80}
 \sum_{k}w_k\max\{\ell,\,a(s_k)^{\mathsf T}x\}
 \;\geq\;
 \max\Bigl\{\ell,\;\sum_{k}w_k\,a(s_k)^{\mathsf T}x\Bigr\},
\end{equation}
with equality only when the forward curve does not cross the bound inside the maturity. The left side averages forwards after censoring each one; the right side censors their average. Forward-level censoring therefore produces weakly higher yields near the bound. The two models are not reparameterizations of each other.

Second, Opschoor and van der Wel defend smoothing as an \emph{economic} specification because yields empirically leave the bound gradually. This licenses interpreting \(\alpha\) as a structural smooth-transition constant rather than as a numerical tolerance. Under the economic interpretation, smoothing describes how yields leave the bound; under the numerical interpretation, it approximates a hard constraint. Section~\ref{sec:shadow-softplus} keeps these readings separate.

Third, the domain-standard estimators in the closure-filter line are all single-Gaussian moment-closure filters, either extended or unscented. These filters summarize state uncertainty with a single Gaussian and are precisely the filters whose branch-mass deficiency Section~\ref{sec:ukf-branchloss} quantifies. None of the shadow-rate papers examined by the survey in this line uses a mixture, particle, or exactly weighted likelihood, and none samples the posterior with gradient MCMC. These absence claims must be scoped to the closure-filter line.

Outside that line, Pericoli and Taboga (2018, Secs.~2, 4--5) estimate a Black-type max shadow-rate model,
\(r_t=\max\{\bar r_t,\,a+bX_t\}\), with a time-varying bound. Their fully Bayesian data-augmentation MCMC uses blockwise random-walk Metropolis over the parameters \emph{and} the entire latent factor path, with no Gaussian moment-closure filter of any kind. Bond prices are evaluated using a neural-network surrogate trained to sub-basis-point accuracy. Their ``nearly exact'' refers to that pricing approximation, not to the sampler. The exact-likelihood (closure-free) route in this literature is therefore already occupied. A UKF-plus-HMC route can claim a gradient-sampler frontier within the closure-filter family, but not a likelihood-fidelity frontier. The exact constructions below are measured against Pericoli--Taboga's closure-free standard, not merely against the filtering papers.

\subsection{A ladder of three targets and the geometry of the
hard-bound variants}
\label{sec:shadow-ladder}

Calling inference ``exact'' has no meaning until the target posterior is named. This shadow-rate model admits three natural versions with \emph{different posteriors}. The names introduced here are used throughout.

\(\mathcal S_1\) denotes the smooth softplus model, i.e.\ the finite program of \eqref{eq:77}--\eqref{eq:78} with \(K=40\) quadrature of the softplus map. This is the model of Section~\ref{sec:shadow-model} exactly as specified. \(\mathcal C_1\) denotes the finite hard-max model: it keeps the same quadrature nodes and weights but replaces \(s_{\ell,\alpha}\) by the hard max \(m_\ell\). \(\mathcal C_0\) denotes the continuous-maturity hard model, defined by the integral \(\tau_i^{-1}\int_0^{\tau_i} m_{\ell_c}(f_c(u;x))\,du\), evaluated by root-bracketed integration to a stated tolerance. These choices determine both how the lower bound is imposed and how the forward curve is averaged over maturity. Table~\ref{tab:ladder} summarizes the ladder and the meaning of exact inference for each target.

\begin{table}[htbp]\small
\begin{tabularx}{\textwidth}{@{}>{\raggedright\arraybackslash}p{0.16\textwidth}XX@{}}
\toprule
Variant & Definition & In what sense inference under it is exact \\
\midrule
\(\mathcal S_1\) (smooth softplus model) &
The finite program of \eqref{eq:77}--\eqref{eq:78}: \(K=40\)
Gauss--Legendre quadrature of the softplus-bounded forward curve. &
Smooth-model inference: any exact filter or sampler is exact for the
softplus posterior, never for a hard-bound posterior
(Section~\ref{sec:shadow-softplus}). \\
\(\mathcal C_1\) (finite hard-max model) &
The same quadrature nodes and weights with \(s_{\ell,\alpha}\)
replaced by the hard max \(m_\ell\). &
Exact relative to the finite hard-quadrature model: the bound is hard
but maturity averaging is the fixed \(K\)-node rule. \\
\(\mathcal C_0\) (continuous-maturity hard model) &
The continuous integral
\(\tau_i^{-1}\int_0^{\tau_i} m_{\ell_c}(f_c(u;x))\,du\), evaluated by
root-bracketed integration to a stated tolerance. &
Exact relative to the declared integral and tolerance: the bound and
the maturity average are both taken at the level of the continuous
model. \\
\bottomrule
\end{tabularx}
\caption{The ladder of three model variants. Each is a different
posterior; every likelihood, sampler, and result in this case study is
stated relative to one of them.}\label{tab:ladder}
\end{table}

\(\mathcal C_1\) is the natural first hard-bound counterfactual because it changes exactly one operation in the smooth model: the softplus becomes a hard maximum. However, \(\mathcal C_0\) and \(\mathcal C_1\) define different posteriors for a structural reason rather than a numerical one. Consider a simple transverse root of the bound gap, where the forward curve crosses the bound with nonzero slope. As the state changes, the crossing moves, changing the integration limits for the two branches. Differentiating these moving limits in the \(\mathcal C_0\) integral produces boundary terms with opposite signs. The terms cancel because the two integrand branches have the same value at the crossing. Consequently, \(y_c^{\mathcal C_0}(\tau_i;x)\) is continuously differentiable in \(x\) there. In contrast, the finite pointwise-max sum of \(\mathcal C_1\) has kinks whenever any quadrature node crosses the bound: the active branch at that node changes, and so does its slope.

Tangencies, endpoint roots, and a curve that lies identically at the bound are separate cases, each requiring its own treatment in the \(\mathcal C_0\) integration scheme. Root-aware \(\mathcal C_0\) evaluation and a \(\mathcal C_0\)-versus-\(\mathcal C_1\) discretization sensitivity experiment are prerequisites before any event- or kink-aware HMC design should be adopted as the final hard-ZLB sampling method. The case study specifies this experiment here but does not run it.

Replacing \(s_{\ell,\alpha}\) by \(m_\ell\) in \eqref{eq:78} gives \(\mathcal C_1\). At each quadrature node, the hard maximum selects either the fixed lower bound or the affine forward rate. Each yield is therefore a nonnegatively weighted sum of maxima of affine functions of the state,
\begin{equation}\label{eq:81}
 y_c^{\max}(\tau_i;x)=\sum_{k=1}^{K}w_k
 \max\bigl\{\ell_c,\;a_c(s_{ik})^{\mathsf T}x^c\bigr\},
 \qquad s_{ik}=\tau_i v_k,
\end{equation}
which is continuous, convex, and piecewise affine in \(x^c\). Continuity follows because the branches agree wherever they meet. Within a region where the active branch at every node stays the same, the sum is affine. Convexity follows because a maximum of affine functions is convex and nonnegative weighting preserves convexity. The FX row \eqref{eq:79} is a difference of these yield functions plus an affine basis term. It remains continuous and piecewise affine, though it is no longer convex.

The complete observation map \(h\) is thus continuous and piecewise affine, with kinks on the hyperplane arrangement \(\bigcup_{c,i,k}\{x:a_c(s_{ik})^{\mathsf T}x^c=\ell_c\}\). These are the states at which a quadrature-node forward rate meets the bound. The arrangement is a Lebesgue-null set: it has zero volume in state space. Slopes can change there even though the observation map itself does not jump.

The transition \eqref{eq:75} is a fixed linear-Gaussian map with no branch, and the observation noise is Gaussian with full support. The Gaussian observation density is a smooth function of its residual. Composing it with the continuous observation map gives likelihood values that agree across each bound crossing; within each affine region, the likelihood is smooth. Together with the specified transition, this means that the joint posterior over \((\vartheta,x_{0:T})\), equivalently over \((\vartheta,x_0,\eta_{1:T})\) in the non-centered chart of \eqref{eq:56}, is continuous, is smooth off a finite hyperplane arrangement, and has neither atoms nor density jumps.

In the classification of Section~\ref{sec:geometries}, these properties place the posterior in geometry one: a kink target, not a jump target. Ordinary Metropolized leapfrog HMC is therefore valid on the \(\mathcal C_1\) target by Section~\ref{sec:kink-validity}. The refraction machinery of Section~\ref{sec:event} degenerates because there is no potential-energy jump (\(\Delta U=0\)). Discontinuous or mixed HMC address structure that this model does not have. The \(\mathcal C_0\) target is smoother still at transverse roots, as the moving-boundary argument above shows; the geometry-one classification covers both hard-bound variants.

This classification applies to the model as specified. It is fragile in three specific directions, each corresponding to an extension a modeler may deliberately adopt. If factor dynamics switch at the bound through a binding-dependent \(\Phi\) or \(Q\), the transition acquires a genuine branch and the target moves to geometry two. Everything that Section~\ref{sec:geometries}, Section~\ref{sec:event}, and Section~\ref{sec:target-det} say about deterministic jumps then applies.

If a policy rate is added as an observation \emph{without} measurement error, with \(i_t=\max\{\ell,f_c(0;x_t)\}\) observed exactly, all underlying rates below the bound produce the same observed rate. The observation law therefore acquires an atom, or point mass, at \(\ell\). The likelihood becomes the censored mixture of Section~\ref{sec:ukf-censored} with \(V\to0\).

If the lower bound \(\ell\) is estimated, as in Lemke and Vladu (2017), the kink location becomes parameter-dependent. Continuity remains intact, but the kink arrangement moves with \(\vartheta\). A modeler must decide which of these extensions, if any, is included. The geometry-one claim for the specified model does not automatically carry over to any of them.

One structural lemma sharpens the analysis that follows by restricting how the forward curve can meet the bound. For a fixed country factor state \(x^c\), the bound gap is the forward rate minus the country's lower bound. Its value and its derivative with respect to the horizon are
\begin{equation}\label{eq:82}
 g(s)=f_c(s;x)-\ell_c=(L-\ell_c)+e^{-\lambda_c s}(S+C\lambda_c s),
 \qquad
 g'(s)=\lambda_c e^{-\lambda_c s}\bigl(C-S-C\lambda_c s\bigr),
\end{equation}
where \((L,S,C)=x^c\). If \(L-\ell_c=0\), \(S=0\), and \(C=0\), every term in the gap vanishes. The gap is then identically zero, so the curve lies exactly on the bound at every horizon. This degenerate case must be handled by convention. It is a Lebesgue-null event under the Gaussian predictive distribution, although synthetic test cases can be constructed to land on it.

Outside this degenerate case, the bracket in \(g'\) is affine in \(s\). The factor multiplying the bracket does not change sign with the horizon, so \(g'\) changes sign at most once. Thus \(g\) is monotone or has a single peak or trough on \((0,\infty)\) and therefore has \emph{at most two zeros}. A nonconstant monotone gap can meet zero at most once; a single peak or trough allows at most one meeting on either side of the turning point; and a nonzero constant gap has no zeros. The binding set \(\{s:f_c(s;x)\leq\ell_c\}\) consists of the horizons at which the gap is nonpositive. Along any Nelson--Siegel forward curve, this set is an interval, the complement of an interval, empty, or the entire horizon domain.

Two consequences follow. First, the reachable binding patterns over the pooled quadrature horizons \(\{s_{ik}\}\) are \emph{interval patterns}: when the horizons are ordered, the binding points form a consecutive block or its complement, including the empty and full cases. Such patterns are determined by at most two boundary locations, so there are \(O(\text{grid}^2)\) of them rather than \(2^{\text{grid}}\), the number of unrestricted binary patterns. Second, the per-country kink cells in factor space are polyhedra: fixing a binding pattern imposes linear inequalities on the country's factors. The constraint normals \(a_c(s)\) span at most that country's three factor dimensions, so per-country cell probabilities reduce to at most three-dimensional Gaussian integrals.

The evidence decomposition below requires the \emph{joint} two-country cell probability, which is generically six-dimensional. The FX row and the measurement update correlate the two factor blocks, so the joint cell mass cannot in general be obtained by multiplying the separate country masses. It factorizes into two three-dimensional integrals only if a factorization of the relevant covariance is proved for the case at hand.

\subsection{An exact cell identity and the bootstrap filter that carries
the likelihood}
\label{sec:shadow-filter}

The observation map of \(\mathcal C_1\) is exactly affine on each polyhedral cell of the kink arrangement. This suggests a beautiful idea: with a Gaussian one-step predictive distribution, Gaussian conjugacy can be applied within each cell, and the cell contributions can be combined into an exact filtering update. This section develops that identity in full because it is exact and organizes everything that can be learned about the model's local structure. The argument proceeds from the ideal identity to three specific obstacles to its implementation as an exact algorithm, then to the robust alternative used for inference. That alternative is the plain bootstrap particle filter. It is stated first so that the reader knows the practical conclusion before examining the cell construction.

For \(\mathcal C_1\), the workhorse likelihood estimator is the \textbf{bootstrap particle filter}. For each particle, it draws the state from the exact Gaussian transition conditional on the selected ancestor, then weights the state by the exact observation density:
\begin{equation}\label{eq:82a}
 x_t^j\sim N\bigl((I-\Phi)\bar\theta+\Phi x_{t-1}^{A_t^j},\,Q\bigr),
 \qquad
 v_t^j=N\bigl(y_t;\,h_{\max}(x_t^j),\,R\bigr),
\end{equation}
where \(h_{\max}\) stacks the \(\mathcal C_1\) yield rows \eqref{eq:81} and the FX row \eqref{eq:79}. Both transition sampling and observation-density evaluation are available in closed form. The likelihood estimator \eqref{eq:13} is nonnegative and unbiased under the standard support and integrability conditions of Section~\ref{sec:pf-generic}. With this estimator, PMMH has the ordinary particle-MCMC extended-target guarantee of Section~\ref{sec:pm-target}. Although less elegant than the cell decomposition, the bootstrap filter is implementable today and, under those conditions, unbiased for the named finite model. Its Monte Carlo variance can be measured directly from likelihood replicates.

Now the ideal. A binding pattern \(B\) assigns bind or slack to every quadrature horizon. Once this pattern is fixed, the observation map is affine, \(h_B(x)=H_Bx+d_B\), on the corresponding polyhedral cell
\begin{equation}\label{eq:83}
 C_B=\bigl\{x:\;a_c(s_{ik})^{\mathsf T}x^c\leq\ell_c\ \text{for}\
 (c,i,k)\in B,\ \ >\ell_c\ \text{otherwise}\bigr\}.
\end{equation}
The binding inequalities include equality; the slack inequalities are strict. Let the one-step predictive distribution be Gaussian, \(x_t\mid\mathcal F\sim N(m,P)\). This specification is exact conditional on an individual particle's ancestor, since the transition \eqref{eq:75} gives \(x_t\mid x_{t-1}\sim N((I-\Phi)\bar\theta+\Phi x_{t-1},Q)\).

To obtain the posterior, multiply the Gaussian predictive density by the observation likelihood within each cell. Because the observation map there is affine, the product factors into a Gaussian observation-density term and an updated Gaussian state density. Restricting each contribution to its cell and summing over cells gives the exact one-step posterior without approximation:
\begin{equation}\label{eq:84}
 p(x_t\mid y_t,\mathcal F)\;\propto\;
 \sum_{B}\mathbf 1\{x_t\in C_B\}\,
 N\!\bigl(y_t;H_Bm+d_B,\,S_B\bigr)\,
 N\!\bigl(x_t;m_B^+,P_B^+\bigr),
\end{equation}
where the indicator confines each Gaussian contribution to the cell on which its observation map applies. Each cell uses the standard conjugate update
\begin{equation}\label{eq:85}
 S_B=H_BPH_B^{\mathsf T}+R,\quad
 K_B=PH_B^{\mathsf T}S_B^{-1},\quad
 m_B^+=m+K_B(y_t-H_Bm-d_B),\quad
 P_B^+=P-K_BS_BK_B^{\mathsf T},
\end{equation}
in which \(S_B\) is invertible for every cell because the measurement covariance \(R\) is positive definite and \(H_BPH_B^{\mathsf T}\) is positive semidefinite. The update combines projected state uncertainty with measurement noise to obtain the predictive observation covariance. The Kalman gain then converts the observation residual into a correction to the predictive state mean, while the covariance update subtracts the corresponding reduction in uncertainty. These are the parameters of the Gaussian factor; the posterior still restricts that factor to its cell.

Integrating the posterior contributions before normalization gives the exact one-step evidence. Each cell contributes its Gaussian observation-density factor multiplied by the probability that its updated Gaussian lies in that cell. Summing these contributions gives the polyhedrally truncated mixture mass
\begin{equation}\label{eq:86}
 p(y_t\mid\mathcal F)=\sum_B Z_B,
 \qquad
 Z_B=N\!\bigl(y_t;H_Bm+d_B,\,S_B\bigr)\,
 \Pr_{x\sim N(m_B^+,P_B^+)}\!\bigl(x\in C_B\bigr).
\end{equation}
Equations \eqref{eq:84}--\eqref{eq:86} are the hard-bound observation-map counterpart of the branch-split update \eqref{eq:36}--\eqref{eq:45}. The branch structure has moved from the transition maps \(F_r\) into the observation cells \(C_B\). The truncated moments of \eqref{eq:21}--\eqref{eq:28} reappear whenever a single constraint dominates. By the crossing lemma of Section~\ref{sec:shadow-ladder}, the sum over \(B\) runs over interval patterns only.

Inside a particle filter, this identity yields an \emph{ideal} fully adapted construction. For particle \(j\) with ancestor \(x_{t-1}^j\), the predictive distribution is exactly Gaussian. Thus \eqref{eq:84}--\eqref{eq:86} hold exactly with \((m,P)=((I-\Phi)\bar\theta+\Phi x_{t-1}^j,\,Q)\). In the ideal construction, the conditionally optimal proposal is the posterior conditional on that ancestor and the current observation. Its incremental weight is the corresponding predictive observation density:
\begin{equation}\label{eq:87}
 q^\star(x_t\mid x_{t-1}^j,y_t)=p(x_t\mid x_{t-1}^j,y_t),
 \qquad
 v_t^j=p(y_t\mid x_{t-1}^j)=\sum_B Z_B\bigl(x_{t-1}^j\bigr),
\end{equation}
where sampling first selects a cell with probability \(\propto Z_B\), normalized over the cells, and then draws from the cell-restricted Gaussian \(N(m_B^+,P_B^+)\mathbf 1_{C_B}\). Conditional on the ancestor, the incremental weight in \eqref{eq:87} does not depend on the newly sampled state. This is the fully adapted property of \eqref{eq:12b}--\eqref{eq:13}. In particular, \emph{if} every \(Z_B\) is computed exactly and the cell-restricted draw is exact, the resulting likelihood estimator is nonnegative and unbiased for the \(\mathcal C_1\) model, by the argument of Section~\ref{sec:pf-generic}.

Three numerical obstacles explain why this ideal construction is not the algorithm used for likelihood estimation. Each concerns implementation rather than a conceptual flaw in the identity. First, the cell masses \(\Pr_{x\sim N(m_B^+,P_B^+)}(x\in C_B)\) are multivariate Gaussian polytope probabilities. They are generically six-dimensional across the two correlated factor blocks, as discussed with the crossing lemma. Evaluating them requires numerical computations in the sense of Genz and Bretz (2009), rather than closed-form operations.

Second, restricting the sum to patterns with ``non-negligible'' mass is a \emph{pruning approximation}. Omitting any cell with nonzero mass changes \(\sum_B Z_B\), the cell-selection probabilities, and the incremental weight. The pruned construction is exact only if the omitted masses are certified to be exactly zero.

Third, exact draws from the cell-restricted Gaussian are themselves nontrivial. Rejection sampling is exact in principle, but can be unusably inefficient in the tails. The exact piecewise-quadratic HMC of Pakman and Paninski (2014) supplies an \emph{invariant kernel} for the truncated Gaussian: the kernel preserves that distribution, but a finite trajectory does not supply an independent exact draw. Substituting one trajectory for an exact cell-restricted draw therefore changes the proposal law and forfeits the constant weight \eqref{eq:87}. Botev's (2017) minimax-tilting method is the modern benchmark for both the polytope probabilities and the cell-restricted Gaussian draws. It nevertheless carries its own numerical error, which must be quantified and incorporated into the proposal before a cell implementation can be described as exact.

A related limitation applies beyond the particle setting. Equations \eqref{eq:84}--\eqref{eq:86} do \emph{not} give a closed-form multi-step filter. The Gaussian predictive specification above is exact conditional on a particle's ancestor; propagating an entire cell-restricted Gaussian through \eqref{eq:75} leaves the Gaussian family. An exact density filter must therefore carry a growing truncated-mixture representation. Any pruning or moment collapse \eqref{eq:47}--\eqref{eq:48} is a declared approximation, exactly as in Section~\ref{sec:ukf-branchwise}.

Until these numerical components are controlled, the cell construction remains a \emph{proposal
and diagnostic layer}. An implementation that prunes cells or approximates draws must use its actual proposal density and the corresponding importance ratio in \eqref{eq:12b}, rather than the ideal constant weight \eqref{eq:87}. The bootstrap filter in \eqref{eq:82a} therefore supplies the likelihood estimator used for inference, as announced at the outset. This positioning transplants the Aruoba et al.\ conditionally-optimal construction from a piecewise-linear transition to a hard-bound observation map. The construction is simpler here because the transition never branches, but it still requires the same control of numerical error.

The relevance of the Tobit literature on censored filtering depends on
whether noise is added before or after censoring. In the
Tobit-type-1 measurement model surveyed by
Geng et al.\ (2021), censoring is applied
\emph{after} the noise: \(y=\max\{c,\,Hx+u\}\).
All noisy values below the censoring limit are mapped to that limit,
creating an atom: positive probability at a single observed value.
This is the setting of the Tobit Kalman filter (Allik et al.\
2016; the authors of the present manuscript cite this work from metadata
because they were unable to obtain its full text).
The present model adds noise to yields \emph{after} the max operation.
That noise spreads observations around the censored values, so the
observed yields have no atom. The mixture structure lies in the state
rather than in the observation support. The two orderings coincide
only in the exactly observed policy-rate extension of
Section~\ref{sec:shadow-ladder}, where the Tobit machinery
would become the appropriate reference.

\subsection{The softplus route: bias bounds and two readings of
\texorpdfstring{$\alpha$}{alpha}}
\label{sec:shadow-softplus}

The discrepancy between the smooth model
\(\mathcal S_1\) and the hard model
\(\mathcal C_1\) can be quantified starting with
the exact pointwise gap between their scalar maps:
\begin{equation}\label{eq:88}
 0\;<\;s_{\ell,\alpha}(u)-m_\ell(u)
 =\alpha\log\!\left(1+e^{-|u-\ell|/\alpha}\right)
 \;\leq\;\alpha\log 2,
\end{equation}
so smoothing raises the scalar map at every input, by no more than the
stated uniform upper bound. The maximum is attained exactly at the
kink, and the gap decays exponentially away from it. For example, when
\(|u-\ell|\geq5\alpha\), the gap is below
\(6.8\times10^{-3}\,\alpha\).

Averaging the scalar inequality \eqref{eq:88}
through the quadrature \eqref{eq:78} preserves
its bounds: each scalar gap is positive and no larger than the same
upper limit. The resulting uniform yield bound is
\begin{equation}\label{eq:89}
 0\;<\;y_c^{\alpha}(\tau_i;x)-y_c^{\max}(\tau_i;x)\;\leq\;\alpha_c\log 2,
\end{equation}
with the upper bound attained only when the whole forward path sits
at the kink. In general, the gap is
\(\alpha\log2\) times the quadrature mass within
\(O(\alpha)\) of the bound, plus an exponentially
small remainder. The quadrature mass measures how much weight the
averaging places near the bound.

With the working example's constants, the worst-case yield gaps are
\(1.04\times10^{-3}\) for domestic yields and
\(0.69\times10^{-3}\) for foreign yields in annualized
rate units, roughly ten and seven basis points, respectively.
The gap is \emph{one-signed}: softplus yields always exceed
hard-censoring yields at the same factors. To offset this increase
near binding, fitted factors and long-run means absorb a systematic
downward shift. This adjustment responds to a directional change in
the yield map rather than to symmetric noise.

For the likelihood comparison, the mean shift is the smooth observation
map minus the hard observation map. The per-observation log-likelihood
perturbation obeys
\begin{equation}\label{eq:90}
 \bigl|\log N(y;h_\alpha,R)-\log N(y;h_{\max},R)\bigr|
 \;\leq\;
 \bigl\|R^{-1/2}(y-h_{\max})\bigr\|\,
 \bigl\|R^{-1/2}\Delta\bigr\|
 +\tfrac12\bigl\|R^{-1/2}\Delta\bigr\|^2,
\end{equation}
with \(0\le\Delta_i\le\alpha\log2\) componentwise on the yield
rows. With this componentwise bound, the stated perturbation is small
per observation. The likelihood bound combines the norms of the
residual and the mean shift, both standardized by the measurement-error
covariance, and adds a quadratic term in the standardized shift.

Defining the mean-map shift as
\(\Delta=h_\alpha-h_{\max}\), the signed log-likelihood
difference is
\begin{equation}\label{eq:90a}
 \log N(y;h_\alpha,R)-\log N(y;h_{\max},R)
 =(y-h_{\max})^{\mathsf T}R^{-1}\Delta
 -\tfrac12\,\Delta^{\mathsf T}R^{-1}\Delta,
\end{equation}
which can take \emph{either sign} even though the map shift
\(\Delta\) is one-sided on the yield rows.
Its sign depends on the residual
\(y-h_{\max}\): the residual-weighted term can be
positive or negative, while the quadratic term is subtracted.
The FX row's \(\Delta\) component is a
domestic-minus-foreign combination, so even the mean shift for that
row can have either sign. The nonnegativity bound for yield rows
therefore does not extend to the FX row.

The one-sided \emph{mean-map} discrepancy on the yield rows may induce
a systematic fitted-factor adjustment near binding. Likelihood and
posterior perturbations, however, can take either sign and should be
reported as computed signed quantities under
\eqref{eq:90}--\eqref{eq:90a}.
With measurement-error standard deviations of a few basis points,
a ten-basis-point systematic map gap is not negligible relative to
the noise scale. Equations
\eqref{eq:90}--\eqref{eq:90a}
quantify its effect in log-likelihood units, precisely the units in
which it affects inference.

Sharpening the softplus by reducing \(\alpha\) may shrink the gap between the smooth and hard-bound posteriors, but it carries a cost for gradient samplers. The scalar curvature is
\begin{equation}\label{eq:91}
 s_{\ell,\alpha}''(u)=\frac{\sigma(z)\bigl(1-\sigma(z)\bigr)}{\alpha}
 \leq\frac1{4\alpha},
 \qquad z=\frac{u-\ell}{\alpha},
\end{equation}
where the rescaled distance measures how far the input lies from the bound in units of the smoothing width. The logistic product is largest at the bound and falls away on either side. Reducing the smoothing width therefore produces ridges of height \(O(1/\alpha)\) in the observation Hessian, concentrated in an \(O(\alpha)\)-neighborhood of the kink. These ridges become taller and narrower, so leapfrog stability suggests a local step size \(\epsilon=O(\sqrt\alpha)\) near the kink.

This scaling is a local stability heuristic under a fixed metric and comparable posterior curvature. It is not a universal necessity theorem: the step size that actually limits the sampler depends on the choice of coordinates (the chart), the mass matrix, the trajectory, and how much time the sampler spends near the bound. These qualifications leave the qualitative comparison intact. The \(\alpha\to0\) limit of the smooth route is \emph{harder} for HMC than the \(\mathcal C_1\) hard-bound model itself. The hard-bound model has bounded curvature on either side of the kink, and its kink cost is the \(O(\epsilon)\) energy error in \eqref{eq:56a}. Sharpening the softplus to approximate the hard bound is therefore the wrong direction. The appropriate choices are either a genuinely smooth model or the exact machinery of Section~\ref{sec:shadow-filter}.

There are two defensible readings of \(\alpha\), and they support different claims. Under the first, the softplus is a \emph{numerical approximation} to the hard-ZLB model, and the choice of \(\alpha\) requires an error protocol. That protocol requires reporting the quantities in \eqref{eq:89}--\eqref{eq:90} alongside binding fractions and rerunning matched synthetic-data estimations at \(\alpha/2\) and \(2\alpha\). The experiments in this survey implement this comparison by rerunning matched synthetic-data estimations at halved and doubled \(\alpha\). Before any hard-ZLB language is used, the posterior functionals of interest must move by less than their Monte Carlo uncertainty in these comparisons.

Under the second reading, the softplus defines a \emph{structural smooth-transition model}, as argued empirically by Opschoor and van der Wel (2024). Their evidence is that yields leave the bound gradually, so the smooth map may be the better description. In this reading, \(\alpha\) is a model constant or even an estimable parameter. No bias language applies, because the smooth model itself is the target; every claim must name it as such. Both readings rule out the hybrid of sampling the smooth model while describing the result as inference under the ZLB: the two posteriors differ by exactly the gaps computed above.

\subsection{Identification at the bound}
\label{sec:shadow-ident}

The bound both complicates the likelihood's geometry and reduces the information that yields provide about the shadow factors. Applying the chain rule to \eqref{eq:78} gives the yield sensitivity to the factors,
\begin{equation}\label{eq:92}
 \frac{\partial y_c(\tau_i;x)}{\partial x^c}
 =\sum_{k=1}^{K}w_k\,
 \sigma\!\left(\frac{f_c(s_{ik};x)-\ell_c}{\alpha_c}\right)
 a_c(s_{ik}),
\end{equation}
where each weighted factor loading is multiplied by the logistic slope of the softplus. The logistic factors are replaced by indicators \(\mathbf 1\{f>\ell\}\) in the \(\mathcal C_1\) hard-bound model. These slopes determine how strongly changes in the shadow factors affect yields. When the forward curve lies far below the bound over the horizons relevant to maturity \(\tau_i\), the logistic slopes tend to zero, while the hard-bound indicators are zero. The corresponding rows of the observation Jacobian \(J\) therefore vanish, and the per-period observed Fisher information \(J^{\mathsf T}R^{-1}J\) about the shadow factors degenerates in exactly the binding directions. During a long binding spell, the prior and the transition \eqref{eq:75} alone drive the shadow-factor posterior. The data inform the binding country's long-run means \(\bar\theta\) only through occasional excursions of the forward curve above the bound.

These identification limits are the formal counterpart of two empirical findings. Bauer and Rudebusch (2016) document that shadow-rate estimates are highly sensitive to the assumed lower-bound value and to model specification. Christensen and Rudebusch (2015) report the same sensitivity for shadow-rate paths. The practical consequences follow directly. Priors on \(\bar\theta\) shape inference during binding spells, so reporting those priors is part of reporting the result.

Identification degeneracy also makes the sampler ill-conditioned in exactly the same binding directions. The loss of information changes the posterior geometry: the HMC mass matrix appropriate during binding spells is very different from the one appropriate off the bound. Step-size and mass-matrix tuning performed on off-bound data, or on another model, horizon, or particle count, will therefore not transfer to binding regimes. That tuning must be redone for the configuration actually being run. Every empirical claim from this model must also be accompanied by identification diagnostics along the binding directions, where the data cease to inform inference and the prior takes over.

\subsection{Five routes and what each delivers}
\label{sec:shadow-routes}

The case study's conclusions give five inference routes. They differ in their target on the ladder in Table~\ref{tab:ladder} and in the strength of their guarantees, as well as in cost. Table~\ref{tab:routes} summarizes what each delivers.

\begin{table}[htbp]\small
\begin{tabularx}{\textwidth}{@{}>{\raggedright\arraybackslash}p{0.16\textwidth}XX@{}}
\toprule
Route & Target & What inference under this route delivers \\
\midrule
(i) SR-UKF marginal likelihood + HMC/NeuTra &
The UKF likelihood as a stated approximate functional of
\(\mathcal S_1\) &
Exact MCMC for its own target, the UKF likelihood rather than the model
likelihood; this mirrors the estimator of Kim and Priebsch (2013). Its
single-Gaussian closure discards branch mass exactly as in
Section~\ref{sec:ukf-branchloss}; its distance from route (iii) on matched
binding scenarios is a required diagnostic, not an assumption. \\
(ii) Softplus model with any exact filter/sampler &
\(\mathcal S_1\) &
Legitimate as a structural model (the Opschoor--van der Wel reading)
or as an approximation under the \eqref{eq:89}--\eqref{eq:90a}
protocol; never a hard-ZLB claim by itself. \\
(iii) Bootstrap hard-bound particle filter \eqref{eq:82a} + PMMH &
\(\mathcal C_1\) &
The likelihood workhorse: exact Gaussian transition draws, exact
observation-density weights, a nonnegative unbiased estimator, and
exact extended-target MCMC by Section~\ref{sec:pm-target}. Delivers valid
inference for the finite hard-quadrature posterior. Supports PMMH, not
pseudo-marginal HMC, absent the differentiability conditions of
Section~\ref{sec:pmhmc}; correlated pseudo-marginal (Deligiannidis, Doucet,
and Pitt 2018) is the variance/mixing comparator, not a
differentiability fix. \\
(iv) Joint \((\vartheta,x_0,\eta_{1:T})\) HMC on \(\mathcal C_1\) &
The \(\mathcal C_1\) model, sampled in the non-centered chart of
\eqref{eq:56} &
Valid by Section~\ref{sec:kink-validity} because the target is a continuous
piecewise-quadratic-plus-smooth kink target; needs no filter at all.
Gradient cost is one simulation sweep, and the linear-Gaussian prior
structure supplies the natural preconditioner. Kink crossings cost
acceptance per \eqref{eq:56a}, not correctness. \\
(v) Cell-adapted fully adapted proposal
\eqref{eq:84}--\eqref{eq:87} inside a particle filter &
\(\mathcal C_1\) &
Conditional: delivers exact fully adapted filtering only once the
Genz--Bretz polytope-probability and Botev truncated-draw components
are benchmarked and carried with quantified numerical error; with
pruning or approximate draws it must use its actual proposal density
and importance ratio, and until then it is a proposal/diagnostic layer
for route (iii). \\
\bottomrule
\end{tabularx}
\caption{Five inference routes for the two-country shadow-rate model,
the ladder target each addresses, and what inference under each
delivers.}\label{tab:routes}
\end{table}

The branch-split mixture UKF retains the role assigned in Section~\ref{sec:ukf}: it is a fast diagnostic and a proposal ingredient. Route (iii) carries the likelihood of record. Routes (iii) and (iv) use PMMH and joint HMC to target the same \(\mathcal C_1\) model. Their results on synthetic data must therefore agree within Monte Carlo error. This comparison is the cheapest available end-to-end correctness test for the model, and it can be run as a small parameter-recovery experiment before any full-data run. The continuous-maturity target \(\mathcal C_0\) enters through the discretization sensitivity experiment of Section~\ref{sec:shadow-ladder}, before any event- or kink-aware design is adopted.

\section{Case study II: a New Keynesian model with an occasionally binding policy rate}
\label{sec:nk}

The second case study turns from term structures to a small New Keynesian economy whose policy rate is truncated at a lower bound and whose equilibrium is computed by a piecewise-linear method. This model class motivates most of the ZLB estimation literature. Likelihoods built on these solutions have a reputation among practitioners for being ``spiky'': reported symptoms include jagged likelihood surfaces, optimizers that stall on cliffs, and gradient-based samplers that reject in bursts. Such behavior makes it tempting to blame the $\max$ operator in the policy rule. The analysis of this section shows why that diagnosis is misplaced.

With a unique verified solution, the piecewise-linear ZLB model is a \emph{kink} target in states and shocks: it is continuous and piecewise smooth. This geometry is well within the reach of the methods of Section~\ref{sec:hmc}. The genuine discontinuities practitioners encounter arise from the solution \emph{correspondence}, which gives the set of verified regime paths for given inputs. That set may be empty because no verified path exists, or it may contain multiple paths. The selection rule imposed on the correspondence is another source of discontinuity. The section derives this distinction in full on the smallest model that exhibits it, then draws out the implications for how an estimation exercise must be posed.

\subsection{Solution operators for the piecewise-linear class}
\label{sec:nk-solvers}

The benchmark model class is the three-equation New Keynesian model with a
Taylor rule truncated at the bound, written in the notation of
\eqref{eq:1}--\eqref{eq:3} and solved by a piecewise-linear method. Three
solution operators dominate the literature for this model class. OccBin
(Guerrieri and Iacoviello 2015, Sec.~2) begins by guessing a regime
path, specifying when the bound binds. It solves the two linear systems
backward from a terminal reference-regime date and then simulates forward.
It revises the guess and repeats these steps until the simulated path satisfies
the inequalities defining the guessed regimes. The conditional filtering
counterpart of this construction is given in \eqref{eq:6}--\eqref{eq:11}.
Holden (2016) reformulates the bound as anticipated news shocks added
to the unconstrained linear solution and implements the search in DynareOBC.
The anticipated-shock approach originates with Holden and Paetz
(2012), cited here for orientation. Boehl (2022) derives a
closed-form expression for the whole trajectory conditional on guessed spell
durations. Because those guesses determine the trajectory through the
closed-form expression, the iteration reduces to a fast search over two
integers. When they converge to a verified path, all three methods produce the
same kind of solution: one that is linear on each regime cell of the
state-shock space, meaning each region where the same regime pattern applies.

Holden (2023, Secs.~2--3) supplies algebra for analyzing the solution
correspondence. This algebra deserves attention because it turns a vague worry
about ``convergence failures'' into precise questions about whether a verified
path exists and whether it is unique. The vector
$q\in\mathbb R^{T}$ stacks the bound violations of the unconstrained solution,
and the matrix $M\in\mathbb R^{T\times T}$ collects the responses of the bounded variable
to unit news shocks. Imposing the bound over a horizon of
$T$ periods is then the linear complementarity problem (LCP)
\begin{equation}\label{eq:93}
 y\geq0,\qquad q+My\geq0,\qquad y^{\mathsf T}(q+My)=0,
\end{equation}
whose solution $y$ gives the news-shock loadings that hold the
variable at the bound during the binding spell. The first inequality requires
nonnegative shock loadings, and the second enforces the bound after the shock
responses have been added. Because both vectors are nonnegative, the zero
inner product requires their product to be zero in every period. Thus a
positive loading can occur only when the variable is at the bound; when the
variable is strictly above the bound, the corresponding loading must be zero.

Existence of a verified path is exactly equivalent to existence of a solution
to \eqref{eq:93}, and uniqueness of the path is exactly equivalent to
uniqueness of that solution. Holden's central result is that the LCP has a
unique solution for \emph{every} $q$ if and only if
$M$ is a P-matrix, meaning that all its principal minors are
positive. Holden also shows that standard New Keynesian models fail this
condition and have multiple perfect-foresight paths that escape the bound.

The P-matrix property establishes uniqueness at each parameter point.
Continuity concerns a separate question: how the solution moves
\emph{between} parameter points. Pointwise uniqueness of the LCP solution does
not, by itself, establish continuity of the solution map
$\vartheta\mapsto y(q(\vartheta),M(\vartheta))$. The usable regularity statement is parametric.
For a \emph{fixed finite horizon} $T$ and a fixed terminal condition,
suppose that $q(\vartheta)$ and $M(\vartheta)$ are continuous on a stated
parameter domain and that $M(\vartheta)$ remains a P-matrix throughout that
domain. Standard parametric-LCP regularity then implies that the unique LCP
solution is continuous in $\vartheta$.

Changing the horizon, changing the terminal condition, or losing the P-matrix
property on the domain removes this continuity guarantee. Any of these changes
can create a genuine discontinuity even when the solution is unique at each
parameter point. An estimation exercise should therefore state which
conditions of the parametric continuity result it has verified and on which
parameter domain.

\subsection{A fully derived multiplicity example and the discontinuity mechanism}
\label{sec:nk-mult}

Holden's (2023, Sec.~2.2) lagged-response economy is the smallest model that shows the multiplicity and discontinuity mechanism examined here. The survey reproduces it in full because the rest of the case study relies on it. With a constant real rate $r>0$, combining the Fisher equation with a truncated Taylor rule that responds to current and lagged inflation gives
\begin{equation}\label{eq:94}
 r+\pi_{t+1}=i_t=\max\{0,\;r+\phi\pi_t-\psi\pi_{t-1}\},
 \qquad \phi-\psi>1,\ \psi\in(0,1),
\end{equation}
where initial inflation $\pi_0$ is given. The maximum enforces the zero lower bound on the nominal interest rate. Away from the bound, the rule is linear and the stable solution is $\pi_t=A\pi_{t-1}$. Substituting this form into the unconstrained equation yields the characteristic condition $A^{2}-\phi A+\psi=0$. Taking $\phi=2$ for transparency gives the stable root $A=1-\sqrt{1-\psi}\in(0,1)$. Rearranging the characteristic condition gives the useful identity $A(\phi-A)=\psi$, which will put the branch admissibility conditions in a common form.

Holden shows that any path that escapes the bound can be at it for at most the first period. Its subsequent evolution therefore follows the stable solution, so $\pi_2=A\pi_1$. Substituting this relation into \eqref{eq:94} at $t=1$ and subtracting the implied nominal interest rate from both sides gives a scalar instance of \eqref{eq:93}:
\begin{equation}\label{eq:95}
 0=\max\bigl\{-r-A\pi_1,\;(\phi-A)\pi_1-\psi\pi_0\bigr\}.
\end{equation}
For the maximum to be zero, both entries must be nonpositive and at least one must be zero. Both branches can therefore be solved explicitly. On the slack branch, setting the second entry to zero gives the fundamental solution $\pi_1^{\mathrm f}=\psi\pi_0/(\phi-A)$. The first entry must then be nonpositive, so this solution is admissible when $r+A\pi_1^{\mathrm f}\geq0$. Using the identity above reduces this condition to $\pi_0\geq-r/A^{2}$. On the binding branch, setting the first entry to zero gives $\pi_1^{\mathrm b}=-r/A$. Requiring the second entry to be nonpositive gives the admissibility condition $(\phi-A)\pi_1^{\mathrm b}\leq\psi\pi_0$. The identity reduces this to the \emph{same} condition $\pi_0\geq-r/A^{2}$. The branch solutions coincide at the shared threshold and are distinct above it; below it, neither branch is admissible. Thus
\begin{equation}\label{eq:96}
\begin{aligned}
 \pi_0>-\frac{r}{A^{2}}&:\ \text{two solutions},\\
 \pi_0=-\frac{r}{A^{2}}&:\ \text{one},\\
 \pi_0<-\frac{r}{A^{2}}&:\ \text{none returning to the standard steady state}.
\end{aligned}
\end{equation}

Multiplicity therefore covers an open half-line of initial states rather than a knife-edge. Holden proves that it becomes pervasive in richer New Keynesian models.

The source of discontinuity can now be located. Each branch map is \emph{continuous} in $(\pi_0,r,\psi)$: $\pi_1^{\mathrm f}$ is linear in $\pi_0$, while $\pi_1^{\mathrm b}$ is constant with respect to initial inflation. The discontinuity is created entirely by \emph{selection}. Any deterministic rule $\sigma$ that picks a branch defines a transition kernel
\begin{equation}\label{eq:97}
 \pi_1=T_{\sigma(\pi_0,\vartheta)}(\pi_0,\vartheta),
 \qquad
 T_{\mathrm f}\neq T_{\mathrm b}\ \text{on the whole multiplicity region},
\end{equation}
which jumps by an $O(1)$ amount wherever $\sigma$ switches branch within the multiplicity region. The jump reflects the gap between the distinct branch values at the switch. A likelihood built on \eqref{eq:97} inherits this jump as a function of the state. The multiplicity boundary $-r/A(\psi)^{2}$ and any solver-dependent selection also move with the parameters $\vartheta$, so the likelihood inherits jumps as a function of the parameters as well.

These state and parameter jumps give formal content to the folklore that piecewise-linear ZLB likelihoods ``spike''. The $\max$ operator does not cause the spike, since its solution branches are individually continuous. The cause is an undeclared selection rule, possibly dependent on solver order, acting on a multiplicity region of positive measure, together with nonexistence of solutions returning to the standard steady state below the threshold. In the taxonomy of Section~\ref{sec:geometries}, the example is geometry four until a selection law is declared: the equilibrium conditions leave the branch unspecified. A deterministic declaration makes it geometry two along the selection switch set, where the selected transition jumps. Only the stochastic completion in the next subsection makes it geometry three, where the mixed-support kernels of Section~\ref{sec:mixed} and Section~\ref{sec:pgas} legitimately apply.

\subsection{The stochastic completion and the sunspot literature}
\label{sec:nk-sunspot}

The statistically coherent repair is to make selection part of the model by attaching a probability law over the branches of every multiplicity region. For the scalar problem in \eqref{eq:95}, this means assigning a probability $p_t$ to the bound branch. The selection law may be Markov in a latent sunspot state $s_t$. The resulting inferential target is the mixed-support posterior \eqref{eq:73}, with $p(s_t\mid s_{t-1},x_{t-1},\vartheta)>0$ on both branches.

Methods for constructing linear models under indeterminacy are well established. Lubik and Schorfheide (2003) show that, under indeterminacy, fundamentals do not uniquely determine endogenous forecast errors. Sunspot shocks enter as additional structural disturbances with loadings that can be estimated. For the ZLB itself, Aruoba, Cuba-Borda, and Schorfheide (2018) provide a published instance. The study is cited for orientation; the survey's authors were unable to obtain its full text. Their New Keynesian model uses a Markov sunspot to select between targeted-inflation and deflation regimes and is estimated with a particle filter. Holden (2023) explicitly contrasts his within-steady-state multiplicity with this literature on switching between steady states.

The consequence for inference is structural. In the setting considered here, sunspot completion is the \emph{only} interpretation under which regime-path Gibbs, PGAS, or mixed HMC are valid, because only this interpretation gives the regime positive conditional probability on both values. Under deterministic selection, the conditioning variables fix the regime, leaving the alternative branch with zero conditional support. These methods then collapse to the zero-support situation of \eqref{eq:5a}.

\subsection{Estimation baselines}
\label{sec:nk-baselines}

Three filter families define current estimation practice for this lower-bound model class. All are approximations relative to the particle approach treated as the authority in Section~\ref{sec:particle}. The piecewise Kalman filter of Section~\ref{sec:linear} conditions on the verified regime path, taking that path as given when filtering. Boehl and Strobel (2023) estimate medium-scale ELB models with an ensemble Kalman filter (EnKF) on top of the Boehl (2022) solver. Their EnKF updates ensemble members by linear shifting rather than reweighting, and they compare this approach against the inversion filter. The ensemble representation does not remove the limitation of Gaussian closure: like every Gaussian-closure method, the EnKF inherits the branch-mass deficiency at the bound described in Section~\ref{sec:ukf-branchloss}. Sampling noise adds to that deficiency. Holden (2017) proposes an extended skew-t cubature Kalman filter with dynamic state-space reduction. This filter tracks third and fourth moments precisely because Gaussian closure loses the asymmetry created by censoring.

These limitations motivate comparing filters on common data. A serious estimation exercise in this class should report the piecewise Kalman filter, an ensemble Kalman filter arm, an inversion-filter arm (Cuba-Borda, Guerrieri, Iacoviello, and Zhong 2019, cited here for orientation), and the bootstrap particle filter on the same data sets. The particle approach should be the authority in this comparison, mirroring the comparison of estimation routes in Section~\ref{sec:shadow-routes} of the first case study.

\subsection{Designing an estimation exercise}
\label{sec:nk-design}

For this model class, the posterior is not well-defined until the modeling choices that determine it have been made explicit. Software defaults usually leave these choices implicit. The specification below uses the notation of \eqref{eq:1}--\eqref{eq:3} because each choice changes the inferential target, rather than merely the algorithm used to estimate an already defined target.

First, the bound and the truncated policy rule must be specified: which variable is bounded, at what level, and how the notional rule behaves while the bound binds. In particular, the specification must say whether the notional rule responds to lagged states. That response is exactly what generated multiplicity in Section~\ref{sec:nk-mult}.

Second, expectation formation and the terminal condition must be specified. This means stating the perfect-foresight horizon and the reference regime to which the path returns. The parametric-LCP regularity result of Section~\ref{sec:nk-solvers} holds for a fixed horizon and terminal condition. Changing either during estimation therefore changes the continuity properties of the likelihood as well.

Third, the solution operator must be specified together with its verification loop. Whichever of the OccBin, news-shock, or closed-form iterations is used, the estimated object is the \emph{verified} path. An unverified guess accepted by a lenient stopping rule is a different, and discontinuously different, object from that verified path.

Fourth, the uniqueness domain must be established using the LCP matrix $M$ of \eqref{eq:93} for the chosen horizon. This requires a computed P-matrix verdict on the parameter region that the posterior will actually explore, so that the domain on which the solution is unique is known rather than assumed.

Fifth, every multiplicity region needs a selection law: either a stated deterministic rule or a stochastic rule with its sunspot process specified. As \eqref{eq:97} shows, the likelihood jumps wherever an unstated selection switches branch. The choice among equilibrium solutions must therefore be part of the stated model.

Sixth, the specification must state what value the likelihood assigns on nonexistence regions, where no verified path returns to the steady state. The alternatives are zero likelihood, a penalty, or an enlarged model in which such histories have positive probability. Each choice produces a different posterior and a different support.

Seventh, the measurement equation and error structure must be specified. These determine whether the state-level kinks and jumps analyzed above survive into the observed-data likelihood at all.

Finally, the estimation design should include a deliberately tiny version of the model, small enough for exhaustive enumeration of its regime paths. This provides an exact-enumeration reference against which to compare the piecewise Kalman filter, the ensemble arm, and the particle authority. Once the modeling choices are frozen, the geometry classification of Section~\ref{sec:geometries} can be applied to the model actually being estimated rather than to its informal description. The guidance of Section~\ref{sec:guidance} then applies.

\section{Solver-induced and code-induced discontinuities}
\label{sec:solver}

Discontinuity in the \emph{model} is only part of the problem. A second, independent family of discontinuities arises in the \emph{evaluation} of the model. These numerical discontinuities matter for gradient samplers even when the intended target is smooth. The distinction must remain clear: a model-level jump changes the target and belongs to the discussions in Section~\ref{sec:geometries} and Section~\ref{sec:targets}; an evaluation-level branch leaves the intended target smooth but makes the \emph{computed} log density or its gradient discontinuous.

Three mechanisms recur in practice. First, matrix-function kernels can involve integer branch decisions. A matrix exponential evaluated by Pad\'e approximation with scaling-and-squaring uses an integer scaling count that jumps at norm thresholds. Automatic differentiation inherits this branch choice. In one production implementation examined by the authors, the derivative is computed with the scaling integer frozen, precisely because no smooth-derivative claim holds across the branch. At such a branch, the \emph{value} is continuous up to the approximation tolerance, whereas the automatic derivative can jump by a finite amount.

Second, factorization conventions can introduce discontinuities in the numerical representation. QR and SVD sign and ordering choices, and square-root filter downdates near rank changes, are branch-discontinuous functions of their inputs even though the objects they represent vary smoothly.

Third, almost-everywhere gradients require a choice at ties. Reverse-mode differentiation of $\max$, elementwise selection, absolute values, and sorting picks one subgradient at a tie. By Section~\ref{sec:kink-validity}, this choice is harmless for validity \emph{provided the value is continuous}.

For HMC, the consequences depend on that continuity condition. If the computed log density is continuous and only its computed gradient jumps at evaluation branches, the situation is exactly the kink case. The Metropolis correction keeps the intended target invariant; the cost is the $O(\epsilon\|\Delta\nabla U\|)$ acceptance penalty of \eqref{eq:56a} at branch crossings. Thus the gradient discontinuity affects acceptance without changing the invariant target under this condition.

If the computed \emph{value} itself jumps, the sampler instead targets a discontinuous computed density without recognizing the discontinuity. Such jumps can arise when a solver switches branches beyond its tolerance, when an iteration is truncated by a state-dependent stopping rule, or when a fixed-point loop finds a different verified path. Rejections then concentrate on branch boundaries, apparent reversibility holds only up to the jump size, and no diagnostic inside the sampler distinguishes this evaluation-level discontinuity from a genuine model jump.

Serious runs therefore require an evaluation-level test of the differences in computed log density across branch boundaries:
\begin{equation}\label{eq:98}
 \sup_{\text{branch boundaries}}
 \bigl|\log\widehat p_{\,\text{branch}\,1}-\log\widehat p_{\,\text{branch}\,2}\bigr|
 \;\ll\;1
\end{equation}
The bound concerns the largest of these discrepancies, measured in Metropolis energy units, and must be verified by test. The implementation must also use branch-stable derivatives, with the fixed-scaling matrix-exponential treatment above as the template, and deterministic tie-breaking. Value-continuity tests must cover every known integer branch: scaling orders, factorization signs, regime-path verification loops, and truncation horizons. These are testing obligations on the likelihood code, not sampler tuning knobs. The next section's guidance assumes that the code passes them.

\section{Practical guidance}
\label{sec:guidance}

The technical development closes with advice for a reader planning Bayesian inference for a ZLB model. The order of the stages below matters because each makes a later stage interpretable. Each also relies on the ladder of standard test models assembled in Table~\ref{tab:testmodels}.

\subsection*{Define the target and its geometry first}

Before any sampler code is written, the model must be stated in the notation of \eqref{eq:1}--\eqref{eq:3}. Its geometry must be identified as a kink, a density jump, or an explicit discrete path using the taxonomy of Section~\ref{sec:geometries}, and the treatment of multiple equilibrium solutions must be specified. Section~\ref{sec:shadow} carries out this exercise for the shadow-rate term-structure case study. Section~\ref{sec:nk-design} lists what must be specified for the New Keynesian case, including a computed LCP uniqueness verdict and a stated selection law. Nothing downstream can compensate for omitting this specification: a sampler cannot be validated against a target that has not been defined.

\subsection*{A ladder of standard test models}
\label{sec:testmodels}

Every stage below presupposes a reference problem whose answer is known by means independent of the code being tested. The literature has converged on a reasonably stable set of such problems, which are worth assembling in one place. They are organized here by the source of the independent answer: a closed form, exact enumeration, a published reference implementation, or agreement between methods that are exact for different reasons. Table~\ref{tab:testmodels} summarizes the ladder; the paragraphs that follow explain what each rung isolates.

\begin{table}[htbp]\small
\begin{tabularx}{\textwidth}{@{}>{\raggedright\arraybackslash}p{0.16\textwidth}>{\raggedright\arraybackslash}p{0.34\textwidth}X@{}}
\toprule
\textbf{Tier} & \textbf{Model} & \textbf{Independent answer and what it isolates} \\
\midrule
0 & Linear-Gaussian state-space model & Kalman filter likelihood in closed form; every particle, sigma-point, and joint-HMC route must reproduce it before a bound is introduced \\
\addlinespace
0 & Conjugate parameter posterior (known states) & Analytic posterior; calibrates the parameter sampler in isolation \\
\addlinespace
1 & Scalar censored-observation model, \eqref{eq:29}--\eqref{eq:35} & Closed-form mixture likelihood; atom versus smoothed-bound handling, branch masses \\
\addlinespace
1 & Affine-boundary truncated Gaussian & Closed-form moments \eqref{eq:21}--\eqref{eq:28}; unit tests for every truncation step \\
\addlinespace
1 & Truncated multivariate Gaussian sampling & Pakman and Paninski (2014) exact HMC and Botev (2017) tilting as reference draws; cell-restricted proposals \\
\addlinespace
2 & Two-state switching linear-Gaussian model, horizon $T\le 12$ & Exact likelihood and smoothing by enumerating all $2^{T}$ regime paths; PF unbiasedness, PGAS invariance, mixed-HMC path moves \\
\addlinespace
2 & Univariate nonlinear growth model (Gordon, Salmond, and Smith 1993) & Dense-grid numerical filter; the canonical particle-filter stress case: severe nonlinearity and multimodality \\
\addlinespace
3 & One-factor shadow-rate model (Gorovoi and Linetsky 2004) & Closed-form bond prices under a hard bound; analytic anchor for the case study of Section~\ref{sec:shadow} \\
\addlinespace
3 & One-node quadrature reduction of $\mathcal C_1$ & Collapses to the Tier-1 censored-scalar likelihood; exact posterior by low-dimensional grid integration \\
\addlinespace
3 & Lagged-rule economy, \eqref{eq:94}--\eqref{eq:96} & Analytic multiplicity boundary $-r/A^{2}$; detection of nonexistence and selection-rule discontinuities \\
\addlinespace
3 & Published OccBin examples (Guerrieri and Iacoviello 2015) & Dynare reference output; regime-path and piecewise-Kalman bookkeeping \\
\addlinespace
3 & Small-scale NK ZLB model (Aruoba et al.\ 2021) & Published COPF likelihood comparisons; conditionally optimal proposal correctness \\
\addlinespace
4 & Joint-distribution tests (Geweke 2004) & Compares two simulators of the same joint law; needs no analytic posterior \\
\addlinespace
4 & Simulation-based calibration (Talts et al.\ 2018) & Rank-uniformity of the true parameter under the fitted posterior across simulated data sets \\
\bottomrule
\end{tabularx}
\caption{A ladder of standard test models for verifying ZLB inference
algorithms, ordered by the source of the independent answer.}
\label{tab:testmodels}
\end{table}

The tiers are ordered so that each new difficulty can be tested on its own. Tier~0 contains no bound at all. For a linear-Gaussian model, the Kalman filter gives the exact marginal likelihood. A bootstrap filter must match that likelihood in expectation at every particle count, where the expectation averages over the filter's randomness. A sigma-point filter must match it to rounding, and joint HMC on the non-centered chart of \eqref{eq:56} must reproduce the analytic posterior. The conjugate benchmark with known states checks the parameter sampler in isolation against an analytic parameter posterior. A method that fails these checks has a bookkeeping error that no ZLB machinery will fix. This is why the rung comes first, even though it tests nothing specific to the bound.

Tier~1 introduces the bound in the smallest settings that still have closed-form answers. The scalar censored-observation model derived in \eqref{eq:29}--\eqref{eq:35} has a closed-form, two-term mixture likelihood. Its separate mixture contributions allow atom handling and branch-mass computation to be checked individually. The same likelihood provides an exact baseline for assessing the bias introduced by smoothing the bound: evaluating the softplus variant on the same data measures the likelihood gap of \eqref{eq:90}--\eqref{eq:90a} directly against the true likelihood. The affine-boundary truncated-Gaussian moments in \eqref{eq:21}--\eqref{eq:28} provide unit tests for every truncation step inside a filter. Exact truncated-Gaussian samplers also provide reference draws against which any cell-restricted proposal can be tested. These comparisons assess whether the proposal produces the intended distribution within the cell; checking only that draws remain in the cell establishes support, not distributional correctness.

Tier~2 replaces closed-form reference answers with exact enumeration. In a two-state switching linear-Gaussian model with a short horizon, conditioning on each of the $2^{T}$ regime paths reduces the model to \eqref{eq:6}--\eqref{eq:11}. The exact likelihood is then a finite sum of conditional likelihoods weighted by regime-path probabilities. The exact smoothing distribution is a finite mixture of the conditional smoothing distributions, weighted by the posterior probabilities of the paths. These independent answers make the model a natural benchmark for checking that the particle estimator \eqref{eq:13} is unbiased in repeated sampling. They also allow a check of PGAS invariance: a PGAS update started from the path posterior must preserve that distribution. Mixed-HMC path moves must accept at the rates predicted by enumeration. In a deterministic-regime variant of the same model, regime-only proposals must be rejected exactly as the support calculation in Section~\ref{sec:targets} requires.

The univariate nonlinear growth model of Gordon, Salmond, and Smith (1993) then adds severe nonlinearity and multimodality, with a dense-grid numerical filter as the reference. Although it contains no bound, it is the standard stress case for particle-filter variance and resampling behavior. A filter that degrades unexpectedly on this model will also degrade near a binding spell.

Tier~3 tests features specific to the lower bound. The one-factor shadow-rate model of Gorovoi and Linetsky (2004) gives closed-form bond prices under a hard bound, making it the natural analytic anchor for the term-structure case study. The quadrature, censoring, and filtering layers of Section~\ref{sec:shadow-model} can each be checked against these prices, whose calculation involves none of those layers.

A complementary check simplifies the $\mathcal C_1$ working example itself. Reducing it to a single quadrature node makes each yield a censored affine function of the factors, so the Tier-1 machinery applies without alteration. The exact posterior is then available by grid integration over the low-dimensional factor space. The bootstrap-filter and joint-HMC methods of Section~\ref{sec:shadow-routes} must agree with that posterior grid before they are trusted at $K=40$.

On the New Keynesian side, the lagged-rule economy of \eqref{eq:94}--\eqref{eq:96} has an analytic multiplicity boundary. This known boundary allows tests of whether the estimation code detects nonexistence, reports multiplicity, and detects selection-rule discontinuities. The code must implement the declared selection rule rather than an accidental rule determined by the order in which the solver finds solutions. No smooth benchmark can provide these tests of nonexistence, multiplicity, and selection. The published OccBin examples and the small-scale New Keynesian model of Aruoba et al.\ (2021) complete the tier with reference implementations for piecewise-linear path iteration and the conditionally optimal proposal, respectively.

Tier~4 tests the sampler as a whole and requires no analytic posterior. Geweke's (2004) joint-distribution test compares two simulators of the same joint law of parameters and data. One generates that law by construction; the other uses the posterior sampler. Their comparison detects errors anywhere in the chain. Simulation-based calibration (Talts et al.\ 2018) checks the rank of the true, data-generating parameter among draws from the fitted posterior. These ranks should be uniform across data sets simulated from the prior. Both tests apply unchanged to ZLB models. They are particularly useful here because they exercise binding and nonbinding regions in prior proportion, including the mixed and boundary cases that a hand-picked test set tends to miss.

A practical validation suite runs the ladder in order and advances a method only when every applicable rung passes. It then retains the rungs as regression tests, because the solver-branch failures of Section~\ref{sec:solver} tend to reappear after unrelated changes.

\subsection*{Validate a linear benchmark}

The implementation should reproduce a published piecewise-linear regime path and its corresponding piecewise-Kalman likelihood, using Dynare/OccBin output as the natural reference. On synthetic data, the comparison should cover every innovation, covariance, regime inequality, and log-likelihood contribution. Agreement on the same calculation rules out an entire class of bookkeeping errors before nonlinear machinery obscures them. The linear method itself cannot be the final answer, for the reasons explained in Section~\ref{sec:linear-limits}. If the verification fixed point is not unique on the test data, that nonuniqueness should be recorded. It is direct evidence about the model's solution correspondence, which Section~\ref{sec:nk-mult} shows drives the estimation difficulties.

\subsection*{Build a particle likelihood authority}

The next step is to implement the bootstrap particle filter of Section~\ref{sec:pf-generic}, retaining particle ancestry and likelihood replicates. The implementation should include PGAS and, where it applies, the conditional-optimal special case of Section~\ref{sec:copf}. The likelihood estimator should first be checked against exact enumeration in a deliberately tiny model. Its variance should then be measured across particle counts.

For a hard-bound observation-map model such as the $\mathcal{C}_1$ target of Section~\ref{sec:shadow}, the bootstrap hard-bound filter is the natural authority for the likelihood calculation. Adapted-proposal and moment-closure constructions have a different role. The cell decomposition and branch-split mixture UKF of Section~\ref{sec:ukf} are proposal and diagnostic layers; they do not replace that authority. These constructions should be unit tested against the affine closed forms of Section~\ref{sec:linear-kalman}.

Posterior sensitivity to nearby model variants also belongs at this stage. The continuous-maturity and discrete hard-max variants of Section~\ref{sec:shadow-ladder} provide an example. This sensitivity must be assessed before a kink-aware sampler design is selected, because it determines whether the design effort is aimed at a posterior that is stable across nearby specifications. A differentiable particle-filter method should remain diagnostic until its extended target and Hamiltonian correction have been derived, as specified in Section~\ref{sec:dpf}.

\subsection*{Only then, target-specific exact kernels}

The support structure established in the first stage determines which exact
kernel is appropriate. Validation evidence does not transfer between the
stochastic regime, deterministic threshold, and continuous kink cases.

For a stochastic regime model, fixed-path HMC should update the continuous
variables conditional on the regime path, with Metropolis or PGAS moves
updating that path (Section~\ref{sec:pgas}, Section~\ref{sec:joint-hmc}). The invariant distribution of this combined sampler must be
validated against exact enumeration.

For a deterministic threshold model, the baseline should use non-gradient
joint-state sampling or particle MCMC. Validation of this baseline must
verify that regime-only moves are rejected exactly as the support calculation
in Section~\ref{sec:targets} predicts. The threshold determines the regime for a given
state, so a sampler that accepts regime-only moves is sampling a different
target.

For a continuous kink target, both the pseudo-marginal and joint-HMC methods
in Section~\ref{sec:pmcmc} are exact. Their agreement on matched synthetic data sets
is the natural acceptance criterion for the implementations. The
solver-branch value-continuity tests in Section~\ref{sec:solver} are a precondition for
trusting any gradient method.

\subsection*{Event-aware HMC only for a genuine switching surface}

Reflection/refraction and discontinuous HMC (Section~\ref{sec:event}, Section~\ref{sec:dhmc})
justify their added complexity only when the first stage has identified an
explicit continuous switching surface across which the density genuinely
jumps. In that case, validation must check whether the trajectory's first
crossing is located accurately and whether the full posterior is evaluated
correctly from each side of the surface. It must also test energy conservation
at reflections and refractions, reversibility, and behavior when a trajectory
grazes the boundary.

If the boundary locator fails on the actual model, the event mechanism has no
located boundary at which to act. Smoothing the surface to make the mechanism
work substitutes a different target, as Section~\ref{sec:relaxations} makes precise.
The event-based approach must then be abandoned or the boundary representation
repaired.

\subsection*{Test value continuity before trusting gradients}

Finally, at every scale, interpreting a gradient-based sampler's diagnostics
requires the evaluation-level tests of Section~\ref{sec:solver}, the branch-boundary
bound \eqref{eq:98}, branch-stable derivatives, and deterministic tie-breaking.
Without these prerequisites, a rejection cluster has no clear interpretation:
it may reflect the model's geometry, or it may arise from a solver branch that
the sampler does not account for.

\section{Conclusions and limits}
\label{sec:conclusion}

The literature provides useful exact building blocks for nonlinear-ZLB HMC,
but it does not provide a finished methodology. These contributions address
different parts of estimation. OccBin and the piecewise Kalman filter solve
the conditional linear problem. Aruoba et al. show how a piecewise-linear
regime mixture can support an efficient particle likelihood. Afshar--Domke
and Nishimura--Dunson--Lu show how Hamiltonian motion can cross true energy
discontinuities. Zhou shows how to preserve a mixed discrete-continuous target
without an artificial ordinal embedding. Childers et al. show the efficiency
of joint latent-state HMC for smooth differentiable state-space models, while
Alenl\"ov et al. state the additional conditions needed for pseudo-marginal
HMC.

The defensible direction therefore depends on the model's support. A
deterministic ZLB calls for a linear piecewise-Kalman benchmark, a nonlinear
particle method as the reference for validation, and eventually an event-aware
continuous kernel. A stochastic regime calls for a mixed discrete-continuous
target, fixed-path HMC, and PGAS or mixed HMC. Softplus and Concrete relaxations
are useful comparison arms, but both change the target.

The case studies turn this dependence on model support into concrete
estimation programs. For the shadow-rate term-structure model of
Section~\ref{sec:shadow}, the hard-bound observation-map construction defines a ladder
of targets. Its hard-bound variants are continuous kink targets: the density
remains continuous where its slope changes. The exact methods for these
variants are the bootstrap hard-bound particle filter with PMMH and joint
kinked-target HMC. They must agree on matched synthetic data.

The fully adapted cell decomposition provides proposals and diagnostics,
conditional on controlled numerical evaluation of polytope probabilities and
simulation of truncated draws. The softplus and UKF methods remain declared
approximations, with stated bias bounds and a sensitivity protocol. The domain
benchmark family is the closure-filter shadow-rate line from Black through
Krippner, Wu--Xia, Kim--Priebsch, Christensen--Rudebusch, and Opschoor--van der
Wel. Any claim of likelihood fidelity must be measured against Pericoli and
Taboga's closure-free data-augmentation MCMC.

For the New Keynesian model of Section~\ref{sec:nk}, the worked linear-complementarity
analysis shows that the genuine discontinuity lies in the solution
correspondence. It finds both nonexistence of verified paths and pervasive
multiplicity. The first modeling step must therefore be to declare a selection
law, either deterministic or sunspot-stochastic. The mixed-support kernels
apply only under the stochastic completion.

A final limitation is that no literature survey establishes the posterior
correctness of a particular implementation. That requires the validation
program sketched in Section~\ref{sec:guidance}: a defined target, a tied-out linear
benchmark, and a particle reference checked against enumeration. The exact
kernels must be validated against that particle reference, and the evaluation
code must undergo value-continuity tests. Every part of this program must be
carried out on the model actually being estimated.


\section*{References}
\addcontentsline{toc}{section}{References}
\begin{list}{}{\setlength{\leftmargin}{2em}\setlength{\itemindent}{-2em}\setlength{\itemsep}{2pt}\setlength{\parsep}{0pt}}

\item Afshar, Hadi Mohasel, and Justin Domke. 2015. ``Reflection, Refraction, and Hamiltonian Monte Carlo.'' \emph{Advances in Neural Information Processing Systems} 28: 3007--3015. \url{http://hdl.handle.net/1885/103828}.

\item Alenl\"ov, Johan, Arnaud Doucet, and Fredrik Lindsten. 2021. ``Pseudo-Marginal Hamiltonian Monte Carlo.'' \emph{Journal of Machine Learning Research} 22(141): 1--45. \url{https://jmlr.org/papers/v22/19-486.html}.

\item Allik, Bethany, Cory Miller, Michael J. Piovoso, and Ryan Zurakowski. 2016. ``The Tobit Kalman Filter: An Estimator for Censored Measurements.'' \emph{IEEE Transactions on Control Systems Technology} 24(1): 365--371. \url{https://doi.org/10.1109/TCST.2015.2432155}.

\item Alspach, Daniel L., and Harold W. Sorenson. 1972. ``Nonlinear Bayesian Estimation Using Gaussian Sum Approximations.'' \emph{IEEE Transactions on Automatic Control} 17(4): 439--448. \url{https://doi.org/10.1109/TAC.1972.1100034}.

\item Andrieu, Christophe, Arnaud Doucet, and Roman Holenstein. 2010. ``Particle Markov Chain Monte Carlo Methods.'' \emph{Journal of the Royal Statistical Society: Series B} 72(3): 269--342. \url{https://doi.org/10.1111/j.1467-9868.2009.00736.x}.

\item Aruoba, S. Bora\u{g}an, Pablo Cuba-Borda, Kenji Higa-Flores, Frank Schorfheide, and Sergio Villalvazo. 2021. ``Piecewise-Linear Approximations and Filtering for DSGE Models with Occasionally Binding Constraints.'' Federal Reserve Bank of Philadelphia Working Paper 20-13. \url{https://doi.org/10.21799/frbp.wp.2020.13}.

\item Aruoba, S. Bora\u{g}an, Pablo Cuba-Borda, and Frank Schorfheide. 2018. ``Macroeconomic Dynamics Near the ZLB: A Tale of Two Countries.'' \emph{Review of Economic Studies} 85(1): 87--118. \url{https://doi.org/10.1093/restud/rdx027}.

\item Bauer, Michael D., and Glenn D. Rudebusch. 2016. ``Monetary Policy Expectations at the Zero Lower Bound.'' \emph{Journal of Money, Credit and Banking} 48(7): 1439--1465. \url{https://doi.org/10.1111/jmcb.12338}.

\item Black, Fischer. 1995. ``Interest Rates as Options.'' \emph{Journal of Finance} 50(5): 1371--1376. \url{https://doi.org/10.1111/j.1540-6261.1995.tb05182.x}.

\item Blom, Henk A. P., and Yaakov Bar-Shalom. 1988. ``The Interacting Multiple Model Algorithm for Systems with Markovian Switching Coefficients.'' \emph{IEEE Transactions on Automatic Control} 33(8): 780--783. \url{https://doi.org/10.1109/9.1299}.

\item Boehl, Gregor. 2022. ``Efficient Solution and Computation of Models with Occasionally Binding Constraints.'' \emph{Journal of Economic Dynamics and Control} 143: 104523. \url{https://doi.org/10.1016/j.jedc.2022.104523}.

\item Boehl, Gregor, and Felix Strobel. 2023. ``Estimation of DSGE Models with the Effective Lower Bound.'' \emph{Journal of Economic Dynamics and Control} 158: 104784. \url{https://doi.org/10.1016/j.jedc.2023.104784}.

\item Botev, Zdravko I. 2017. ``The Normal Law Under Linear Restrictions: Simulation and Estimation via Minimax Tilting.'' \emph{Journal of the Royal Statistical Society: Series B} 79(1): 125--148. \url{https://doi.org/10.1111/rssb.12162}. [Cited from publication metadata.]

\item Cha\^ari, Lotfi, Jean-Yves Tourneret, Caroline Chaux, and Hadj Batatia. 2016. ``A Hamiltonian Monte Carlo Method for Non-Smooth Energy Sampling.'' \emph{IEEE Transactions on Signal Processing} 64(21): 5585--5594. \url{https://doi.org/10.1109/TSP.2016.2585120}.

\item Childers, David, Jes\'us Fern\'andez-Villaverde, Jesse Perla, Christopher Rackauckas, and Peifan Wu. 2022. ``Differentiable State-Space Models and Hamiltonian Monte Carlo Estimation.'' NBER Working Paper 30573. \url{https://doi.org/10.3386/w30573}.

\item Christensen, Jens H. E., and Glenn D. Rudebusch. 2015. ``Estimating Shadow-Rate Term Structure Models with Near-Zero Yields.'' \emph{Journal of Financial Econometrics} 13(2): 226--259. \url{https://doi.org/10.1093/jjfinec/nbu010}.

\item Cuba-Borda, Pablo, Luca Guerrieri, Matteo Iacoviello, and Molin Zhong. 2019. ``Likelihood Evaluation of Models with Occasionally Binding Constraints.'' \emph{Journal of Applied Econometrics} 34(7): 1073--1085. \url{https://doi.org/10.1002/jae.2733}. [Cited from publication metadata.]

\item Corenflos, Adrien, James Thornton, George Deligiannidis, and Arnaud Doucet. 2021. ``Differentiable Particle Filtering via Entropy-Regularized Optimal Transport.'' \emph{Proceedings of the 38th International Conference on Machine Learning}, PMLR 139: 2100--2111. \url{https://proceedings.mlr.press/v139/corenflos21a.html}.

\item Deligiannidis, George, Arnaud Doucet, and Michael K. Pitt. 2018. ``The Correlated Pseudomarginal Method.'' \emph{Journal of the Royal Statistical Society: Series B} 80(5): 839--870. \url{https://doi.org/10.1111/rssb.12280}. [Cited from publication metadata.]

\item Garc\'ia-Fern\'andez, \'Angel F., Mark R. Morelande, and Jes\'us Grajal. 2012a. ``Truncated Unscented Kalman Filtering.'' \emph{IEEE Transactions on Signal Processing} 60(7): 3372--3386. \url{https://doi.org/10.1109/TSP.2012.2193393}.

\item Garc\'ia-Fern\'andez, \'Angel F., Mark R. Morelande, and Jes\'us Grajal. 2012b. ``Mixture Truncated Unscented Kalman Filtering.'' In \emph{Proceedings of the 15th International Conference on Information Fusion (FUSION)}, 479--486. Singapore.

\item Geng, Hang, Hongjian Liu, Lifeng Ma, and Xiaojian Yi. 2021. ``Multi-Sensor Filtering Fusion Meets Censored Measurements Under a Constrained Network Environment: Advances, Challenges and Prospects.'' \emph{International Journal of Systems Science} 52(16): 3410--3436. \url{https://doi.org/10.1080/00207721.2021.2005178}.

\item Genz, Alan, and Frank Bretz. 2009. \emph{Computation of Multivariate Normal and t Probabilities}. Lecture Notes in Statistics 195. Berlin: Springer. \url{https://doi.org/10.1007/978-3-642-01689-9}. [Cited from publication metadata.]

\item Geweke, John. 2004. ``Getting It Right: Joint Distribution Tests of Posterior Simulators.'' \emph{Journal of the American Statistical Association} 99(467): 799--804. \url{https://doi.org/10.1198/016214504000001132}. [Cited from publication metadata.]

\item Giovannini, Massimo, Philipp Pfeiffer, and Marco Ratto. 2021. ``Efficient and Robust Inference of Models with Occasionally Binding Constraints.'' JRC Working Papers in Economics and Finance 2021/3. \url{https://hdl.handle.net/10419/249365}.

\item Gokce, Murat, and Mustafa Kuzuoglu. 2015. ``Unscented Kalman Filter-Aided Gaussian Sum Filter.'' \emph{IET Radar, Sonar \& Navigation} 9(5): 589--599. \url{https://doi.org/10.1049/iet-rsn.2014.0088}.

\item Gordon, Neil J., David J. Salmond, and Adrian F. M. Smith. 1993. ``Novel Approach to Nonlinear/Non-Gaussian Bayesian State Estimation.'' \emph{IEE Proceedings F (Radar and Signal Processing)} 140(2): 107--113. \url{https://doi.org/10.1049/ip-f-2.1993.0015}. [This reference is cited from publication metadata.]

\item Gorovoi, Viatcheslav, and Vadim Linetsky. 2004. ``Black's Model of Interest Rates as Options, Eigenfunction Expansions and Japanese Interest Rates.'' \emph{Mathematical Finance} 14(1): 49--78. \url{https://doi.org/10.1111/j.0960-1627.2004.00181.x}. [This reference is cited from publication metadata.]

\item Guerrieri, Luca, and Matteo Iacoviello. 2015. ``OccBin: A Toolkit for Solving Dynamic Models with Occasionally Binding Constraints Easily.'' \emph{Journal of Monetary Economics} 70: 22--38. \url{https://doi.org/10.1016/j.jmoneco.2014.08.005}.

\item Holden, Tom D. 2016. ``Computation of Solutions to Dynamic Models with Occasionally Binding Constraints.'' Unpublished working paper, University of Surrey. \url{https://github.com/tholden/dynareOBC}.

\item Holden, Tom D., and Michael Paetz. 2012. ``Efficient Simulation of DSGE Models with Inequality Constraints.'' Quantitative Macroeconomics Working Papers 21207b, Hamburg University. [This reference is cited from publication metadata. The work presents the original anticipated-shock approach.]

\item Holden, Tom D. 2017. ``Tractable Estimation and Smoothing of Highly Nonlinear Dynamic State-Space Models.'' Unpublished working paper, University of Surrey.

\item Holden, Tom D. 2023. ``Existence and Uniqueness of Solutions to Dynamic Models with Occasionally Binding Constraints.'' \emph{Review of Economics and Statistics} 105(6): 1481--1499. \url{https://doi.org/10.1162/rest_a_01122}.

\item Kandepu, Rambabu, Lars Imsland, and Bjarne A. Foss. 2008. ``Constrained State Estimation Using the Unscented Kalman Filter.'' In \emph{Proceedings of the 16th Mediterranean Conference on Control and Automation}, 1453--1458. Ajaccio: IEEE. \url{https://doi.org/10.1109/MED.2008.4602001}.

\item Kim, Don H., and Marcel Priebsch. 2013. ``Estimation of Multi-Factor Shadow-Rate Term Structure Models.'' Preliminary draft, October 9. Washington: Board of Governors of the Federal Reserve System.

\item Krippner, Leo. 2012. ``Modifying Gaussian Term Structure Models When Interest Rates Are Near the Zero Lower Bound.'' Reserve Bank of New Zealand Discussion Paper DP2012/02.

\item Lemke, Wolfgang, and Andreea Liliana Vladu. 2017. ``Below the Zero Lower Bound: A Shadow-Rate Term Structure Model for the Euro Area.'' ECB Working Paper No. 1991. Frankfurt: European Central Bank.

\item Lindsten, Fredrik, Michael I. Jordan, and Thomas B. Sch\"on. 2014. ``Particle Gibbs with Ancestor Sampling.'' \emph{Journal of Machine Learning Research} 15: 2145--2184. \url{https://jmlr.org/papers/v15/lindsten14a.html}.

\item Lubik, Thomas A., and Frank Schorfheide. 2003. ``Computing Sunspot Equilibria in Linear Rational Expectations Models.'' \emph{Journal of Economic Dynamics and Control} 28(2): 273--285. \url{https://doi.org/10.1016/S0165-1889(02)00153-7}.

\item Neal, Radford M. 2011. ``MCMC Using Hamiltonian Dynamics.'' In \emph{Handbook of Markov Chain Monte Carlo}, edited by Steve Brooks, Andrew Gelman, Galin L. Jones, and Xiao-Li Meng, 113--162. Chapman and Hall/CRC. \url{https://doi.org/10.1201/b10905-6}.

\item Nishimura, Akihiko, David B. Dunson, and Jianfeng Lu. 2020. ``Discontinuous Hamiltonian Monte Carlo for Discrete Parameters and Discontinuous Likelihoods.'' \emph{Biometrika} 107(2): 365--380. \url{https://doi.org/10.1093/biomet/asz083}.

\item Opschoor, Daan, and Michel van der Wel. 2024. ``A Smooth Shadow-Rate Dynamic Nelson-Siegel Model for Yields at the Zero Lower Bound.'' \emph{Journal of Business \& Economic Statistics} 43(2): 298--311. \url{https://doi.org/10.1080/07350015.2024.2365779}.

\item Pakman, Ari, and Liam Paninski. 2014. ``Exact Hamiltonian Monte Carlo for Truncated Multivariate Gaussians.'' \emph{Journal of Computational and Graphical Statistics} 23(2): 518--542. \url{https://doi.org/10.1080/10618600.2013.788448}.

\item Pericoli, Marcello, and Marco Taboga. 2018. ``Nearly Exact Bayesian Estimation of Non-Linear No-Arbitrage Term-Structure Models.'' Banca d'Italia Temi di discussione (Working Papers) 1189, September 2018. [The citation refers to the 2018 working-paper version.]

\item Poyiadjis, George, Arnaud Doucet, and Sumeetpal S. Singh. 2011. ``Particle Approximations of the Score and Observed Information Matrix in State Space Models with Application to Parameter Estimation.'' \emph{Biometrika} 98(1): 65--80. \url{https://doi.org/10.1093/biomet/asq062}.

\item Priebsch, Marcel A. 2013. ``Computing Arbitrage-Free Yields in Multi-Factor Gaussian Shadow-Rate Term Structure Models.'' Finance and Economics Discussion Series 2013-63. Washington: Board of Governors of the Federal Reserve System. \url{https://doi.org/10.17016/feds.2013.63}.

\item \'Scibior, Adam, and Frank Wood. 2021. ``Differentiable Particle Filtering without Modifying the Forward Pass.'' arXiv:2106.10314v2. \url{https://arxiv.org/abs/2106.10314}.

\item Talts, Sean, Michael Betancourt, Daniel Simpson, Aki Vehtari, and Andrew Gelman. 2018. ``Validating Bayesian Inference Algorithms with Simulation-Based Calibration.'' arXiv:1804.06788. \url{https://arxiv.org/abs/1804.06788}. [This reference is cited from publication metadata.]

\item Teixeira, Bruno O. S., Leonardo A. B. T\^orres, Luis A. Aguirre, and Dennis S. Bernstein. 2010. ``On Unscented Kalman Filtering with State Interval Constraints.'' \emph{Journal of Process Control} 20(1): 45--57. \url{https://doi.org/10.1016/j.jprocont.2009.10.007}.

\item Torgander, Jakob, M\aa{}ns Magnusson, and Jonas Wallin. 2024. ``Hamiltonian Monte Carlo with Categorical Parameters Using the Concrete Distribution.'' Workshop at the 6th Symposium on Advances in Approximate Bayesian Inference, non-archival, 1--12.

\item van der Merwe, Rudolph, Arnaud Doucet, Nando de Freitas, and Eric A. Wan. 2000. ``The Unscented Particle Filter.'' In \emph{Advances in Neural Information Processing Systems} 13: 584--590.

\item Wang, Yanbo, Fasheng Wang, Jianjun He, and Fuming Sun. 2020. ``Iterative Truncated Unscented Particle Filter.'' \emph{Information} 11(4): 214. \url{https://doi.org/10.3390/info11040214}.

\item Wu, Jing Cynthia, and Fan Dora Xia. 2016. ``Measuring the Macroeconomic Impact of Monetary Policy at the Zero Lower Bound.'' \emph{Journal of Money, Credit and Banking} 48(2--3): 253--291. \url{https://doi.org/10.1111/jmcb.12300}.

\item Zhang, Hongwei, Xiaohu Zhang, Weixin Xie, and Xia Yang. 2020. ``Constrained Multiple Model UK Filter.'' In \emph{Proceedings of the 15th IEEE International Conference on Signal Processing (ICSP)}, 48--51. Beijing: IEEE. \url{https://doi.org/10.1109/ICSP48669.2020.9320991}.

\item Zhou, Guangyao. 2020. ``Mixed Hamiltonian Monte Carlo for Mixed Discrete and Continuous Variables.'' \emph{Advances in Neural Information Processing Systems} 33. \url{https://arxiv.org/abs/1909.04852}.

\end{list}


\end{document}
